\documentclass[a4paper,11pt]{amsart}

% =========================================================================
%  Electromagnetismo I — Tema 0: Herramientas matemáticas
%  Apuntes formalizados. Compilar con: pdflatex + biber + pdflatex ×2
%  (o latexmk -pdf main.tex).  Todo el preámbulo está en preambulo.tex.
% =========================================================================
% =========================================================================
%  Electromagnetismo I — Tema 0: Herramientas matemáticas
%  Estilo: artículo científico minimalista, heredado del informe «Efecto
%  Doppler con ultrasonidos» (amsart, Palatino, versalitas, enunciados
%  numerados por sección, paleta de azules y grises sobrios), con la
%  notación vectorial de «Cálculo integral multivariable» (L. Navas) y las
%  expresiones del formulario del curso.
% =========================================================================

% --- IDIOMA Y CODIFICACIÓN ------------------------------------------------
\usepackage[utf8]{inputenc}
\usepackage[T1]{fontenc}
\usepackage[spanish,es-tabla,es-nodecimaldot]{babel}

% --- TIPOGRAFÍA (Palatino, como el documento de referencia) ---------------
\usepackage{mathpazo}
\linespread{1.05}
\usepackage[final]{microtype}

% --- MATEMÁTICAS -------------------------------------------------------------
\usepackage{amsmath, amssymb, mathtools}
\usepackage{bm}
\usepackage{mathrsfs}

% --- PÁGINA, COLORES Y GRÁFICOS --------------------------------------------
\usepackage[a4paper,margin=2.6cm]{geometry}
\usepackage{graphicx}
\usepackage{xcolor}
% Paleta base (idéntica a la del informe de ecuaciones diferenciales)
\definecolor{azulprof}{HTML}{1B365D}    % azul marino
\definecolor{azulacero}{HTML}{3B6EA8}   % azul acero
\definecolor{azulclaro}{HTML}{B7CBE3}   % azul hielo
\definecolor{grispizarra}{HTML}{5B6770} % gris pizarra
\definecolor{grisclaro}{HTML}{E4E8EC}   % gris perla
% Paleta ampliada para campos escalares y vectoriales (magnitud creciente:
% índigo → púrpura → carmín → naranja → ámbar), sobria y legible en impresión
\definecolor{cm0}{RGB}{38,30,112}
\definecolor{cm1}{RGB}{112,36,140}
\definecolor{cm2}{RGB}{190,52,118}
\definecolor{cm3}{RGB}{236,104,66}
\definecolor{cm4}{RGB}{252,184,42}
% Colores de apoyo para diagramas geométricos
\definecolor{rojoteo}{HTML}{B3262E}     % trayectorias y circuitos
\definecolor{verdeteo}{HTML}{2E8B57}    % vector unitario / normal
\definecolor{ocreteo}{HTML}{C98A12}     % superficies y volúmenes

\usepackage{tikz}
\usepackage{tikz-3dplot}
\usetikzlibrary{babel, arrows.meta, calc, positioning, decorations.pathmorphing,
                decorations.markings, patterns, 3d, intersections, angles, quotes,
                fadings, shapes.geometric}
\usepackage{pgfplots}
\pgfplotsset{compat=1.17}
\usepgfplotslibrary{colormaps, patchplots, fillbetween}
\pgfplotsset{colormap={campo}{
  rgb255=(38,30,112) rgb255=(112,36,140) rgb255=(190,52,118)
  rgb255=(236,104,66) rgb255=(252,184,42) rgb255=(255,226,110)}}
\pgfplotsset{colormap={campoclaro}{
  rgb255=(232,236,240) rgb255=(183,203,227) rgb255=(59,110,168) rgb255=(27,54,93)}}
\pgfplotsset{colormap={verdeamarillo}{
  rgb255=(120,200,70) rgb255=(190,220,80) rgb255=(250,225,100) rgb255=(250,190,90)}}
\usepackage{booktabs}
\usepackage{array}
\usepackage{enumitem}
\setlist[itemize]{leftmargin=1.6em, itemsep=2pt, topsep=4pt}
\setlist[enumerate]{leftmargin=1.8em, itemsep=2pt, topsep=4pt}

% --- ENLACES Y REFERENCIAS CRUZADAS ----------------------------------------
\usepackage[colorlinks=true, linkcolor=azulacero, citecolor=azulacero,
            urlcolor=azulacero, plainpages=false, pdfusetitle=false,
            pdftitle={Tema 0. Herramientas matemáticas},
            pdfauthor={Daniel Martín Guerrero}]{hyperref}
\usepackage[spanish,nameinlink]{cleveref}

% --- BIBLIOGRAFÍA -----------------------------------------------------------
\usepackage{csquotes}
\usepackage[backend=biber, style=numeric, sorting=none, giveninits=true,
            maxbibnames=6, minbibnames=3, doi=true, url=true,
            isbn=false]{biblatex}
\addbibresource{referencias.bib}
\DeclareFieldFormat[article,techreport,unpublished,book]{title}{#1\isdot}
\DeclareFieldFormat{url}{\url{#1}}
\renewcommand*{\bibfont}{\small}
\setlength{\bibitemsep}{3pt}

% --- CABECERAS ---------------------------------------------------------------
\pagestyle{headings}
\makeatletter
\def\ps@headings{%
  \let\@mkboth\markboth
  \def\@evenhead{\small\scshape\normalfont\thepage\hfil Electromagnetismo I}%
  \def\@oddhead{\small\scshape\normalfont Tema 0. Herramientas matemáticas\hfil\thepage}%
  \let\@evenfoot\@empty \let\@oddfoot\@empty}
\makeatother
\numberwithin{equation}{section}
\setcounter{tocdepth}{2}
\makeatletter
\def\@@and{y}
\renewcommand{\keywordsname}{Palabras clave}
\makeatother

% =========================================================================
%               NOTACIÓN Y MACROS DE USO FRECUENTE
% =========================================================================
% --- Diferencial, derivadas parciales y números -----------------------------
\newcommand{\dif}{\mathrm{d}}
\newcommand{\pd}[2]{\frac{\partial #1}{\partial #2}}
\newcommand{\pdd}[2]{\frac{\partial^{2} #1}{\partial #2^{2}}}
\newcommand{\pdm}[3]{\frac{\partial^{2} #1}{\partial #2\,\partial #3}}
\newcommand{\R}{\mathbb{R}}
\newcommand{\N}{\mathbb{N}}
\newcommand{\Z}{\mathbb{Z}}
\newcommand{\eul}{\mathrm{e}}

% --- Vectores, versores y operador nabla ---------------------------------
\newcommand{\vc}[1]{\vec{#1}}                       % vector con flecha
\newcommand{\uh}[1]{\hat{u}_{#1}}                   % versor û_#1
\newcommand{\ux}{\uh{x}}\newcommand{\uy}{\uh{y}}\newcommand{\uz}{\uh{z}}
\newcommand{\ur}{\uh{r}}\newcommand{\uth}{\uh{\theta}}\newcommand{\uph}{\uh{\varphi}}
\newcommand{\urho}{\uh{\rho}}\newcommand{\uR}{\uh{R}}\newcommand{\un}{\hat{n}}
\newcommand{\nab}{\vec{\nabla}}                     % ∇⃗
\newcommand{\nabp}{\vec{\nabla}'}                   % ∇⃗′ (coordenadas fuente)
\newcommand{\Lap}{\nabla^{2}}                       % laplaciano
\newcommand{\Lapp}{\nabla'^{2}}                      % laplaciano en coordenadas fuente
\newcommand{\norm}[1]{\lVert #1\rVert}
\newcommand{\abs}[1]{\lvert #1\rvert}
\newcommand{\prodesc}[2]{\langle #1,#2\rangle}
\newcommand{\rp}{\vc{r}\,'}                         % r⃗′ punto fuente
\newcommand{\dl}{\dif\vc{\ell}}                     % elemento de línea dirigido
\newcommand{\dS}{\dif\vc{S}}                        % elemento de área dirigido
\newcommand{\dsc}{\dif S}                           % elemento de área escalar
\newcommand{\dV}{\dif V}                            % elemento de volumen
\newcommand{\defeq}{\mathrel{\overset{\mathrm{def}}{=}}}
\newcommand{\bordeS}{\partial S}
\newcommand{\bordeV}{\partial V}
\DeclareMathOperator{\grad}{grad}
\DeclareMathOperator{\dive}{div}
\DeclareMathOperator{\rot}{rot}
\DeclareMathOperator{\sen}{sen}
\DeclareMathOperator{\senh}{senh}
\DeclareMathOperator{\sgn}{sgn}
\DeclareMathOperator{\Tr}{Tr}
\renewcommand{\sin}{\sen}                            % «sen» en todo el texto

% Integral de superficie cerrada (doble integral con círculo)
\newcommand{\oiint}{\mathop{\vcenter{\hbox{$\displaystyle\oint\mkern-13mu\oint$}}}\nolimits}
\newcommand{\oiintt}{\mathop{\vcenter{\hbox{$\textstyle\oint\mkern-11mu\oint$}}}\nolimits}

% Tensor de Levi-Civita y delta de Kronecker
\newcommand{\eps}{\varepsilon}
\newcommand{\dpart}[1]{\partial_{#1}}

% Texto de relleno discreto para datos pendientes
\newcommand{\pend}[1]{\textcolor{grispizarra}{\colorbox{grisclaro}{\strut\,#1\,}}}

% =========================================================================
%        ENTORNOS FORMALES (enunciados numerados, estilo artículo)
% =========================================================================
\usepackage{amsthm}

% Enunciados en cursiva (resultados generales)
\newtheoremstyle{cienplano}{10pt}{10pt}{\itshape}{}%
  {\bfseries\color{azulprof}}{.}{0.6em}%
  {\thmname{#1}\thmnumber{ #2}\thmnote{\normalfont\upshape\ (#3)}}
% Definiciones, cálculos y ejemplos (cuerpo en redonda)
\newtheoremstyle{ciendef}{10pt}{10pt}{\normalfont}{}%
  {\bfseries\color{azulprof}}{.}{0.6em}%
  {\thmname{#1}\thmnumber{ #2}\thmnote{\normalfont\upshape\ (#3)}}
% Observaciones y notas (encabezado en cursiva)
\newtheoremstyle{cienobs}{8pt}{8pt}{\normalfont}{}%
  {\itshape\color{azulprof}}{.}{0.6em}%
  {\thmname{#1}\thmnumber{ #2}\thmnote{\normalfont\upshape\ (#3)}}

\theoremstyle{cienplano}
\newtheorem{proposicion}{Proposición}[section]
\newtheorem{teorema}[proposicion]{Teorema}
\newtheorem{lema}[proposicion]{Lema}
\newtheorem{corolario}[proposicion]{Corolario}
\theoremstyle{ciendef}
\newtheorem{definicion}[proposicion]{Definición}
\newtheorem{notacion}[proposicion]{Notación}
\newtheorem{ejemplo}[proposicion]{Ejemplo}
\newtheorem{ejercicio}[proposicion]{Ejercicio}
\theoremstyle{cienobs}
\newtheorem{observacion}[proposicion]{Observación}

% Resolución de ejercicios: misma maquinaria que «proof», con otro rótulo
\renewcommand{\qedsymbol}{\ensuremath{\blacksquare}}
\newenvironment{resolucion}[1][Resolución]
  {\begin{proof}[\itshape\color{azulprof}#1]}{\end{proof}}

\crefname{proposicion}{proposición}{proposiciones}
\Crefname{proposicion}{Proposición}{Proposiciones}
\crefname{teorema}{teorema}{teoremas}
\Crefname{teorema}{Teorema}{Teoremas}
\crefname{lema}{lema}{lemas}
\Crefname{lema}{Lema}{Lemas}
\crefname{corolario}{corolario}{corolarios}
\Crefname{corolario}{Corolario}{Corolarios}
\crefname{definicion}{definición}{definiciones}
\Crefname{definicion}{Definición}{Definiciones}
\crefname{notacion}{notación}{notaciones}
\Crefname{notacion}{Notación}{Notaciones}
\crefname{ejemplo}{ejemplo}{ejemplos}
\Crefname{ejemplo}{Ejemplo}{Ejemplos}
\crefname{ejercicio}{ejercicio}{ejercicios}
\Crefname{ejercicio}{Ejercicio}{Ejercicios}
\crefname{observacion}{observación}{observaciones}
\Crefname{observacion}{Observación}{Observaciones}
\crefname{section}{sección}{secciones}
\Crefname{section}{Sección}{Secciones}
\crefname{subsection}{sección}{secciones}
\Crefname{subsection}{Sección}{Secciones}
\crefname{table}{tabla}{tablas}
\Crefname{table}{Tabla}{Tablas}
\crefname{figure}{figura}{figuras}
\Crefname{figure}{Figura}{Figuras}
\crefname{equation}{ecuación}{ecuaciones}
\Crefname{equation}{Ecuación}{Ecuaciones}
\crefformat{equation}{ecuación~(#2#1#3)}
\Crefformat{equation}{Ecuación~(#2#1#3)}
\crefrangeformat{equation}{ecuaciones~(#3#1#4)--(#5#2#6)}
\crefmultiformat{equation}{ecuaciones~(#2#1#3)}{ y~(#2#1#3)}{, (#2#1#3)}{ y~(#2#1#3)}

% Pies de figura y tabla, en la línea del documento de referencia
\usepackage[font=small, labelfont={sc,color=azulprof}, labelsep=period,
            margin=0.6cm, justification=justified]{caption}

% =========================================================================
%        ESTILOS GRÁFICOS COMUNES (TikZ / pgfplots)
% =========================================================================
% =========================================================================
%  Estilos y macros gráficas comunes
%  - \mapcolor{t}: escala cromática «campo» (t en [0,1]); define el color mapc
%  - \flechas{u}{v}{lista x}{lista y}{escala}{|F|max}: campo vectorial 2D
%  - \flechasIII{u}{v}{w}{...}: campo vectorial 3D (con tikz-3dplot)
% =========================================================================

% --- Escala cromática (5 tramos entre cm0 y cm4) ------------------------
\newcommand{\mapcolor}[1]{%
  \pgfmathsetmacro{\mt}{min(max(#1,0),0.9999)*4}%
  \pgfmathtruncatemacro{\mi}{floor(\mt)}%
  \pgfmathtruncatemacro{\mj}{\mi+1}%
  \pgfmathsetmacro{\mf}{100*(\mt-\mi)}%
  \colorlet{mapc}{cm\mj!\mf!cm\mi}}

% --- Estilos TikZ generales ------------------------------------------------
\tikzset{
  >=Stealth,
  eje/.style={->, line width=0.6pt, grispizarra},
  ejeprof/.style={->, line width=0.7pt, azulprof},
  vect/.style={->, line width=1pt, azulprof},
  vectrojo/.style={->, line width=1pt, rojoteo},
  vectverde/.style={->, line width=1pt, verdeteo},
  fino/.style={line width=0.4pt, grispizarra},
  punteado/.style={densely dotted, line width=0.5pt, grispizarra},
  etiq/.style={font=\small, text=black!85},
  cuadro/.style={draw=grispizarra!60, line width=0.4pt, fill=white, inner sep=2pt},
  trayecto/.style={line width=1.1pt, rojoteo},
  orienta/.style={decoration={markings, mark=at position 0.55 with {\arrow{Stealth[length=5pt]}}}, postaction={decorate}},
}

% --- Marco con marcas para gráficas de campos 2D ---------------------------
% \marco2D{xmin}{xmax}{ymin}{ymax}{xticks}{yticks}  (unidades TikZ = unidades del campo)
\newcommand{\marcoD}[6]{%
  \draw[grispizarra, line width=0.5pt] (#1,#3) rectangle (#2,#4);
  \foreach \t in {#5}{%
     \draw[grispizarra, line width=0.4pt] (\t,#3) -- ++(0,0.08);
     \draw[grispizarra, line width=0.4pt] (\t,#4) -- ++(0,-0.08);
     \node[below=1pt, font=\scriptsize] at (\t,#3) {\pgfmathprintnumber[fixed,precision=1,use comma,zerofill=false]{\t}};}%
  \foreach \t in {#6}{%
     \draw[grispizarra, line width=0.4pt] (#1,\t) -- ++(0.08,0);
     \draw[grispizarra, line width=0.4pt] (#2,\t) -- ++(-0.08,0);
     \node[left=1pt, font=\scriptsize] at (#1,\t) {\pgfmathprintnumber[fixed,precision=1,use comma,zerofill=false]{\t}};}%
}

% --- Campo vectorial plano: flechas centradas en cada nudo de la malla ---
% #1,#2: componentes u(x,y), v(x,y) (en términos de \x e \y)
% #3,#4: listas de abscisas y ordenadas; #5: factor de escala de longitud
% #6: módulo máximo (normaliza el color)
\newcommand{\flechas}[6]{%
  \foreach \x in {#3}{\foreach \y in {#4}{%
    \pgfmathsetmacro{\fu}{#1}\pgfmathsetmacro{\fv}{#2}%
    \pgfmathsetmacro{\fm}{sqrt(\fu*\fu+\fv*\fv)}%
    \pgfmathsetmacro{\ft}{\fm/(#6)}%
    \mapcolor{\ft}%
    \ifdim\fm pt>0.001pt
      \pgfmathsetmacro{\hl}{2.2+4.2*\ft}%
      \pgfmathsetmacro{\lw}{0.35+0.5*\ft}%
      \draw[-{Stealth[length=\hl pt,width=\hl pt]}, line width=\lw pt, mapc]
        (\x-0.5*#5*\fu,\y-0.5*#5*\fv) -- ++(#5*\fu,#5*\fv);
    \else
      \fill[mapc] (\x,\y) circle (0.7pt);
    \fi}}}

% --- Campo vectorial espacial (tikz-3dplot): lo mismo en 3D ----------------
\newcommand{\flechasIII}[8]{%
  \foreach \x in {#4}{\foreach \y in {#5}{\foreach \z in {#6}{%
    \pgfmathsetmacro{\fu}{#1}\pgfmathsetmacro{\fv}{#2}\pgfmathsetmacro{\fw}{#3}%
    \pgfmathsetmacro{\fm}{sqrt(\fu*\fu+\fv*\fv+\fw*\fw)}%
    \pgfmathsetmacro{\ft}{\fm/(#8)}%
    \mapcolor{\ft}%
    \ifdim\fm pt>0.001pt
      \pgfmathsetmacro{\hl}{2.0+3.6*\ft}%
      \pgfmathsetmacro{\lw}{0.35+0.5*\ft}%
      \draw[-{Stealth[length=\hl pt,width=\hl pt]}, line width=\lw pt, mapc]
        (\x-0.5*#7*\fu,\y-0.5*#7*\fv,\z-0.5*#7*\fw) -- ++(#7*\fu,#7*\fv,#7*\fw);
    \else
      \fill[mapc] (\x,\y,\z) circle (0.7pt);
    \fi}}}}

% --- Estilo base de ejes pgfplots (como en la figura del efecto Doppler) --
\pgfplotsset{
  tema/.style={
    tick label style={font=\scriptsize}, label style={font=\small},
    axis lines=left, axis line style={grispizarra},
    grid=major, grid style={grisclaro, line width=0.4pt},
    every tick/.style={grispizarra},
    legend style={font=\scriptsize, draw=none, fill=white, fill opacity=0.85,
                  text opacity=1},
    /pgf/number format/use comma,
  },
  mapa/.style={
    tick label style={font=\scriptsize}, label style={font=\small},
    axis line style={grispizarra}, every tick/.style={grispizarra},
    /pgf/number format/use comma,
    colormap name=campo,
  },
  superficie/.style={
    tick label style={font=\scriptsize}, label style={font=\small},
    axis line style={grispizarra}, every tick/.style={grispizarra},
    grid=none, /pgf/number format/use comma,
    colormap name=campo,
  },
}

% Rótulos de ejes para marcoD:  \ejesD{xmin}{xmax}{ymin}{ymax}
\newcommand{\ejesD}[4]{%
  \node[font=\small, anchor=north] at ($({(#1+#2)/2},#3)+(0,-13pt)$) {$x$ (u.\,arb.)};
  \node[font=\small, rotate=90, anchor=south] at ($(#1,{(#3+#4)/2})+(-17pt,0)$) {$y$ (u.\,arb.)};}

% Recorta el contenido #5 al rectángulo [#1,#2]x[#3,#4]
\newcommand{\clipD}[5]{\begin{scope}\clip (#1,#3) rectangle (#2,#4); #5 \end{scope}}

% Versor dibujado en el punto (#1,#2,#3) con desplazamiento #4 (coordenadas 3D)
\newcommand{\uvect}[5]{%  versor en P: origen (#1,#2,#3), dirección (#4) y rótulo (#5)
  \draw[-{Stealth[length=4.5pt,width=4pt]}, rojoteo, line width=1.1pt] (#1,#2,#3) -- ++#4 node[#5] {};}

% Mapa de calor de f=exp(-r^2) (anillos concéntricos) con curvas de nivel y marco; #1 = extras
\newcommand{\mapaGauss}[1]{%
  \begin{scope}
    \clip (-2,-2) rectangle (2,2);
    \foreach \k in {0,...,90}{%
      \pgfmathsetmacro{\rk}{3.0*(1-\k/91)}%
      \pgfmathsetmacro{\fk}{exp(-\rk*\rk)}%
      \mapcolor{\fk}%
      \fill[mapc] (0,0) circle (\rk);}
    \foreach \lv in {0.9,0.8,0.7,0.6,0.5,0.4,0.3,0.2,0.1}{%
      \pgfmathsetmacro{\rl}{sqrt(-ln(\lv))}%
      \draw[black!75, line width=0.45pt] (0,0) circle (\rl);}
    #1
  \end{scope}
  \marcoD{-2}{2}{-2}{2}{-2,-1,0,1,2}{-2,-1,0,1,2}
  \ejesD{-2}{2}{-2}{2}}




% =========================================================================
%                     VARIABLES PARA LA PORTADA
% =========================================================================
\newcommand{\Universidad}{Universidad de Salamanca}
\newcommand{\Facultad}{Facultad de Ciencias}
\newcommand{\Titulacion}{Grado en Física}
\newcommand{\Asignatura}{Electromagnetismo I}
\newcommand{\NumTema}{Tema 0}
\newcommand{\TituloTema}{Herramientas\\matemáticas}
\newcommand{\tipotrabajo}{Apuntes formalizados}
\newcommand{\miNombre}{Daniel Martín Guerrero}
\newcommand{\miContenido}{Cálculo vectorial\\Coordenadas curvilíneas\\Teoremas integrales}
\newcommand{\miUbicacion}{Salamanca, curso 2025/2026}

% =========================================================================
%                     DATOS DEL ARTÍCULO (cabecera)
% =========================================================================
\title{Herramientas matemáticas para el electromagnetismo}
\author{Daniel Martín Guerrero}
\address{Grado en Física, Facultad de Ciencias, Universidad de Salamanca,
         Salamanca, España}
\date{}

% =========================================================================
%                                DOCUMENTO
% =========================================================================
\begin{document}
% --- PORTADA ----------------------------------------------------------------
\pagenumbering{alph}
%%%%%%%%%%%%%%%%%%%%%%%%%%%%%%%%%%%%%%%%%%%%%%%%%%%%%%%%%%%%%%%%%%%%%%%%
% Portada — estilo artículo científico, motivo: curvas de nivel y campo
% gradiente de f(x,y)=exp(-(x^2+y^2)).  Requiere (definidos en main.tex)
% los colores azulprof, azulacero, azulclaro, grispizarra, grisclaro y las
% macros \Universidad, \Facultad, \Titulacion, \Asignatura, \NumTema,
% \TituloTema, \tipotrabajo, \miNombre, \miContenido, \miUbicacion.
%%%%%%%%%%%%%%%%%%%%%%%%%%%%%%%%%%%%%%%%%%%%%%%%%%%%%%%%%%%%%%%%%%%%%%%%

\begin{titlepage}
\thispagestyle{empty}
\pagecolor{white}

\begin{tikzpicture}[remember picture, overlay]
  \coordinate (O) at (current page.south west);

  % ----------------------------------------------------------------
  % Motivo gráfico: curvas de nivel de f = exp(-rho^2/s^2) (equiespaciadas
  % en f) y campo gradiente (flechas hacia el máximo, longitud ~ |grad f|).
  % ----------------------------------------------------------------
  \begin{scope}
    \clip (O) rectangle ($(O)+(21cm,24.6cm)$);
    % curvas de nivel
    \foreach \k in {1,...,14}{
      \pgfmathsetmacro{\fk}{\k/15}
      \pgfmathsetmacro{\rr}{5.2*sqrt(-ln(\fk))}
      \pgfmathsetmacro{\pc}{100-6.2*\k}
      \pgfmathsetmacro{\lw}{0.25+0.045*\k}
      \draw[azulprof!\pc!azulclaro, line width=\lw pt]
        ($(O)+(14.2cm,17.2cm)$) circle (\rr cm);
    }
    % flechas del gradiente a lo largo de 24 rayos
    \foreach \ang in {0,15,...,345}{
      \foreach \rho in {1.6,2.9,4.3,6.1,8.3}{
        \pgfmathsetmacro{\g}{(\rho/5.2)*exp(-(\rho/5.2)*(\rho/5.2))}
        \pgfmathsetmacro{\len}{3.2*\g}
        \pgfmathsetmacro{\pc}{30+250*\g}
        \draw[-{Stealth[length=3.2pt,width=3pt]}, line width=0.5pt,
              azulacero!\pc!azulclaro]
          ($(O)+(14.2cm,17.2cm)+(\ang:\rho cm)$) -- ++(\ang+180:\len cm);
      }
    }
    \fill[azulprof] ($(O)+(14.2cm,17.2cm)$) circle (2.2pt);
  \end{scope}

  % ----------------------------------------------------------------
  % Cabecera: logotipos y filete
  % ----------------------------------------------------------------
  \node[anchor=north west, inner sep=0] at ($(O)+(2.2cm,27.9cm)$)
    {\includegraphics[height=1.42cm]{logo_usal.png}};
  \node[anchor=north east, inner sep=0] at ($(O)+(18.8cm,27.9cm)$)
    {\includegraphics[height=1.42cm]{logo_facultad.png}};
  \draw[grispizarra!60, line width=0.4pt]
    ($(O)+(2.2cm,26.0cm)$) -- ($(O)+(18.8cm,26.0cm)$);
  \node[anchor=north west, inner sep=0, text=grispizarra, font=\scshape\small]
    at ($(O)+(2.2cm,25.75cm)$)
    {\textls[80]{\Universidad\ \textbar\ \Facultad}};
  \node[anchor=north east, inner sep=0, text=grispizarra, font=\scshape\small]
    at ($(O)+(18.8cm,25.75cm)$)
    {\textls[80]{\Titulacion}};

  % ----------------------------------------------------------------
  % Bloque del título (sobre un velo blanco para que el motivo no tape el texto)
  % ----------------------------------------------------------------
  \shade[top color=white!0, bottom color=white, opacity=0.95]
    ($(O)+(0,12.3cm)$) rectangle ($(O)+(21cm,13.4cm)$);
  \fill[white, opacity=0.95] ($(O)+(0,5.35cm)$) rectangle ($(O)+(21cm,12.3cm)$);

  % Perfil de f a lo largo de un diámetro, como filete ondulado
  \begin{scope}[shift={(O)}]
    \draw[azulacero, line width=0.9pt, domain=2.2:18.8, samples=300, smooth]
      plot ({\x}, {12.72 + 0.62*exp(-((\x-10.5)/2.4)^2)});
  \end{scope}

  \node[anchor=north west, inner sep=0, text width=16.6cm, align=left]
    at ($(O)+(2.2cm,12.0cm)$)
    {{\color{azulacero}\scshape\large\textls[120]{\Asignatura\ \textbar\ \NumTema}}\\[12pt]
     {\color{azulprof}\fontsize{35}{39}\selectfont\bfseries\TituloTema}\\[14pt]
     {\color{azulacero}\rule{2.5cm}{1.4pt}}\\[9pt]
     {\color{grispizarra}\itshape\large
      $\displaystyle \int_{\vc{a}}^{\vc{b}}\nab f\cdot\dl \;=\; f(\vc{b})-f(\vc{a})$}};

  % ----------------------------------------------------------------
  % Pie: tipo de trabajo, autor, contenidos, lugar y fecha
  % ----------------------------------------------------------------
  \draw[grispizarra!60, line width=0.4pt]
    ($(O)+(2.2cm,5.1cm)$) -- ($(O)+(18.8cm,5.1cm)$);

  \node[anchor=north west, inner sep=0, text width=4.6cm, align=left]
    at ($(O)+(2.2cm,4.75cm)$)
    {{\scshape\small\color{azulacero}\textls[80]{Documento}}\\[3pt]
     {\large\bfseries\color{azulprof}\tipotrabajo}};

  \node[anchor=north west, inner sep=0, text width=5.6cm, align=left]
    at ($(O)+(7.2cm,4.75cm)$)
    {{\scshape\small\color{azulacero}\textls[80]{Autor}}\\[3pt]
     {\large\color{black!85}\miNombre}};

  \node[anchor=north west, inner sep=0, text width=5cm, align=left]
    at ($(O)+(13.8cm,4.75cm)$)
    {{\scshape\small\color{azulacero}\textls[80]{Contenidos}}\\[3pt]
     {\large\color{black!85}\miContenido}};

  \node[anchor=south east, inner sep=0, text=grispizarra, font=\scshape\small]
    at ($(O)+(18.8cm,1.45cm)$)
    {\textls[100]{\miUbicacion}};
  \fill[azulprof] ($(O)+(0,0)$) rectangle ($(O)+(21cm,0.45cm)$);
  \fill[azulacero] ($(O)+(0,0)$) rectangle ($(O)+(7cm,0.45cm)$);
  \fill[azulclaro] ($(O)+(0,0)$) rectangle ($(O)+(3cm,0.45cm)$);
\end{tikzpicture}

\end{titlepage}

\clearpage
\pagenumbering{arabic}
\setcounter{page}{1}

\begin{abstract}
Recopilamos y formalizamos las herramientas matemáticas que sustentan el
curso de Electromagnetismo~I: campos escalares y vectoriales, el operador
gradiente y las derivadas vectoriales de primer y segundo orden, los tres
teoremas integrales fundamentales del cálculo vectorial (gradiente, Stokes
y Gauss), los sistemas de coordenadas cartesianas, cilíndricas y esféricas
con los operadores diferenciales en cada uno de ellos, y las funciones de
la posición relativa $\vc{R}=\vc{r}-\vc{r}\,'$ junto con la delta de Dirac.
Cada resultado se enuncia con rigor y se acompaña de su demostración o de
una justificación geométrica; los seminarios finales resuelven los
ejercicios de verificación de los teoremas, incluida la demostración
de identidades vectoriales con el tensor de Levi-Civita.
\end{abstract}

\maketitle
\tableofcontents

% =========================================================================
\section{Introducción}\label{sec:introduccion}
% =========================================================================
El electromagnetismo clásico es, en esencia, una teoría de \emph{campos}: a
cada punto del espacio se le asigna una magnitud (un número, un vector) que
varía de un punto a otro y, en general, también con el tiempo. El campo
eléctrico $\vc{E}(\vc{r})$ y el campo magnético $\vc{B}(\vc{r})$ son campos
vectoriales; el potencial electrostático $V(\vc{r})$ y la densidad de carga
$\rho(\vc{r})$ son campos escalares. Para hablar de cómo cambian esos campos
(\emph{derivadas}), de cuánto «atraviesan» una superficie o «circulan» a lo
largo de un circuito (\emph{integrales}), y de cómo conectar ambas ideas,
necesitamos un pequeño conjunto de herramientas matemáticas, que son las que
reúne este tema~\cite{NAVAS,GRIFFITHS,ARFKEN}. Las figuras reproducen, ahora como
diagramas vectoriales en TikZ y gráficas en pgfplots, las del material de
figuras del curso~\cite{FIGURAS-T0}, que se prepararon con Mathematica, y las completan con
diagramas propios de los ejercicios de verificación.

Estas herramientas se pueden agrupar en tres bloques que se apoyan unos en
otros:
\begin{itemize}
  \item \textbf{Cálculo diferencial vectorial} (\cref{sec:campos,sec:operadores}):
        el gradiente de un campo escalar, y la divergencia y el rotacional de un
        campo vectorial, con sus reglas de derivación y las identidades que
        satisfacen entre sí.
  \item \textbf{Cálculo integral vectorial} (\cref{sec:integral}): las integrales
        de línea y de flujo, y los tres teoremas integrales fundamentales
        (del gradiente, de Stokes y de Gauss), todos ellos generalizaciones del
        teorema fundamental del cálculo.
  \item \textbf{Geometría de las coordenadas} (\cref{sec:coordenadas,sec:posicion}):
        los sistemas cartesiano, cilíndrico y esférico con sus operadores
        diferenciales, y las funciones de la posición relativa
        $\vc{R}=\vc{r}-\vc{r}\,'$ que aparecen siempre que una distribución
        de fuentes genera un campo, junto con la delta de Dirac.
\end{itemize}
El documento termina con unos \emph{seminarios prácticos} (\cref{sec:seminarios})
donde se resuelven, con todo detalle, los ejercicios de verificación de los
teoremas y las demostraciones de identidades vectoriales con el tensor de
Levi-Civita. Los resultados se enuncian con el rigor de un artículo, pero
procuramos acompañar cada fórmula de una explicación geométrica o física que
permita \emph{ver} por qué es cierta.

% -------------------------------------------------------------------------
\subsection{Notación y convenios}\label{sec:notacion}
% -------------------------------------------------------------------------
Seguimos la notación vectorial de L.~Navas~\cite{NAVAS} y las expresiones del
formulario del curso~\cite{FORMULARIO}. Las flechas distinguen los
vectores de los escalares, y los versores (vectores unitarios) se escriben con
sombrero.

\begin{table}[ht]
\centering
\caption{Resumen de la notación empleada.}
\label{tab:notacion}
\small
\renewcommand{\arraystretch}{1.35}
\begin{tabular}{@{}p{4.6cm}p{9.2cm}@{}}
\toprule
Símbolo & Significado\\
\midrule
$\vc{r}=x\,\ux+y\,\uy+z\,\uz$ & Vector de posición del \emph{punto campo};
  $\vc{r}\,'$ es el del \emph{punto fuente}.\\
$\ux,\uy,\uz$;\ $\ur,\uth,\uph$ & Versores de las bases cartesiana y esférica
  (en la base cilíndrica, $\urho,\uph,\uz$).\\
$f,\,g,\,V$ & Campos escalares, $f:W\subseteq\R^{n}\to\R$.\\
$\vc{F},\,\vc{A},\,\vc{B}$ & Campos vectoriales, $\vc{F}:W\subseteq\R^{n}\to\R^{n}$,
  con componentes $(F_x,F_y,F_z)$.\\
$\nab=\ux\,\partial_x+\uy\,\partial_y+\uz\,\partial_z$ & Operador «nabla»
  (o «del»), un operador diferencial; $\nab'$ actúa sobre las coordenadas fuente.\\
$\nab f,\ \nab\cdot\vc{F},\ \nab\times\vc{F},\ \Lap f$ & Gradiente, divergencia,
  rotacional y laplaciano.\\
$\dl,\ \dS,\ \dV$ & Elemento de línea dirigido, de área dirigido
  ($\dS=\un\,\dsc$) y de volumen.\\
$\partial S,\ \partial V$ & Borde (orientado) de una superficie $S$ y de un
  volumen $V$.\\
$\displaystyle\int_{C}\!,\ \oint_{C},\ \iint_{S},\ \oiint_{S},\ \iiint_{V}$ &
  Integrales sobre curvas abiertas y cerradas, sobre superficies abiertas y
  cerradas, y sobre volúmenes.\\
$\varepsilon_{ijk},\ \delta_{ij}$ & Tensor de Levi-Civita y delta de Kronecker.\\
$\delta^{3}(\vc{r})$ & Delta de Dirac tridimensional.\\
\bottomrule
\end{tabular}
\end{table}

\begin{notacion}[Convenios generales]\label{not:convenios}
Fijamos los siguientes convenios.\begin{enumerate}[label=(\alph*)]
  \item Salvo mención en contrario, los campos son de clase $C^{1}$ en la región
        considerada (de clase $C^{2}$ cuando intervienen derivadas segundas), y
        las curvas y superficies son regulares a trozos.
  \item Las \emph{regiones} $W\subseteq\R^{n}$ son abiertos conexos.
  \item Las curvas y las superficies están \emph{orientadas}. El borde de una
        superficie lleva la orientación compatible con la de la superficie
        según la regla de la mano derecha (\cref{sec:stokes}), y el borde de
        un sólido se orienta con el vector normal exterior.
  \item En los desarrollos con índices se aplica el convenio de sumación de
        Einstein: todo índice repetido en un mismo monomio se suma de $1$ a $3$,
        y $\partial_i\equiv\partial/\partial x_i$ con
        $(x_1,x_2,x_3)=(x,y,z)$.
  \item Escribimos $\sen$ para el seno y $\tan^{-1}$ para la función arcotangente.
\end{enumerate}
\end{notacion}

\begin{observacion}[Sobre el símbolo $\nab$]\label{obs:nabla}
El símbolo $\nab$ se manipula como si fuese un vector de componentes
$(\partial_x,\partial_y,\partial_z)$, y esa analogía es muy útil para recordar
fórmulas: $\nab f$ «es» el vector $\nab$ multiplicado por el escalar $f$,
$\nab\cdot\vc{F}$ su producto escalar con $\vc{F}$, y $\nab\times\vc{F}$ su
producto vectorial. Pero conviene no olvidar que $\nab$ \emph{no es un vector}
hasta que opera sobre un campo: es un operador diferencial, y en particular no
conmuta con los campos ($\nab\cdot\vc{A}\neq\vc{A}\cdot\nab$). La analogía
sugiere fórmulas, pero no las demuestra.
\end{observacion}

% =========================================================================
\section{Campos escalares, diferencial y gradiente}\label{sec:campos}
% =========================================================================

% -------------------------------------------------------------------------
\subsection{Campos escalares y su diferencial}\label{sec:escalares}
% -------------------------------------------------------------------------
\begin{definicion}[Campo escalar]\label{def:escalar}
Sea $W\subseteq\R^{n}$ una región. Un \emph{campo escalar} es una función
$f:W\to\R$ que asigna a cada punto $\vc{r}=(x,y,z)\in W$ un número
$f(\vc{r})=f(x,y,z)$.
\end{definicion}

Son campos escalares la temperatura $T(\vc{r})$ en el interior de una
habitación, la altura $h(x,y)$ sobre el nivel del mar en un mapa topográfico
(aquí $n=2$) o el potencial electrostático $V(\vc{r})$. En una sola
variable, la diferencial de $f$ recoge cuánto cambia la función cuando la
variable cambia en una cantidad infinitesimal:
\begin{equation}\label{eq:df1d}
  \dif f=\frac{\dif f}{\dif x}\,\dif x
\end{equation}
(\cref{fig:diferencial1d}): es el incremento a lo largo de la \emph{recta
tangente}, de pendiente $\tan\alpha=f'(x)$. Para un campo escalar de tres
variables hay tres direcciones independientes en las que moverse y la
diferencial total es la suma de las tres contribuciones,
\begin{equation}\label{eq:df3d}
  \dif f=\pd{f}{x}\,\dif x+\pd{f}{y}\,\dif y+\pd{f}{z}\,\dif z .
\end{equation}

\begin{figure}[ht]
\centering
\begin{tikzpicture}[>=Stealth, x=1cm, y=1cm]
  % ejes
  \draw[eje] (-0.3,0) -- (7.6,0) node[right, text=black] {$x$};
  \draw[eje] (0,-0.3) -- (0,4.6) node[above, text=black] {$f$};
  % curva
  \draw[azulprof, line width=1.2pt, domain=0.4:7.2, samples=100, smooth]
    plot (\x, {0.55+3.6/(1+exp(-1.35*(\x-3.7)))});
  % punto, tangente y triángulo diferencial
  \def\xo{3.2}
  \pgfmathsetmacro{\yo}{0.55+3.6/(1+exp(-1.35*(\xo-3.7)))}
  \pgfmathsetmacro{\sl}{3.6*1.35*exp(-1.35*(\xo-3.7))/(1+exp(-1.35*(\xo-3.7)))^2}
  \draw[azulacero, dashed, line width=0.7pt, domain=\xo-1.7:\xo+2.4]
    plot (\x, {\yo+\sl*(\x-\xo)});
  \fill[azulprof] (\xo,\yo) circle (2pt);
  \draw[grispizarra, densely dashed] (\xo,0) -- (\xo,\yo) -- (0,\yo);
  \draw[rojoteo, line width=0.9pt] (\xo,\yo) -- node[below, text=rojoteo] {$\dif x$} (\xo+1.3,\yo);
  \draw[rojoteo, line width=0.9pt] (\xo+1.3,\yo) -- node[right, text=rojoteo] {$\dif f$}
        (\xo+1.3,{\yo+1.3*\sl});
  \draw[grispizarra, densely dashed] (\xo+1.3,0) -- (\xo+1.3,\yo);
  \node[below, font=\small] at (\xo,0) {$x$};
  \node[below, font=\small] at (\xo+1.3,0) {$x+\dif x$};
  \draw[grispizarra, thick] (\xo+0.75,\yo) arc (0:{atan(\sl)}:0.75);
  \node[font=\small, text=black!80] at (\xo+1.18,\yo+0.25) {$\alpha$};
  \node[azulacero, font=\small, anchor=north east] at (\xo-0.55,{\yo-0.55*\sl-0.1}) {$\tan\alpha=f'(x)$};
  \node[azulprof, font=\small, anchor=south] at (6.7,4.3) {$f(x)$};
\end{tikzpicture}
\caption{Diferencial de una función de una variable: $\dif f=f'(x)\,\dif x$
es el incremento a lo largo de la tangente, y coincide con el incremento real
de $f$ salvo infinitésimos de orden superior.}
\label{fig:diferencial1d}
\end{figure}

Si escribimos el desplazamiento infinitesimal
$\dl=\dif x\,\ux+\dif y\,\uy+\dif z\,\uz$ (el \emph{elemento de línea}, un
vector), la \cref{eq:df3d} puede leerse como un producto escalar:
\begin{equation}\label{eq:df-gradiente}
  \dif f=\Bigl(\pd{f}{x}\,\ux+\pd{f}{y}\,\uy+\pd{f}{z}\,\uz\Bigr)\cdot
         \bigl(\dif x\,\ux+\dif y\,\uy+\dif z\,\uz\bigr).
\end{equation}
El primer factor es el protagonista de esta sección.

\begin{definicion}[Gradiente]\label{def:gradiente}
Sea $f:W\subseteq\R^{3}\to\R$ un campo escalar derivable. Su \emph{gradiente}
es el campo vectorial
\[
  \nab f\equiv\grad f=\pd{f}{x}\,\ux+\pd{f}{y}\,\uy+\pd{f}{z}\,\uz,
\]
y la diferencial se escribe, de forma compacta e independiente de la
descripción por componentes,
\begin{equation}\label{eq:df-compacta}
  \boxed{\;\dif f=\nab f\cdot\dl\;}
\end{equation}
\end{definicion}

\begin{observacion}[El gradiente no depende de las coordenadas]\label{obs:grad-intrinseco}
La fórmula de la \cref{def:gradiente} parece depender de los ejes elegidos, pero
no es así: el gradiente es el único vector que satisface, para todo vector
$\vc{v}$, la relación $\nab f(P)\cdot\vc{v}=D_{\vc{v}}f(P)$ con la derivada
direccional que definimos a continuación~\cite{NAVAS}. En
física se acostumbra a escribir los campos conservativos como $\vc{F}=-\nab f$
cuando $f$ es una energía potencial; aquí no cambiamos el signo.
\end{observacion}

% -------------------------------------------------------------------------
\subsection{Derivada direccional y significado geométrico}\label{sec:direccional}
% -------------------------------------------------------------------------
Necesitamos primero una versión de la regla de la cadena: cómo cambia $f$ si
nos movemos a lo largo de una curva.

\begin{proposicion}[Derivada de un campo escalar a lo largo de una curva]\label{prop:cadena}
Sea $f$ un campo escalar de clase $C^{1}$ y $\gamma:[a,b]\to W$ una curva
derivable. Entonces
\begin{equation}\label{eq:cadena}
  (f\circ\gamma)'(t)=\nab f\bigl(\gamma(t)\bigr)\cdot\gamma'(t).
\end{equation}
\end{proposicion}
\begin{proof}
La matriz jacobiana de $f$ es la matriz fila $\bigl(\partial_xf,\partial_yf,\partial_zf\bigr)$,
la de $\gamma$ es la columna $\bigl(x'(t),y'(t),z'(t)\bigr)^{\!\top}$ y la de la
composición $f\circ\gamma$ es la matriz $1\times1$ formada por su derivada
ordinaria. La regla de la cadena multivariable dice que la jacobiana de la
composición es el producto de las dos primeras, es decir,
$(f\circ\gamma)'=\partial_xf\,x'+\partial_yf\,y'+\partial_zf\,z'$, que es
$\nab f\cdot\gamma'$.
\end{proof}

\begin{definicion}[Derivada direccional]\label{def:direccional}
Sea $\hat u$ un vector unitario. La \emph{derivada direccional} de $f$ en el
punto $P$ en la dirección $\hat u$ es el ritmo de variación de $f$ al avanzar
sobre la recta que pasa por $P$ con dirección $\hat u$:
\[
  D_{\hat u}f(P)\defeq\lim_{h\to0}\frac{f(P+h\,\hat u)-f(P)}{h}.
\]
\end{definicion}

\begin{proposicion}\label{prop:direccional}
$D_{\hat u}f(P)=\nab f(P)\cdot\hat u=\norm{\nab f(P)}\cos\epsilon$, donde
$\epsilon$ es el ángulo que forman $\nab f(P)$ y $\hat u$.
\end{proposicion}
\begin{proof}
Aplicamos la \cref{prop:cadena} a la recta $\gamma(t)=P+t\,\hat u$, cuya
velocidad es $\gamma'(t)=\hat u$: $D_{\hat u}f(P)=(f\circ\gamma)'(0)=\nab f(P)\cdot\hat u$.
La segunda igualdad es la definición geométrica del producto escalar con
$\norm{\hat u}=1$.
\end{proof}

\begin{proposicion}[Dirección de máximo crecimiento]\label{prop:maximo}
Sea $\nab f(P)\neq\vc{0}$. Entre todas las direcciones unitarias,
\[
  \max_{\norm{\hat u}=1}D_{\hat u}f(P)=\norm{\nab f(P)}
  \qquad\text{y se alcanza si y solo si}\qquad
  \hat u=\frac{\nab f(P)}{\norm{\nab f(P)}}.
\]
\end{proposicion}
\begin{proof}
Por la desigualdad de Cauchy--Schwarz,
$D_{\hat u}f(P)=\nab f(P)\cdot\hat u\le\norm{\nab f(P)}\,\norm{\hat u}=\norm{\nab f(P)}$,
con igualdad exactamente cuando $\hat u$ es paralelo y del mismo sentido que
$\nab f(P)$, es decir, cuando $\cos\epsilon=\cos0=1$.
\end{proof}

\begin{observacion}[Las dos preguntas que responde el gradiente]\label{obs:dos-preguntas}
La \cref{prop:maximo} dice que en cada punto el gradiente contesta a dos
preguntas a la vez:
\[
  \nab f\;\longrightarrow\;
  \begin{cases}
    \text{\emph{dirección} en la que $f$ crece más deprisa,}\\[2pt]
    \norm{\nab f}=\text{\emph{valor} de esa tasa máxima de crecimiento.}
  \end{cases}
\]
En la dirección opuesta, $-\nab f$, la función decrece lo más rápido posible,
y en las direcciones perpendiculares a $\nab f$ no cambia (a primer orden).
En términos de la \cref{eq:df-compacta}: $\dif f=\norm{\nab f}\,\norm{\dl}\cos\theta$ (con $\theta$ el ángulo entre $\nab f$ y $\dl$),
y por tanto, para un desplazamiento de módulo fijo, $\dif f$ es máximo cuando
$\theta=0$, esto es, cuando $\dl$ es paralelo a $\nab f$ y entonces
$\dif f=\norm{\nab f}\,\dif\ell$.
\end{observacion}

\begin{proposicion}[El gradiente es ortogonal a las superficies de nivel]\label{prop:ortogonal}
Sea $c\in\R$ y $\Sigma_c=\{\vc{r}\in W: f(\vc{r})=c\}$ la superficie de nivel
(\emph{equipotencial}, o curva de nivel si $n=2$). Si $\gamma$ es una curva
contenida en $\Sigma_c$ con $\gamma(t_0)=P$, entonces
$\nab f(P)\cdot\gamma'(t_0)=0$: el gradiente es perpendicular a $\Sigma_c$ en $P$.
\end{proposicion}
\begin{proof}
Como $f(\gamma(t))=c$ para todo $t$, su derivada es nula. Por la
\cref{prop:cadena}, $0=(f\circ\gamma)'(t_0)=\nab f(P)\cdot\gamma'(t_0)$, y
$\gamma'(t_0)$ recorre todos los vectores tangentes a $\Sigma_c$ en $P$.
\end{proof}

Geométricamente (\cref{fig:gradiente-geo,fig:gradiente-superficie}): si $h(x,y)$ es la altura de un
terreno, sus curvas de nivel son las líneas de «misma cota» del mapa
topográfico, y el gradiente apunta cuesta arriba, cortando a esas curvas en
ángulo recto. Es la línea de máxima pendiente: la que seguiría una gota de agua
(en sentido opuesto) al deslizarse por la ladera.

% Figura: interpretación geométrica del gradiente
\begin{figure}[ht]
\centering
\begin{minipage}[t]{0.47\textwidth}\centering
\begin{tikzpicture}[>=Stealth, baseline=(current bounding box.north)]
  \coordinate (P) at (0,0);
  \draw[vect, line width=1.4pt, acento] (P) -- (4.0,1.6) node[above right, text=acento] {$\nab f$};
  \draw[vect, line width=1.2pt, azulprof] (P) -- (2.2,-0.9) node[right, text=azulprof] {$\dl$};
  % ángulo
  \draw[grispizarra, line width=0.5pt] (1.1,{1.1*1.6/4.0}) arc (21.8:-22.2:1.1);
  \node[font=\small] at (1.45,0.0) {$\theta$};
  % proyección de dl sobre grad f
  \pgfmathsetmacro{\pr}{(2.2*4.0-0.9*1.6)/(4.0*4.0+1.6*1.6)}
  \draw[densely dashed, grispizarra] (2.2,-0.9) -- ({\pr*4.0},{\pr*1.6});
  \draw[line width=2.4pt, azulclaro] (P) -- ({\pr*4.0},{\pr*1.6});
  \node[font=\footnotesize, text=azulacero, rotate=21.8, anchor=south] at ({\pr*2.0},{\pr*0.8}) {$\dl\cos\theta$};
  \draw[grispizarra, line width=0.5pt] ({\pr*4.0-0.18*0.928},{\pr*1.6-0.18*0.371})
     -- ++(0.371*0.18,-0.928*0.18) -- ++(0.928*0.18,0.371*0.18);
  \fill (P) circle (1.6pt);
\end{tikzpicture}\\[4pt]{\small\scshape (a)}\; $\dif f=\norm{\nab f}\,\norm{\dl}\cos\theta$
\end{minipage}\hfill
\begin{minipage}[t]{0.47\textwidth}\centering
\begin{tikzpicture}[>=Stealth, baseline=(current bounding box.north), scale=0.9, transform shape]
  \foreach \c/\pc in {1/100,2/65,3/35}{%
    \draw[azulprof!\pc!azulclaro, line width=0.9pt] (0,0) ellipse ({1.0*\c} and {0.5*\c});}
  \fill[azulprof] (0,0) circle (1.4pt);
  
  % gradientes en puntos de las elipses (apuntan hacia el centro)
  \foreach \c/\t in {2/60,2/200,2/320}{%
    \pgfmathsetmacro{\px}{1.0*\c*cos(\t)}\pgfmathsetmacro{\py}{0.5*\c*sin(\t)}%
    \pgfmathsetmacro{\gx}{-\px/1.0}\pgfmathsetmacro{\gy}{-4*\py/1.0}% -grad(x^2/1+4y^2)/... (dir. al centro)
    \pgfmathsetmacro{\gn}{sqrt(\gx*\gx+\gy*\gy)}%
    \pgfmathsetmacro{\ux}{0.62*\gx/\gn}\pgfmathsetmacro{\uy}{0.62*\gy/\gn}%
    \draw[vect, acento, line width=1pt] (\px,\py) -- ++(\ux,\uy);
    \fill (\px,\py) circle (1pt);}
  \node[font=\footnotesize, text=black!80, anchor=west] at (3.2,1.3) {curvas de nivel};
  
  \node[font=\footnotesize, text=azulprof, anchor=west] at (3.2,-1.3) {$h_3<h_2<h_1$};
\end{tikzpicture}\\[4pt]{\small\scshape (b)}\; $\nab h\perp\{h=\text{cte}\}$
\end{minipage}
\caption{Interpretación geométrica del gradiente. (a)~El incremento
$\dif f$ es el producto de $\norm{\nab f}$ por la proyección del desplazamiento
sobre la dirección del gradiente: es máximo si $\dl\parallel\nab f$ y nulo si
$\dl\perp\nab f$. (b)~Mapa topográfico de una colina: el gradiente de la altura
$h$ es perpendicular a las curvas de nivel y apunta hacia la cumbre (máxima
pendiente).}
\label{fig:gradiente-geo}
\end{figure}


% Figura: el gradiente sobre la superficie z=f(x,y): dirección de máximo ascenso
\begin{figure}[ht]
\centering
\begin{tikzpicture}
  \begin{axis}[superficie, width=0.92\textwidth, height=0.66\textwidth,
      xlabel={$x$}, ylabel={$y$}, zlabel={$z=f(x,y)$},
      xmin=-2, xmax=2, ymin=-2, ymax=2, zmin=0, zmax=1.15,
      xtick={-2,-1,0,1,2}, ytick={-2,-1,0,1,2}, ztick={0,0.5,1},
      grid=major, grid style={grisclaro, very thin},
      axis line style={azulprof!70}, every tick/.style={azulprof!70},
      colormap name=campoclaro, point meta min=0, point meta max=1.3,
      view={-38}{26}, clip=false]
    % superficie con malla
    \addplot3[surf, shader=faceted interp, faceted color=white!40!azulclaro, line width=0.25pt,
              samples=30, domain=-2:2, opacity=0.78] {exp(-(x^2+y^2))};
    % curva de nivel por P (f = e^{-0.89}) sobre la superficie y su proyección
    \addplot3[acento, line width=1.1pt, domain=0:360, samples=90, samples y=0]
      ({0.9434*cos(x)},{0.9434*sin(x)},{0.4107});
    \addplot3[acento, densely dashed, line width=0.6pt, domain=0:360, samples=90, samples y=0]
      ({0.9434*cos(x)},{0.9434*sin(x)},{0});
    % P, su proyección y el gradiente (en el plano z=0)
    \addplot3[only marks, mark=*, mark size=2pt, azulprof] coordinates {(0.8,0.5,0.4107)};
    \addplot3[densely dotted, grispizarra, line width=0.6pt] coordinates {(0.8,0.5,0) (0.8,0.5,0.4107)};
    \draw[-{Stealth[length=6pt]}, line width=1.6pt, acento] (axis cs:0.8,0.5,0) -- (axis cs:0.376,0.235,0);
    % vector tangente a la superficie, de máximo ascenso
    \draw[-{Stealth[length=6pt]}, line width=1.6pt, azulprof] (axis cs:0.8,0.5,0.4107) -- (axis cs:0.376,0.235,0.7977);
    \node[font=\small, azulprof, anchor=east, fill=white, fill opacity=0.85, text opacity=1, inner sep=1.5pt] at (axis cs:-0.1,-0.05,1.0) {máximo ascenso};
    \node[font=\small, azulprof, anchor=south west, fill=white, fill opacity=0.85, text opacity=1, inner sep=1pt] at (axis cs:0.85,0.5,0.45) {$P$};
    \node[font=\small, acento, anchor=north east, fill=white, fill opacity=0.85, text opacity=1, inner sep=1.5pt] at (axis cs:0.3,0.1,0.0) {$\nab f(P)$};
    \node[font=\small, acento, anchor=north west, fill=white, fill opacity=0.85, text opacity=1, inner sep=1.5pt] at (axis cs:-0.5,-1.9,0.0) {curva de nivel $f=c$};
  \end{axis}
\end{tikzpicture}
\caption{Interpretación del gradiente sobre la superficie $z=f(x,y)$. En el
punto $P$, la curva de nivel (en ocre) y el gradiente $\nab f(P)$ (en el plano
del dominio) son perpendiculares. El vector tangente a la superficie con
la misma proyección horizontal apunta en la dirección de máximo ascenso y su
pendiente es $\norm{\nab f(P)}$. El gradiente \emph{no} apunta hacia abajo: en el
plano $xy$ señala hacia donde $f$ aumenta, aquí hacia el origen.}
\label{fig:gradiente-superficie}
\end{figure}


\begin{ejemplo}[Una campana gaussiana]\label{ej:gaussiana}
Sea $f(x,y)=\eul^{-(x^{2}+y^{2})}$, que se parece a una gaussiana. Con la regla
de la cadena $[g(f(x))]'=g'(f(x))f'(x)$, cada derivada parcial produce
un factor $\eul^{-(x^{2}+y^{2})}$ por la derivada del exponente:
\[
  \nab f=\pd{f}{x}\,\ux+\pd{f}{y}\,\uy
  =-2x\,\eul^{-(x^{2}+y^{2})}\,\ux-2y\,\eul^{-(x^{2}+y^{2})}\,\uy
  =-2\,\eul^{-(x^{2}+y^{2})}\,(x\,\ux+y\,\uy).
\]
Es una función con un máximo $f(0,0)=1$ en el origen que decae al alejarnos de
él (\cref{fig:gauss}). Sus curvas de nivel son circunferencias
$x^{2}+y^{2}=-\ln c$, y el gradiente apunta radialmente hacia dentro, hacia el
máximo, perpendicular a esas circunferencias, como exige la
\cref{prop:ortogonal}. Obsérvese además que $\nab f(0,0)=\vc{0}$ (en la cima la
pendiente es nula) y que $\norm{\nab f}\to0$ lejos del origen (la función se
aplana): la pendiente máxima ha de estar en algún punto intermedio, que se
calcula en el \cref{ejer:pendiente}.
\end{ejemplo}

% Figura: f(x,y)=exp(-(x^2+y^2)) y su gradiente (mapa de calor + curvas de
% nivel, superficie 3D, y el campo gradiente en ambas representaciones).
% Los mapas planos se dibujan con anillos concéntricos (f solo depende de r).
\begin{figure}[ht]
\centering
\begin{tikzpicture}[x=1.15cm, y=1.15cm]
  % (a) mapa de calor con curvas de nivel
  \begin{scope}
    \mapaGauss{%
      \foreach \lv/\px/\py in {0.9/0.111/0.305, 0.7/0.561/0.204, 0.5/0.721/-0.416,
                               0.3/0.190/-1.081, 0.2/-0.815/0.972, 0.1/-1.494/-0.264}{%
        \node[font=\tiny, fill=white, fill opacity=0.85, text opacity=1, inner sep=0.6pt]
          at (\px,\py) {\pgfmathprintnumber[fixed,precision=1,use comma]{\lv}};}}
    \node[font=\small\scshape, anchor=south] at (0,2.15) {(a)\; $f(x,y)=\eul^{-(x^{2}+y^{2})}$};
  \end{scope}
  % (c) gradiente sobre el mapa
  \begin{scope}[shift={(0,-6.1)}]
    \mapaGauss{%
      \flechas{-2*exp(-(\x*\x+\y*\y))*\x}{-2*exp(-(\x*\x+\y*\y))*\y}%
        {-2,-1.6,...,2}{-2,-1.6,...,2}{0.75}{0.86}}
    \node[font=\small\scshape, anchor=south] at (0,2.15) {(c)\; campo $\nab f$};
  \end{scope}
\end{tikzpicture}
\hfill
\begin{tikzpicture}
  \pgfplotsset{gsup/.style={superficie, width=0.58\textwidth, height=0.48\textwidth,
      xlabel={$x$}, ylabel={$y$}, zlabel={$f(x,y)$},
      xmin=-2, xmax=2, ymin=-2, ymax=2, zmin=0, zmax=1,
      xtick={-2,-1,0,1,2}, ytick={-2,-1,0,1,2}, ztick={0,0.5,1},
      grid=major, grid style={grisclaro, very thin},
      axis line style={azulprof!70}, every tick/.style={azulprof!70},
      view={-38}{26}, point meta min=0, point meta max=1,
      title style={yshift=3pt}}}
  % (b) superficie con malla
  \begin{axis}[gsup, name=B, title={\small\scshape (b)\; superficie $z=f(x,y)$}]
    \addplot3[surf, shader=faceted interp, faceted color=white!85!azulclaro, line width=0.3pt,
              samples=22, domain=-2:2, opacity=0.97] {exp(-(x^2+y^2))};
  \end{axis}
  % (d) gradiente sobre la superficie
  \begin{axis}[gsup, at={(B.south)}, anchor=north, yshift=-2.4cm,
      title={\small\scshape (d)\; $\nab f$ sobre la superficie}]
    \addplot3[surf, shader=faceted interp, faceted color=white!85!azulclaro, line width=0.3pt,
              samples=22, domain=-2:2, opacity=0.55] {exp(-(x^2+y^2))};
    \addplot3[quiver={u={-1.1*x*exp(-(x^2+y^2))}, v={-1.1*y*exp(-(x^2+y^2))},
                      w={1.1*2*(x^2+y^2)*exp(-2*(x^2+y^2))}, colored=mapped color},
              -{Stealth[length=3pt]}, line width=0.6pt, point meta={2*sqrt(x^2+y^2)*exp(-(x^2+y^2))},
              domain=-1.8:1.8, samples=8, samples y=8] {exp(-(x^2+y^2))+0.015};
  \end{axis}
\end{tikzpicture}
\caption{Campo escalar $f(x,y)=\eul^{-(x^{2}+y^{2})}$ y su gradiente
$\nab f=-2\,\eul^{-(x^{2}+y^{2})}(x\,\ux+y\,\uy)$. (a)~Mapa de calor con
curvas de nivel $f=0{,}1,\,0{,}2,\ldots,0{,}9$ (equipotenciales).
(b)~La misma función vista como superficie $z=f(x,y)$.
(c)~El gradiente apunta hacia el máximo, es ortogonal a las curvas de nivel y
su módulo se anula en el origen y decae a grandes distancias; el color
codifica $\norm{\nab f}$ (de índigo a ámbar). (d)~Sobre la superficie, los
vectores son tangentes a ella y señalan la dirección de máxima pendiente.}
\label{fig:gauss}
\end{figure}


% -------------------------------------------------------------------------
\subsection{Campos vectoriales y campos gradiente}\label{sec:vectoriales}
% -------------------------------------------------------------------------
\begin{definicion}[Campo vectorial]\label{def:vectorial}
Un \emph{campo vectorial} en una región $W\subseteq\R^{n}$ es una función
$\vc{F}:W\to\R^{n}$ que asigna a cada punto un vector,
\[
  \vc{F}(\vc{r})=F_x(\vc{r})\,\ux+F_y(\vc{r})\,\uy+F_z(\vc{r})\,\uz ,
\]
donde las componentes $F_x,F_y,F_z$ son campos escalares.
\end{definicion}

Se visualizan dibujando en cada punto una flecha cuya longitud (y, a menudo, su
color) es proporcional al módulo del campo; esa es la convención que usamos en
todas las figuras de este documento, con una escala cromática de índigo
(módulo pequeño) a ámbar (módulo grande). El campo de velocidades de un fluido,
el campo eléctrico y el gravitatorio son ejemplos físicos. El gradiente de un
campo escalar es siempre un campo vectorial, pero no todo campo vectorial es un
gradiente, lo que motiva la siguiente definición.

\begin{definicion}[Campo gradiente y potencial]\label{def:campo-gradiente}
Un campo vectorial $\vc{F}:W\to\R^{n}$ es un \emph{campo gradiente}
(o \emph{conservativo}) si existe un campo escalar $f:W\to\R$ tal que
$\vc{F}(\vc{r})=\nab f(\vc{r})$ para todo $\vc{r}\in W$. En ese caso $f$ es un
\emph{potencial} (escalar) de $\vc{F}$ en $W$.
\end{definicion}

Los potenciales generalizan la noción de antiderivada de una variable. Hay una
diferencia esencial: toda función continua de una variable en un intervalo
tiene antiderivada, pero para $n\ge2$ existen campos vectoriales continuos
\emph{sin} potencial. Ahora bien, cuando existe, el potencial es esencialmente
único.

\begin{proposicion}[Unicidad del potencial salvo constante]\label{prop:unicidad-potencial}
Si $\vc{F}=\nab f=\nab g$ en una región conexa $W$, entonces $g=f+k$ con $k$
constante.
\end{proposicion}
\begin{proof}
En un dominio conexo, una función $h$ de clase $C^{1}$ tiene gradiente nulo si y
solo si es constante. Como $\nab(f-g)=\vc{0}$, la diferencia $f-g$ es
constante en $W$.
\end{proof}

% =========================================================================
\section{Operadores diferenciales: divergencia, rotacional y laplaciano}\label{sec:operadores}
% =========================================================================

% -------------------------------------------------------------------------
\subsection{El operador nabla}\label{sec:nabla}
% -------------------------------------------------------------------------
Todo el cálculo diferencial vectorial se organiza alrededor del operador
simbólico
\begin{equation}\label{eq:nabla}
  \nab=\pd{}{x}\,\ux+\pd{}{y}\,\uy+\pd{}{z}\,\uz=\sum_{i=1}^{3}\uh{i}\,\partial_i ,
\end{equation}
que se combina con los campos de las tres maneras en que un vector puede
combinarse con otros objetos (\cref{obs:nabla}):
\begin{center}
\renewcommand{\arraystretch}{1.25}
\begin{tabular}{@{}lcll@{}}
$\vc{A}\,f$ & $\longrightarrow$ & $\nab f$ & (gradiente), aplicado a un campo escalar;\\
$\vc{A}\cdot\vc{B}$ & $\longrightarrow$ & $\nab\cdot\vc{B}$ & (divergencia), aplicado a un campo vectorial;\\
$\vc{A}\times\vc{B}$ & $\longrightarrow$ & $\nab\times\vc{B}$ & (rotacional), aplicado a un campo vectorial.
\end{tabular}
\end{center}
El gradiente ya se estudió en la \cref{sec:campos}; veamos los otros dos.

% -------------------------------------------------------------------------
\subsection{Divergencia}\label{sec:divergencia}
% -------------------------------------------------------------------------
\begin{definicion}[Divergencia]\label{def:divergencia}
Sea $\vc{F}=F_x\,\ux+F_y\,\uy+F_z\,\uz$ un campo vectorial derivable. Su
\emph{divergencia} es el campo escalar
\begin{equation}\label{eq:div}
  \nab\cdot\vc{F}(\vc{r})=\pd{F_x}{x}+\pd{F_y}{y}+\pd{F_z}{z}
  =\nab\cdot\bigl(F_x\,\ux+F_y\,\uy+F_z\,\uz\bigr).
\end{equation}
\end{definicion}

\emph{¿Qué mide?} La divergencia dice cuánto «sale» o «entra» del campo en un
entorno infinitesimal del punto. Si pensamos en $\vc{F}$ como el campo de
velocidades de un fluido, $\nab\cdot\vc{F}(\vc{r})$ mide la tendencia del fluido
a expandirse o a contraerse en $\vc{r}$. En la \cref{sec:gauss} daremos la
definición intrínseca (independiente de las coordenadas) como densidad de flujo
por unidad de volumen:
\begin{center}
\begin{tabular}{@{}lll@{}}
$\nab\cdot\vc{F}(\vc{r})>0$ & $\Rightarrow$ & el campo \emph{se expande}: $\vc{r}$ es una \emph{fuente};\\
$\nab\cdot\vc{F}(\vc{r})<0$ & $\Rightarrow$ & el campo \emph{se contrae}: $\vc{r}$ es un \emph{sumidero};\\
$\nab\cdot\vc{F}(\vc{r})=0$ & $\Rightarrow$ & campo \emph{solenoidal} (ni se expande ni se contrae).
\end{tabular}
\end{center}

\begin{ejemplo}[Divergencia de campos de referencia]\label{ej:div}
Consideremos los campos de las \cref{fig:radial,fig:campos_div}:
\begin{enumerate}[label=(\roman*)]
  \item $\vc{F}(\vc{r})=x\,\ux+y\,\uy+z\,\uz=\vc{r}$ (campo radial):
        $\nab\cdot\vc{F}=\pd{x}{x}+\pd{y}{y}+\pd{z}{z}=1+1+1=3>0$. Se expande en
        todos los puntos, porque las flechas crecen alejándose del origen.
  \item $\vc{F}_1(\vc{r})=\ux+\uy$ (campo uniforme): $\nab\cdot\vc{F}_1=0$.
  \item $\vc{F}_2(\vc{r})=x\,\ux$: $\nab\cdot\vc{F}_2=\pd{x}{x}=1>0$.
  \item $\vc{F}_3(\vc{r})=y\,\ux$: $\nab\cdot\vc{F}_3=\pd{y}{x}=0$. El módulo
        crece con $y$, pero el campo apunta según $\ux$ y no hay variación del
        módulo en esa dirección.
\end{enumerate}
\end{ejemplo}

\begin{observacion}[Cuándo hay divergencia]\label{obs:cuando-div}
Para que un campo tenga divergencia no nula, la componente $F_i$ tiene que
\emph{variar con respecto a la coordenada $x_i$ correspondiente}: el módulo
del campo debe cambiar en la dirección hacia la que apunta. Si el módulo solo
cambia en una dirección perpendicular al campo (como en $\vc{F}_3=y\,\ux$), la
divergencia es nula: lo que entra por un lado, sale por el otro.
\end{observacion}

% Figura: campo radial F(r)= x ux + y uy + z uz en 3D, en tres planos de
% corte y en el plano xy.
\newcommand{\flechaIII}[6]{%  componentes de posición #1,#2,#3 y del campo #4,#5,#6
  \pgfmathsetmacro{\fm}{sqrt((#4)*(#4)+(#5)*(#5)+(#6)*(#6))}%
  \ifdim\fm pt>0.01pt
    \pgfmathsetmacro{\ft}{\fm/3.465}\mapcolor{\ft}%
    \pgfmathsetmacro{\hl}{2.0+3.4*\ft}%
    \draw[-{Stealth[length=\hl pt,width=\hl pt]}, line width={0.35+0.5*\ft}, mapc]
      (#1-0.17*#4,#2-0.17*#5,#3-0.17*#6) -- ++(0.34*#4,0.34*#5,0.34*#6);
  \fi}
\begin{figure}[ht]
\centering
\tdplotsetmaincoords{68}{128}
\begin{minipage}[t]{0.325\textwidth}\centering
\begin{tikzpicture}[tdplot_main_coords, scale=0.62, baseline=(current bounding box.north)]
  \draw[eje] (0,0,0) -- (3.0,0,0) node[below left] {\small $x$};
  \draw[eje] (0,0,0) -- (0,3.0,0) node[right] {\small $y$};
  \draw[eje] (0,0,0) -- (0,0,3.0) node[above] {\small $z$};
  \foreach \z in {-2,-1,0,1,2}{\foreach \y in {2,1,0,-1,-2}{\foreach \x in {-2,-1,0,1,2}{%
    \flechaIII{\x}{\y}{\z}{\x}{\y}{\z}}}}
\end{tikzpicture}\\[2pt]{\small\scshape (a)}
\end{minipage}\hfill
\begin{minipage}[t]{0.325\textwidth}\centering
\begin{tikzpicture}[tdplot_main_coords, scale=0.62, baseline=(current bounding box.north)]
  \fill[grispizarra!22, opacity=0.55] (-2.7,-2.7,0) -- (2.7,-2.7,0) -- (2.7,2.7,0) -- (-2.7,2.7,0) -- cycle;
  \fill[grispizarra!22, opacity=0.55] (0,-2.7,-2.7) -- (0,2.7,-2.7) -- (0,2.7,2.7) -- (0,-2.7,2.7) -- cycle;
  \fill[grispizarra!22, opacity=0.55] (-2.7,0,-2.7) -- (2.7,0,-2.7) -- (2.7,0,2.7) -- (-2.7,0,2.7) -- cycle;
  \draw[eje] (0,0,0) -- (3.0,0,0) node[below left] {\small $x$};
  \draw[eje] (0,0,0) -- (0,3.0,0) node[right] {\small $y$};
  \draw[eje] (0,0,0) -- (0,0,3.0) node[above] {\small $z$};
  \foreach \x in {-2,-1,0,1,2}{\foreach \y in {-2,-1,0,1,2}{\flechaIII{\x}{\y}{0}{\x}{\y}{0}}}
  \foreach \y in {-2,-1,0,1,2}{\foreach \z in {-2,-1,0,1,2}{\flechaIII{0}{\y}{\z}{0}{\y}{\z}}}
  \foreach \x in {-2,-1,0,1,2}{\foreach \z in {-2,-1,0,1,2}{\flechaIII{\x}{0}{\z}{\x}{0}{\z}}}
\end{tikzpicture}\\[2pt]{\small\scshape (b)}
\end{minipage}\hfill
\begin{minipage}[t]{0.325\textwidth}\centering
\begin{tikzpicture}[x=0.50cm, y=0.50cm, baseline=(current bounding box.north)]
  \marcoD{-2.7}{2.7}{-2.7}{2.7}{-2,-1,0,1,2}{-2,-1,0,1,2}
  \ejesD{-2.7}{2.7}{-2.7}{2.7}
  \clipD{-2.7}{2.7}{-2.7}{2.7}{\flechas{\x}{\y}{-2.4,-1.8,...,2.4}{-2.4,-1.8,...,2.4}{0.5}{3.4}}
\end{tikzpicture}\\[2pt]{\small\scshape (c)}
\end{minipage}
\caption{Campo radial $\vc{F}(\vc{r})=x\,\ux+y\,\uy+z\,\uz=\vc{r}$.
(a)~Representación tridimensional: las flechas se alejan del origen y su
módulo crece linealmente con la distancia (color de índigo a ámbar).
(b)~Cortes por los tres planos coordenados. (c)~Corte por el plano $z=0$.
Es el prototipo de campo con divergencia positiva ($\nab\cdot\vc{F}=3$):
el origen es una fuente.}
\label{fig:radial}
\end{figure}

% Figura: tres campos planos de referencia: uniforme, F=x ux (divergencia
% positiva) y F=y ux (cizalla: rotacional no nulo, divergencia nula).
\begin{figure}[ht]
\centering
\begin{minipage}[t]{0.325\textwidth}\centering
\begin{tikzpicture}[x=0.68cm, y=0.68cm, baseline=(current bounding box.north)]
  \marcoD{-2.7}{2.7}{-2.7}{2.7}{-2,-1,0,1,2}{-2,-1,0,1,2}
  \ejesD{-2.7}{2.7}{-2.7}{2.7}
  \foreach \x in {-2.4,-1.8,...,2.4}{\foreach \y in {-2.4,-1.8,...,2.4}{%
    \draw[-{Stealth[length=4.5pt,width=4.5pt]}, line width=0.7pt, azulprof]
      (\x-0.17,\y-0.17) -- ++(0.34,0.34);}}
\end{tikzpicture}\\[3pt]{\small $\vc{F}_1=\ux+\uy$}
\end{minipage}\hfill
\begin{minipage}[t]{0.325\textwidth}\centering
\begin{tikzpicture}[x=1.12cm, y=1.12cm, baseline=(current bounding box.north)]
  \marcoD{0}{3.2}{0}{3.2}{0,1,2,3}{0,1,2,3}
  \ejesD{0}{3.2}{0}{3.2}
  \clipD{0}{3.2}{0}{3.2}{\flechas{\x}{0}{0.2,0.575,...,3.1}{0.2,0.575,...,3.1}{0.30}{3.1}}
\end{tikzpicture}\\[3pt]{\small $\vc{F}_2=x\,\ux$}
\end{minipage}\hfill
\begin{minipage}[t]{0.325\textwidth}\centering
\begin{tikzpicture}[x=1.12cm, y=1.12cm, baseline=(current bounding box.north)]
  \marcoD{0}{3.2}{0}{3.2}{0,1,2,3}{0,1,2,3}
  \ejesD{0}{3.2}{0}{3.2}
  \clipD{0}{3.2}{0}{3.2}{\flechas{\y}{0}{0.2,0.575,...,3.1}{0.2,0.575,...,3.1}{0.30}{3.1}}
\end{tikzpicture}\\[3pt]{\small $\vc{F}_3=y\,\ux$}
\end{minipage}
\caption{Tres campos vectoriales planos. $\vc{F}_1=\ux+\uy$ es uniforme:
ni diverge ni rota. En $\vc{F}_2=x\,\ux$ el módulo crece \emph{en la dirección
del campo}, lo que da $\nab\cdot\vc{F}_2=1>0$ (expansión) y
$\nab\times\vc{F}_2=\vc{0}$. En $\vc{F}_3=y\,\ux$ el módulo crece
\emph{perpendicularmente} al campo: $\nab\cdot\vc{F}_3=0$ pero
$\nab\times\vc{F}_3=-\uz$ (cizalla). El color codifica $\norm{\vc{F}}$.}
\label{fig:campos_div}
\end{figure}


% -------------------------------------------------------------------------
\subsection{Rotacional}\label{sec:rotacional}
% -------------------------------------------------------------------------
\begin{definicion}[Rotacional]\label{def:rotacional}
El \emph{rotacional} de un campo vectorial derivable $\vc{F}$ es el campo
vectorial
\begin{equation}\label{eq:rot}
  \nab\times\vc{F}\equiv\rot\vc{F}=
  \begin{vmatrix}
    \ux & \uy & \uz\\[2pt]
    \partial_x & \partial_y & \partial_z\\[2pt]
    F_x & F_y & F_z
  \end{vmatrix}
  =\Bigl(\pd{F_z}{y}-\pd{F_y}{z}\Bigr)\ux
  +\Bigl(\pd{F_x}{z}-\pd{F_z}{x}\Bigr)\uy
  +\Bigl(\pd{F_y}{x}-\pd{F_x}{y}\Bigr)\uz ,
\end{equation}
donde el determinante se desarrolla por la primera fila haciendo que las
derivadas (a la izquierda) actúen sobre las componentes (a la derecha). Un campo
de rotacional nulo se llama \emph{irrotacional}.
\end{definicion}

\emph{¿Qué mide?} Si sumergimos en el campo una pequeña ruedecita con
paletas, el rotacional mide la tendencia del campo a hacerla girar:
la \emph{dirección} de $\nab\times\vc{F}$ es el eje de giro (con sentido dado
por la regla de la mano derecha) y su \emph{módulo} es proporcional a la
velocidad angular de la rueda. La definición precisa como densidad de
circulación se da en la \cref{sec:stokes}.

\begin{ejemplo}[Rotacional de campos de referencia]\label{ej:rot}
\begin{enumerate}[label=(\roman*)]
  \item \emph{Giro rígido.} $\vc{F}(\vc{r})=-y\,\ux+x\,\uy$. Solo sobrevive la
        componente $z$ de la \cref{eq:rot}:
        \[
          \nab\times\vc{F}=\Bigl[\pd{(x)}{x}-\pd{(-y)}{y}\Bigr]\uz=[1-(-1)]\,\uz=2\,\uz .
        \]
        El campo «gira» en sentido antihorario, y el rotacional es un vector
        que apunta a lo largo del eje de giro (\cref{fig:campos_rot}, izquierda).
  \item \emph{Cizalla.} $\vc{F}(\vc{r})=y\,\ux$:
        \[
          \nab\times\vc{F}=\Bigl[\pd{(0)}{x}-\pd{(y)}{y}\Bigr]\uz=-\uz .
        \]
        No hay líneas de corriente cerradas, pero el módulo del campo crece en
        la dirección \emph{perpendicular} al campo, y eso basta para hacer girar
        la rueda (en sentido horario).
  \item \emph{Campo radial.} $\nab\times\vc{r}=\vc{0}$: es irrotacional (en
        efecto, $\vc{r}=\nab(r^{2}/2)$ es un gradiente; véase la
        \cref{prop:fundamentales}).
\end{enumerate}
\end{ejemplo}

\begin{observacion}[Cuándo hay rotacional]\label{obs:cuando-rot}
Complementariamente a la \cref{obs:cuando-div}: el rotacional es no nulo
cuando el módulo de una componente del campo cambia en una dirección
\emph{perpendicular} a ella, o cuando el campo gira alrededor de un eje.
\end{observacion}

% Figura: campos con rotacional: F=-y ux + x uy (giro) y F=y ux (cizalla).
\begin{figure}[ht]
\centering
\begin{minipage}[t]{0.47\textwidth}\centering
\begin{tikzpicture}[x=0.58cm, y=0.58cm, baseline=(current bounding box.north)]
  \marcoD{-2.7}{2.7}{-2.7}{2.7}{-2,-1,0,1,2}{-2,-1,0,1,2}
  \ejesD{-2.7}{2.7}{-2.7}{2.7}
  \clipD{-2.7}{2.7}{-2.7}{2.7}{\flechas{-\y}{\x}{-2.4,-1.8,...,2.4}{-2.4,-1.8,...,2.4}{0.34}{3.4}}
\end{tikzpicture}\\[3pt]{\small $\vc{F}(\vc{r})=-y\,\ux+x\,\uy$}
\end{minipage}\hfill
\begin{minipage}[t]{0.47\textwidth}\centering
\begin{tikzpicture}[x=0.58cm, y=0.58cm, baseline=(current bounding box.north)]
  \marcoD{-2.7}{2.7}{-2.7}{2.7}{-2,-1,0,1,2}{-2,-1,0,1,2}
  \ejesD{-2.7}{2.7}{-2.7}{2.7}
  \clipD{-2.7}{2.7}{-2.7}{2.7}{\flechas{\y}{0}{-2.4,-1.8,...,2.4}{-2.4,-1.8,...,2.4}{0.34}{2.5}}
\end{tikzpicture}\\[3pt]{\small $\vc{F}(\vc{r})=y\,\ux$}
\end{minipage}
\caption{Campos con rotacional no nulo. (Izquierda) Giro rígido
$\vc{F}=-y\,\ux+x\,\uy$ (sentido antihorario): $\nab\times\vc{F}=2\,\uz$. Una
pequeña rueda colocada en cualquier punto gira con velocidad angular
constante. (Derecha) Cizalla $\vc{F}=y\,\ux$, de rotacional $-\uz$: no hay
líneas de corriente cerradas, pero la rueda gira igualmente (sentido horario)
porque el flujo es más rápido por un lado que por el otro.}
\label{fig:campos_rot}
\end{figure}


% -------------------------------------------------------------------------
\subsection{Derivadas de productos}\label{sec:productos}
% -------------------------------------------------------------------------
Al tratarse de derivadas, cabe esperar versiones de la regla del producto.
Con dos campos escalares $f,g$ y dos vectoriales $\vc{A},\vc{B}$ se pueden
formar dos productos que dan un escalar ($fg$ y $\vc{A}\cdot\vc{B}$) y dos que
dan un vector ($f\vc{A}$ y $\vc{A}\times\vc{B}$), y sobre cada uno actúa el
operador $\nab$ de la manera que corresponde a su carácter (gradiente sobre
escalares, divergencia y rotacional sobre vectores).

\begin{notacion}[El operador $\vc{A}\cdot\nab$]\label{not:A-nabla}
Dado un campo vectorial $\vc{A}$, el operador escalar
$\vc{A}\cdot\nab=A_x\,\partial_x+A_y\,\partial_y+A_z\,\partial_z$ puede aplicarse
a un campo escalar o a un campo vectorial (componente a componente). Es distinto
de $\nab\cdot\vc{A}$ (\cref{obs:nabla}).
\end{notacion}

\begin{proposicion}[Reglas de derivación del producto]\label{prop:productos}
Sean $f,g$ campos escalares y $\vc{A},\vc{B}$ campos vectoriales de clase
$C^{1}$. Entonces
\begin{align}
  \nab(fg)&=f\,\nab g+g\,\nab f, \label{eq:prod-a}\\
  \nab(\vc{A}\cdot\vc{B})&=\vc{A}\times(\nab\times\vc{B})+\vc{B}\times(\nab\times\vc{A})
     +(\vc{A}\cdot\nab)\vc{B}+(\vc{B}\cdot\nab)\vc{A}, \label{eq:prod-b}\\
  \nab\cdot(f\vc{A})&=f\,\nab\cdot\vc{A}+\vc{A}\cdot\nab f, \label{eq:prod-c}\\
  \nab\cdot(\vc{A}\times\vc{B})&=\vc{B}\cdot(\nab\times\vc{A})-\vc{A}\cdot(\nab\times\vc{B}), \label{eq:prod-d}\\
  \nab\times(f\vc{A})&=\nab f\times\vc{A}+f\,\nab\times\vc{A}, \label{eq:prod-e}\\
  \nab\times(\vc{A}\times\vc{B})&=(\nab\cdot\vc{B})\vc{A}-(\nab\cdot\vc{A})\vc{B}
     +(\vc{B}\cdot\nab)\vc{A}-(\vc{A}\cdot\nab)\vc{B}. \label{eq:prod-f}
\end{align}
\end{proposicion}
\begin{proof}[Demostración de \cref{eq:prod-a,eq:prod-c}]
La primera es la regla de Leibniz aplicada a cada componente: $\partial_i(fg)=f\,\partial_ig+g\,\partial_if$.
Para la segunda, desarrollamos la divergencia en coordenadas y usamos la regla
del producto en cada sumando:
\begin{align*}
  \nab\cdot(f\vc{A})&=\pd{}{x}(fA_x)+\pd{}{y}(fA_y)+\pd{}{z}(fA_z)\\
  &=f\Bigl(\pd{A_x}{x}+\pd{A_y}{y}+\pd{A_z}{z}\Bigr)
   +A_x\pd{f}{x}+A_y\pd{f}{y}+A_z\pd{f}{z}
   =f\,\nab\cdot\vc{A}+\vc{A}\cdot\nab f.
\end{align*}
Las identidades \eqref{eq:prod-d} y \eqref{eq:prod-e} se demuestran con el tensor de
Levi-Civita en el \cref{sec:levi-civita}
(\cref{ejer:div-prod-vect,ejer:rot-fA}); \eqref{eq:prod-b} y \eqref{eq:prod-f} se
obtienen de la misma manera con la identidad $\varepsilon$--$\delta$
(\cref{lema:eps-delta}).
\end{proof}

\begin{observacion}\label{obs:orden}
En las \cref{eq:prod-b,eq:prod-c,eq:prod-f} importa el \emph{orden} de los factores:
$\vc{A}\cdot\nab f$ es la derivada direccional de $f$ a lo largo de $\vc{A}$ (un
escalar), mientras que $\nab\cdot\vc{A}$ es la divergencia. Del mismo modo,
$\vc{B}\cdot\nab$ es un \emph{operador} que se aplica a lo que tiene a su
derecha.
\end{observacion}

% -------------------------------------------------------------------------
\subsection{Operadores de segundo orden}\label{sec:segundo-orden}
% -------------------------------------------------------------------------
Aplicando dos veces $\nab$ se obtienen cinco combinaciones. La más importante
en física es el laplaciano.

\begin{definicion}[Laplaciano]\label{def:laplaciano}
El \emph{laplaciano} de un campo escalar $f$ de clase $C^{2}$ es la divergencia
de su gradiente,
\begin{equation}\label{eq:laplaciano}
  \Delta f\defeq\nab^{2}f=\nab\cdot(\nab f)=\operatorname{div}(\operatorname{grad}f)
  =\sum_{i=1}^{n}\pdd{f}{x_i}\ ;\qquad
  \text{en }\R^{3}:\ \ \Lap f=\pdd{f}{x}+\pdd{f}{y}+\pdd{f}{z}.
\end{equation}
El \emph{laplaciano vectorial} de un campo $\vc{A}$ se define, en coordenadas
cartesianas, componente a componente:
\begin{equation}\label{eq:laplaciano-vect}
  \Lap\vc{A}\defeq(\Lap A_x)\,\ux+(\Lap A_y)\,\uy+(\Lap A_z)\,\uz .
\end{equation}
\end{definicion}

\begin{observacion}
La definición \eqref{eq:laplaciano-vect} es válida \emph{solo} en cartesianas.
En otros sistemas de coordenadas el laplaciano vectorial no es el laplaciano
escalar de cada componente, porque los versores $\ur,\uth,\uph,\ldots$
cambian de punto a punto (\cref{sec:coordenadas}); su definición general es
$\Lap\vc{A}=\nab(\nab\cdot\vc{A})-\nab\times(\nab\times\vc{A})$.
\end{observacion}

Las otras cuatro combinaciones se reducen a dos identidades fundamentales y a
una relación entre ellas.

\begin{proposicion}[Relaciones fundamentales]\label{prop:fundamentales}
Para un campo escalar $f$ y un campo vectorial $\vc{A}$, ambos de clase $C^{2}$,
\begin{equation}\label{eq:rotgrad}
  \nab\times(\nab f)=\vc{0},
  \qquad\qquad
  \nab\cdot(\nab\times\vc{A})=0 .
\end{equation}
\end{proposicion}
\begin{proof}
Ambas se reducen a la simetría de las derivadas parciales segundas
cruzadas (teorema de Schwarz), válida para funciones de clase $C^{2}$.
Desarrollando el rotacional del gradiente,
\[
  \nab\times\nab f=
  \begin{vmatrix}
    \ux & \uy & \uz\\
    \partial_x & \partial_y & \partial_z\\
    \partial_xf & \partial_yf & \partial_zf
  \end{vmatrix}
  =\Bigl(\pdm{f}{y}{z}-\pdm{f}{z}{y}\Bigr)\ux
   +\Bigl(\pdm{f}{z}{x}-\pdm{f}{x}{z}\Bigr)\uy
   +\Bigl(\pdm{f}{x}{y}-\pdm{f}{y}{x}\Bigr)\uz=\vc{0}.
\]
Del mismo modo, para $\vc{A}=(A_x,A_y,A_z)$,
\begin{align*}
  \nab\cdot(\nab\times\vc{A})
  &=\pd{}{x}\Bigl(\pd{A_z}{y}-\pd{A_y}{z}\Bigr)
   +\pd{}{y}\Bigl(\pd{A_x}{z}-\pd{A_z}{x}\Bigr)
   +\pd{}{z}\Bigl(\pd{A_y}{x}-\pd{A_x}{y}\Bigr)\\
  &=\Bigl(\pdm{A_z}{x}{y}-\pdm{A_z}{y}{x}\Bigr)
   +\Bigl(\pdm{A_x}{y}{z}-\pdm{A_x}{z}{y}\Bigr)
   +\Bigl(\pdm{A_y}{z}{x}-\pdm{A_y}{x}{z}\Bigr)=0 .
\end{align*}
En el \cref{sec:levi-civita} se da otra demostración, con índices, que
explica por qué ocurre: un objeto simétrico contraído con otro antisimétrico
es siempre cero.
\end{proof}

\begin{observacion}[Reglas mnemotécnicas, con cautela]\label{obs:mnemo}
El producto vectorial cumple $\vc{v}\times\vc{v}=\vc{0}$, lo que sugiere
$\nab\times\nab=\vc{0}$ y, aplicado a $f$, la primera relación. Cumple también
$\vc{v}\cdot(\vc{v}\times\vc{w})=0$ (el producto vectorial es ortogonal a sus
factores), lo que sugiere la segunda. Pero, como $\nab$ no es un vector, esta
manipulación \emph{sugiere} sin \emph{demostrar}, y puede llevar a error:
aunque $(\vc{v}\times\vc{w})\cdot\vc{w}=0$ para vectores, en general
$(\nab\times\vc{F})\cdot\vc{F}\neq0$.
\end{observacion}

\begin{corolario}[Los gradientes son irrotacionales]\label{cor:grad-irrot}
Si $\vc{F}=\nab f$ es un campo gradiente de clase $C^{1}$ ($f\in C^{2}$), entonces
$\nab\times\vc{F}=\vc{0}$. En particular, ser irrotacional es una condición
\emph{necesaria} para ser conservativo. Es también suficiente en regiones
estrelladas (lema de Poincaré)~\cite{NAVAS}.
\end{corolario}

La combinación restante que se usa en electromagnetismo es la del doble
rotacional.

\begin{proposicion}[Rotacional del rotacional]\label{prop:rotrot}
Para un campo vectorial $\vc{A}$ de clase $C^{2}$,
\begin{equation}\label{eq:rotrot}
  \nab\times(\nab\times\vc{A})=\nab(\nab\cdot\vc{A})-\Lap\vc{A}.
\end{equation}
\end{proposicion}
\begin{proof}
Véase el \cref{ejer:rotrot}, donde se obtiene con el tensor de Levi-Civita y la
identidad $\varepsilon$--$\delta$.
\end{proof}

El resto de combinaciones, $\nab(\nab\cdot\vc{A})$, no tiene un nombre ni un uso
especial en este curso; solo aparece en la \cref{eq:rotrot}.

% =========================================================================
\section{Cálculo integral vectorial: los teoremas integrales fundamentales}\label{sec:integral}
% =========================================================================

% -------------------------------------------------------------------------
\subsection{Del teorema fundamental del cálculo a los teoremas vectoriales}\label{sec:tfc}
% -------------------------------------------------------------------------
Recordemos el teorema fundamental del cálculo en una variable: para una
función $f$ derivable en $[a,b]$,
\begin{equation}\label{eq:tfc}
  \int_{a}^{b}\frac{\dif f}{\dif x}\,\dif x=f(b)-f(a).
\end{equation}
La idea es la de una \emph{suma de variaciones}: si dividimos $[a,b]$ en
intervalitos de longitud $\dif x$, en cada uno de ellos la función cambia en
$\dif f$; al sumar todos los cambios infinitesimales, las contribuciones
intermedias se cancelan por parejas y solo queda la diferencia entre lo que
vale la función en los extremos (\cref{fig:tfc}). Es la primera instancia de un
principio que se repite en todo el cálculo vectorial:

\begin{quote}\itshape
La integral de una «derivada» sobre una región es igual a la integral de la
función sobre el borde de esa región.
\end{quote}

\begin{figure}[ht]
\centering
\begin{tikzpicture}[>=Stealth, x=1cm, y=1cm]
  \draw[eje] (-0.3,0) -- (8.6,0) node[right, text=black] {$x$};
  \draw[eje] (0,-0.3) -- (0,4.4) node[above, text=black] {$f(x)$};
  \def\fx#1{0.9+0.55*(#1)+0.9*sin(deg(0.55*(#1)))}
  % barras de variación
  \foreach \xx in {1.0,1.6,...,7.0}{%
    \pgfmathsetmacro{\ya}{\fx{\xx}}\pgfmathsetmacro{\yb}{\fx{\xx+0.6}}%
    \draw[azulclaro, line width=0.7pt] (\xx,\ya) -- (\xx+0.6,\ya);
    \draw[rojoteo!70, line width=1.0pt] (\xx+0.6,\ya) -- (\xx+0.6,\yb);}
  \draw[azulprof, line width=1.3pt, domain=1.0:7.6, samples=120, smooth]
    plot (\x, {\fx{\x}});
  \foreach \xx/\lab in {1.0/a,7.6/b}{%
    \pgfmathsetmacro{\yy}{\fx{\xx}}%
    \draw[densely dashed, grispizarra] (\xx,0) -- (\xx,\yy);
    \fill[azulprof] (\xx,\yy) circle (2pt);
    \node[below] at (\xx,0) {$\lab$};}
  \node[font=\small, text=rojoteo, anchor=west] at (3.0,3.7) {$\dif f$};
  \draw[->, rojoteo, line width=0.5pt] (3.0,3.65) -- (2.55,3.0);
  \node[font=\small, text=azulacero, anchor=north] at (3.9,0.35) {$\dif x$};
  \draw[|<->|, azulacero, line width=0.5pt] (3.4,0.3) -- (4.0,0.3);
  \node[font=\small, text=azulprof, rotate=90, anchor=south] at (9.1,{(\fx{1.0}+\fx{7.6})/2}) {$f(b)-f(a)=\textstyle\sum\dif f$};
  \draw[|<->|, grispizarra, line width=0.5pt] (8.5,{\fx{1.0}}) -- (8.5,{\fx{7.6}});
  \draw[densely dashed, grispizarra] (7.6,{\fx{7.6}}) -- (8.5,{\fx{7.6}});
  \draw[densely dashed, grispizarra] (1.0,{\fx{1.0}}) -- (8.5,{\fx{1.0}});
\end{tikzpicture}
\caption{Teorema fundamental del cálculo como suma de variaciones: la suma de
los incrementos $\dif f$ (en rojo) sobre una partición de $[a,b]$ es el
incremento total $f(b)-f(a)$.}
\label{fig:tfc}
\end{figure}

El «borde» de un intervalo $[a,b]$ son sus dos extremos, a los que se asigna una
orientación: $b$ cuenta con signo $+$ (punto de salida) y $a$ con signo $-$
(punto de entrada). Al aumentar la dimensión del dominio de integración se
obtienen tres generalizaciones, que son los teoremas de este apartado:
\begin{center}
\renewcommand{\arraystretch}{1.3}
\begin{tabular}{@{}llll@{}}
\toprule
Dominio $\Omega$ & Dimensión & «Derivada» & Borde $\partial\Omega$\\
\midrule
curva $C$ & 1 & gradiente $\nab f$ & dos puntos $\{\vc{a}^{-},\vc{b}^{+}\}$\\
superficie $S$ & 2 & rotacional $\nab\times\vc{A}$ & curva cerrada $L=\partial S$\\
sólido $V$ & 3 & divergencia $\nab\cdot\vc{A}$ & superficie cerrada $S=\partial V$\\
\bottomrule
\end{tabular}
\end{center}

% -------------------------------------------------------------------------
\subsection{Integral de línea}\label{sec:linea}
% -------------------------------------------------------------------------
\begin{definicion}[Integral de línea]\label{def:linea}
Sean
\begin{itemize}
  \item $C\subseteq\R^{n}$ una curva regular parametrizada por una aplicación
        $\gamma:[a,b]\to\R^{n}$ de clase $C^{1}$ con $\gamma'(t)\neq\vc{0}$, y
  \item $\vc{F}:C\to\R^{n}$ un campo vectorial continuo sobre la curva.
\end{itemize}
Se define la \emph{integral de línea} (o \emph{de camino}, o \emph{de trabajo})
de $\vc{F}$ a lo largo de $C$ como la integral ordinaria
\begin{equation}\label{eq:linea}
  \int_{C}\vc{F}\cdot\dl\defeq\int_{a}^{b}\bigl\langle\vc{F}(\gamma(t)),\gamma'(t)\bigr\rangle\,\dif t
  =\int_{a}^{b}\bigl[(\vc{F}\circ\gamma)\cdot\gamma'\bigr](t)\,\dif t .
\end{equation}
El símbolo $\dl\defeq\gamma'(t)\,\dif t$ es el \emph{elemento de longitud
dirigido}: un desplazamiento infinitesimal tangente a la curva,
$\dl=\sum_i\dif x_i\,\uh{i}$ con $\dif x_i=x_i'(t)\,\dif t$. Si $C$ es cerrada se
escribe $\oint_C$ y se habla de \emph{circulación}.
\end{definicion}

Físicamente, $\int_C\vc{F}\cdot\dl$ es el \emph{trabajo} que realiza el campo de
fuerzas $\vc{F}$ sobre una partícula que recorre $C$: en cada tramo infinitesimal
solo cuenta la componente de $\vc{F}$ a lo largo del desplazamiento.

\begin{notacion}[Procedimiento para calcular una integral de línea]\label{not:procedimiento-linea}
Se procede en cinco pasos.\begin{enumerate}[label=\textbf{\arabic*.}, itemsep=1pt]
  \item \emph{Parametrizar} la curva: $\vc{r}=\gamma(t)=(x_1(t),\ldots,x_n(t))$, $t\in[a,b]$.
  \item \emph{Derivar}: calcular el vector tangente $\gamma'(t)=(x_1'(t),\ldots,x_n'(t))$.
  \item \emph{Restringir} el campo a la trayectoria: $(\vc{F}\circ\gamma)(t)=\vc{F}(\gamma(t))$.
  \item \emph{Proyectar}: producto escalar $\vc{F}(\gamma(t))\cdot\gamma'(t)$, que es una función de una sola variable.
  \item \emph{Integrar} esa función de $t$ sobre $[a,b]$.
\end{enumerate}
\end{notacion}

\paragraph{En coordenadas.} Si $\vc{F}=(f_1,\ldots,f_n)$, la integral de línea es
la suma de $n$ integrales ordinarias:
\begin{equation}\label{eq:linea-coord}
  \int_{C}\vc{F}\cdot\dl=\int_{a}^{b}\sum_{k=1}^{n}(f_k\circ\gamma)(t)\,x_k'(t)\,\dif t
  =\sum_{k=1}^{n}\int_{a}^{b}f_k\bigl(\gamma(t)\bigr)\,x_k'(t)\,\dif t
\end{equation}
(aditividad de la integral).

\begin{definicion}[Notación diferencial y 1-formas]\label{def:diferencial-linea}
Con el cambio de variable $\dif x_k=x_k'(t)\,\dif t$ se tiene
$f_k(\vc{r})\,\dif x_k=f_k(x_1(t),\ldots,x_n(t))\,x_k'(t)\,\dif t$ y la integral de
línea queda
\[
  \int_{C}\vc{F}\cdot\dl=\int_{C}f_1\,\dif x_1+f_2\,\dif x_2+\cdots+f_n\,\dif x_n ,
\]
donde cada sumando se integra por separado. A una expresión de este tipo,
$\omega=f_1\,\dif x_1+\cdots+f_n\,\dif x_n=\sum_i f_i\,\dif x_i$ con $f_k$
campos escalares, se le llama \emph{1-forma diferencial}.
\end{definicion}

En dos y tres dimensiones, con $\vc{F}=(P,Q)$ y $\vc{F}=(P,Q,R)$ respectivamente,
\begin{equation}\label{eq:linea-2d3d}
  \int_{C}\vc{F}\cdot\dl=\int_{C}P\,\dif x+Q\,\dif y
  \qquad\text{y}\qquad
  \int_{C}\vc{F}\cdot\dl=\int_{C}P\,\dif x+Q\,\dif y+R\,\dif z,
\end{equation}
donde, por ejemplo, $\int_{C}P\,\dif x=\int_a^bP\bigl(x(t),y(t),z(t)\bigr)\,x'(t)\,\dif t$.

\begin{ejemplo}[Segmento rectilíneo]\label{ej:segmento}
Sea $C=[\vc{p},\vc{q}]$ el segmento que une dos puntos de $\R^{n}$. Su
parametrización estándar es $\sigma(t)=\vc{p}+t\,\vc{h}$ con
$\vc{h}=\vc{q}-\vc{p}$ y $t\in[0,1]$: cumple $\sigma(0)=\vc{p}$,
$\sigma(1)=\vc{q}$ y $\sigma'(t)=\vc{h}$ constante. Para cualquier campo
continuo,
\[
  \int_{[\vc{p},\vc{q}]}\vc{F}\cdot\dl
  =\int_{0}^{1}\vc{F}(\vc{p}+t\,\vc{h})\cdot\vc{h}\,\dif t
  =\int_{0}^{1}\vc{F}\bigl(\vc{p}+t(\vc{q}-\vc{p})\bigr)\cdot(\vc{q}-\vc{p})\,\dif t .
\]
\end{ejemplo}

\begin{proposicion}[Propiedades de la integral de línea]\label{prop:prop-linea}
\begin{enumerate}[label=(\alph*)]
  \item \emph{Linealidad:} $\displaystyle\int_C(\alpha\vc{F}+\beta\vc{G})\cdot\dl=\alpha\int_C\vc{F}\cdot\dl+\beta\int_C\vc{G}\cdot\dl$.
  \item \emph{Aditividad:} si $C=C_1\uplus C_2$ (la segunda curva empieza donde acaba la primera), $\int_C=\int_{C_1}+\int_{C_2}$.
  \item \emph{Cambio de orientación:} si $C^{\mathrm{op}}$ es la misma curva recorrida en sentido contrario, $\int_{C^{\mathrm{op}}}\vc{F}\cdot\dl=-\int_C\vc{F}\cdot\dl$.
  \item \emph{Invariancia por reparametrización:} si $\alpha:[a,b]\to C$ y
        $\beta:[A,B]\to C$ son dos parametrizaciones regulares de la misma curva y
        recorren $C$ en el mismo sentido, las integrales \eqref{eq:linea}
        coinciden.
\end{enumerate}
\end{proposicion}
\begin{proof}
Las tres primeras se siguen de las propiedades de la integral de Riemann.
Para (d), sea $\varphi=\beta^{-1}\circ\alpha:[a,b]\to[A,B]$, un cambio de
variable de clase $C^{1}$ y biyectivo, de modo que $\alpha=\beta\circ\varphi$ y
$\alpha'(t)=\beta'(\varphi(t))\,\varphi'(t)$. Entonces
\[
  \int_a^b\vc{F}(\alpha(t))\cdot\alpha'(t)\,\dif t
  =\int_a^b\vc{F}\bigl(\beta(\varphi(t))\bigr)\cdot\beta'(\varphi(t))\,\varphi'(t)\,\dif t
  =\int_A^B\vc{F}(\beta(s))\cdot\beta'(s)\,\dif s ,
\]
con $s=\varphi(t)$, siempre que $\varphi'>0$ (mismo sentido). Si $\varphi'<0$, los
límites de integración se invierten y la integral cambia de signo, que es (c).
\end{proof}

La integral de línea está bien definida: depende de la curva y de su
orientación, pero no de cómo se parametrice.

% -------------------------------------------------------------------------
\subsection{El teorema del gradiente}\label{sec:teorema-gradiente}
% -------------------------------------------------------------------------
\begin{teorema}[Teorema del gradiente o fundamental de las integrales de línea]\label{teo:gradiente}
Sea $f$ un campo escalar de clase $C^{1}$ en una región $W\subseteq\R^{n}$ y sea
$C\subset W$ un arco regular orientado cuyo borde es $\partial C=\{\vc{a}^{-},\vc{b}^{+}\}$
(punto de partida $\vc{a}$, punto de llegada $\vc{b}$). Entonces
\begin{equation}\label{eq:teo-gradiente}
  \boxed{\;\int_{C}\nab f\cdot\dl=f(\vc{b})-f(\vc{a})\;}
\end{equation}
\end{teorema}
\begin{proof}
Sea $\gamma:[a,b]\to C$ una parametrización compatible con la orientación, con
$\gamma(a)=\vc{a}$ y $\gamma(b)=\vc{b}$. Por la \cref{prop:cadena},
$\nab f(\gamma(t))\cdot\gamma'(t)=(f\circ\gamma)'(t)$, y por el teorema
fundamental del cálculo en una variable, \eqref{eq:tfc},
\[
  \int_C\nab f\cdot\dl=\int_a^b\nab f(\gamma(t))\cdot\gamma'(t)\,\dif t
  =\int_a^b(f\circ\gamma)'(t)\,\dif t=f(\gamma(b))-f(\gamma(a))=f(\vc{b})-f(\vc{a}).\qedhere
\]
\end{proof}

El valor de la integral queda determinado única y exclusivamente por la
diferencia de potencial en los extremos del camino. Con la notación de la
\cref{sec:tfc}: la integral de la «derivada» $\nab f$ sobre $C$ es igual a $f$
evaluada en el borde $\partial C$, contando $+f(\vc{b})$ y $-f(\vc{a})$. Dos
consecuencias inmediatas son de gran importancia en física.

\begin{corolario}[Independencia del camino y circulación nula]\label{cor:gradiente}
Sea $\vc{F}=\nab f$ un campo gradiente en $W$.
\begin{enumerate}[label=(\alph*)]
  \item El trabajo de $\vc{F}$ entre dos puntos es el mismo sobre cualquier curva
        que los una: depende solo de los extremos, no del camino.
  \item Para toda curva cerrada $C\subset W$ (con $\vc{r}(a)=\vc{r}(b)$),
        \begin{equation}\label{eq:circulacion-nula}
          \oint_C\nab f\cdot\dl=0 .
        \end{equation}
\end{enumerate}
\end{corolario}
\begin{proof}
(a) Es inmediato de la \cref{eq:teo-gradiente}. (b) Si la curva es cerrada,
$\vc{b}=\vc{a}$ y $f(\vc{b})-f(\vc{a})=0$.
\end{proof}

\begin{definicion}[Independencia del camino y circulación nula]\label{def:indep-camino}
Un campo $\vc{F}$ en $W$ tiene la propiedad de \emph{independencia del camino}
si $\int_\alpha\vc{F}\cdot\dl=\int_\beta\vc{F}\cdot\dl$ para cualquier par de
curvas $\alpha,\beta\subset W$ con los mismos extremos orientados; y tiene
\emph{circulación nula} si $\oint_C\vc{F}\cdot\dl=0$ para toda curva cerrada
$C\subset W$.
\end{definicion}

\begin{proposicion}\label{prop:equiv-indep}
Un campo continuo tiene independencia del camino en $W$ si y solo si tiene
circulación nula en $W$.
\end{proposicion}
\begin{proof}
($\Leftarrow$) Si $A$ y $B$ unen los mismos extremos, $C=A\uplus B^{\mathrm{op}}$ es
un circuito, y por la \cref{prop:prop-linea},
$0=\oint_C\vc{F}\cdot\dl=\int_A\vc{F}\cdot\dl-\int_B\vc{F}\cdot\dl$.
($\Rightarrow$) Si $C$ es cerrada, elegimos dos puntos distintos de $C$, que la
dividen en dos arcos $A$ y $B$ con los mismos extremos y $C=A\uplus B^{\mathrm{op}}$;
entonces $\oint_C\vc{F}\cdot\dl=\int_A\vc{F}\cdot\dl-\int_B\vc{F}\cdot\dl=0$.
\end{proof}

\begin{proposicion}[Recíproco del teorema del gradiente]\label{prop:reciproco}
Un campo continuo $\vc{F}$ en una región $W$ con independencia del camino es un
campo gradiente: fijado $\vc{p}\in W$, la función
$f(\vc{r})=\int_{\vc{p}}^{\vc{r}}\vc{F}\cdot\dl$ (calculada sobre cualquier
curva de $W$ que una $\vc{p}$ con $\vc{r}$) cumple $\nab f=\vc{F}$.
\end{proposicion}
\begin{proof}
Por independencia del camino, $f$ está bien definida. Para $h$ pequeño,
$f(\vc{r}+h\,\uh{i})-f(\vc{r})$ es la integral de $\vc{F}$ sobre el segmento que
une $\vc{r}$ con $\vc{r}+h\,\uh{i}$, es decir,
$\int_0^hF_i(\vc{r}+t\,\uh{i})\,\dif t$. Dividiendo por $h$ y haciendo
$h\to0$, la continuidad de $F_i$ da $\partial_if(\vc{r})=F_i(\vc{r})$.
\end{proof}

Reuniendo todo: para campos continuos, \emph{ser gradiente $\Leftrightarrow$
independencia del camino $\Leftrightarrow$ circulación nula}; y ser irrotacional es
condición necesaria (\cref{cor:grad-irrot}), y suficiente en regiones
estrelladas.

% Figura: teorema del gradiente. (a) Trabajo de un gradiente entre a y b por
% dos caminos; (b) circulación nula en un circuito cerrado.
\newcommand{\curvasNivelElipt}{%
  \foreach \c/\pc in {0.6/100,1.2/80,1.8/60,2.4/40}{%
    \draw[azulprof!\pc!azulclaro, line width=0.8pt] (0,0) ellipse ({2*\c} and {\c});}
  \fill[azulprof] (0,0) circle (1.4pt);}
\begin{figure}[ht]
\centering
\begin{minipage}[t]{0.49\textwidth}\centering
\begin{tikzpicture}[>=Stealth, x=0.74cm, y=0.74cm, baseline=(current bounding box.north)]
  \clip (-5.1,-3.2) rectangle (5.1,3.2);
  \curvasNivelElipt
  % camino principal: x(t)=-4.2+6t, y(t)=-2.2+3.0t-1.4 sin(pi t)
  \draw[trayecto, domain=0:1, samples=60, smooth, variable=\t]
    plot ({-4.0+5.8*\t},{-2.0+3.0*\t-1.2*sin(180*\t)});
  % camino alternativo: segmento
  \draw[acento!60, densely dashed, line width=0.8pt] (-4.0,-2.0) -- (1.8,1.0);
  % vectores en puntos del camino principal
  \foreach \t in {0.18,0.52,0.86}{%
    \pgfmathsetmacro{\px}{-4.0+5.8*\t}\pgfmathsetmacro{\py}{-2.0+3.0*\t-1.2*sin(180*\t)}%
    \pgfmathsetmacro{\dx}{5.8}\pgfmathsetmacro{\dy}{3.0-1.2*pi*cos(180*\t)}%
    \pgfmathsetmacro{\dn}{sqrt(\dx*\dx+\dy*\dy)}%
    \pgfmathsetmacro{\gx}{-\px/2}\pgfmathsetmacro{\gy}{-2*\py}%
    \pgfmathsetmacro{\gn}{sqrt(\gx*\gx+\gy*\gy)}%
    \draw[->, azulprof, line width=1pt] (\px,\py) -- ++({0.8*\dx/\dn},{0.8*\dy/\dn});
    \draw[->, line width=0.9pt, cm2] (\px,\py) -- ++({1.2*\gx/\gn},{1.2*\gy/\gn});
    \fill (\px,\py) circle (0.8pt);}
  \fill[acento] (-4.0,-2.0) circle (1.2pt) node[below left=0pt, text=black] {$\vc{a}$};
  \fill[acento] (1.8,1.0) circle (1.2pt) node[above right=0pt, text=black] {$\vc{b}$};
  \node[font=\small, text=acento] at (-0.9,-2.1) {$C$};
  
  \node[fill=white, fill opacity=0.85, text opacity=1, inner sep=2pt, font=\footnotesize, anchor=south west]
       at (-5.0,2.45) {\textcolor{azulprof}{$\to\ \dl$}\quad\textcolor{cm2}{$\to\ \nab f$}};
\end{tikzpicture}\\[3pt]{\small\scshape (a)}\; $\displaystyle\int_C\nab f\cdot\dl=f(\vc{b})-f(\vc{a})$
\end{minipage}\hfill
\begin{minipage}[t]{0.49\textwidth}\centering
\begin{tikzpicture}[>=Stealth, x=0.74cm, y=0.74cm, baseline=(current bounding box.north)]
  \clip (-5.1,-3.2) rectangle (5.1,3.2);
  \curvasNivelElipt
  % circuito cerrado: elipse desplazada, sentido antihorario
  \def\cx{0.5}\def\cy{0.1}
  \draw[trayecto] ({\cx},{\cy}) ellipse (3.5 and 2.2);
  \foreach \s in {20,140,260}{%
    \pgfmathsetmacro{\px}{\cx+3.5*cos(\s)}\pgfmathsetmacro{\py}{\cy+2.2*sin(\s)}%
    \pgfmathsetmacro{\dx}{-3.5*sin(\s)}\pgfmathsetmacro{\dy}{2.2*cos(\s)}%
    \pgfmathsetmacro{\dn}{sqrt(\dx*\dx+\dy*\dy)}%
    \pgfmathsetmacro{\gx}{-\px/2}\pgfmathsetmacro{\gy}{-2*\py}%
    \pgfmathsetmacro{\gn}{sqrt(\gx*\gx+\gy*\gy)}%
    \draw[->, azulprof, line width=1pt] (\px,\py) -- ++({0.8*\dx/\dn},{0.8*\dy/\dn});
    \draw[->, line width=0.9pt, cm2] (\px,\py) -- ++({1.2*\gx/\gn},{1.2*\gy/\gn});
    \fill (\px,\py) circle (0.8pt);}
  \node[font=\small, text=acento] at (4.4,2.55) {$C$};
  \node[fill=white, fill opacity=0.85, text opacity=1, inner sep=2pt, font=\footnotesize, anchor=south west]
       at (-5.0,2.45) {\textcolor{azulprof}{$\to\ \dl$}\quad\textcolor{cm2}{$\to\ \nab f$}};
\end{tikzpicture}\\[3pt]{\small\scshape (b)}\; $\displaystyle\oint_C\nab f\cdot\dl=0$
\end{minipage}
\caption{Teorema del gradiente sobre el mapa de curvas de nivel de una
colina $f$ (ahora la cima es el máximo). (a)~Sobre la curva $C$ (y sobre
cualquier otra con los mismos extremos, p.\,ej. el segmento discontinuo) la
integral suma, en cada tramo, la proyección del gradiente sobre el desplazamiento
$\dl$: tramos cuesta arriba suman, tramos cuesta abajo restan, y el balance es
solo $f(\vc{b})-f(\vc{a})$. (b)~En un circuito cerrado, lo que se sube se baja:
la circulación es cero.}
\label{fig:teor_grad}
\end{figure}


% -------------------------------------------------------------------------
\subsection{Integral de flujo}\label{sec:flujo}
% -------------------------------------------------------------------------
Pasamos a la integral que formaliza la noción de \emph{flujo} de un campo a
través de una superficie en $\R^{3}$. Ahora el vector relevante es la normal.

\begin{definicion}[Superficie parametrizada y vector normal]\label{def:superficie}
Una \emph{superficie parametrizada} $S\subset\R^{3}$ es la imagen de una
aplicación $\sigma:S^{*}\subseteq\R^{2}\to\R^{3}$, $\sigma(u,v)=(x(u,v),y(u,v),z(u,v))$,
de clase $C^{1}$ y regular ($\sigma_u\times\sigma_v\neq\vc{0}$). Su \emph{vector
normal asociado}, su normal unitaria y sus elementos de área son
\[
  \vc{n}_\sigma=\frac{\partial\sigma}{\partial u}\times\frac{\partial\sigma}{\partial v},
  \qquad
  \un=\frac{\vc{n}_\sigma}{\norm{\vc{n}_\sigma}},
  \qquad
  \dsc=\norm{\vc{n}_\sigma}\,\dif u\,\dif v,
  \qquad
  \dS=\un\,\dsc=\vc{n}_\sigma\,\dif u\,\dif v .
\]
La elección de $\un$ (de entre sus dos posibles sentidos) es una
\emph{orientación} de $S$.
\end{definicion}

\begin{definicion}[Integral de flujo]\label{def:flujo}
Dado un campo vectorial $\vc{F}$ continuo sobre una superficie orientada
$S=\sigma(S^{*})$, el \emph{flujo} de $\vc{F}$ a través de $S$ es
\begin{equation}\label{eq:flujo}
  \iint_{S}\vc{F}\cdot\dS\defeq\iint_{S^{*}}\vc{F}\bigl(\sigma(u,v)\bigr)\cdot\vc{n}_\sigma(u,v)\,\dif u\,\dif v
  =\iint_{S}\vc{F}\cdot\un\,\dsc ,
\end{equation}
donde la segunda igualdad (con $\un\,\dsc=\dS$) dice que el flujo es la integral
de \emph{área} de la componente \emph{normal} de $\vc{F}$. Para una superficie
cerrada se escribe $\oiint_S\vc{F}\cdot\dS$ y, por convenio, se toma la normal
exterior.
\end{definicion}

Solo la componente normal contribuye: un campo que corre \emph{a lo largo} de la
superficie, sin atravesarla, no aporta flujo (\cref{fig:flujo}). Ahora bien,
\emph{cualquier superficie es localmente plana}: en un entorno infinitesimal de
cada punto, el área orientada $\dS$ es un pequeño parche plano de normal $\un$,
y el flujo infinitesimal es $\vc{F}\cdot\dS=\norm{\vc{F}}\,\dsc\cos\theta$.

\begin{notacion}[Procedimiento para calcular un flujo]\label{not:procedimiento-flujo}
Se procede en seis pasos.\begin{enumerate}[label=\textbf{\arabic*.}, itemsep=1pt]
  \item \emph{Parametrizar} la superficie: $\sigma(u,v)=(x(u,v),y(u,v),z(u,v))$.
  \item \emph{Derivar}: $\sigma_u=\partial\sigma/\partial u$ y $\sigma_v=\partial\sigma/\partial v$.
  \item \emph{Normal}: $\vc{n}_\sigma=\sigma_u\times\sigma_v$ (comprobar la orientación).
  \item \emph{Restringir} el campo: $\vc{F}\circ\sigma$.
  \item \emph{Proyectar}: $(\vc{F}\circ\sigma)\cdot\vc{n}_\sigma$, en función de $u,v$.
  \item \emph{Integrar} sobre la región plana $(u,v)\in S^{*}$.
\end{enumerate}
\end{notacion}

\begin{observacion}[Notación diferencial: $2$-formas]\label{obs:2formas}
En coordenadas, $\vc{n}_\sigma=\bigl(\tfrac{\partial(y,z)}{\partial(u,v)},
\tfrac{\partial(z,x)}{\partial(u,v)},\tfrac{\partial(x,y)}{\partial(u,v)}\bigr)$,
donde $\tfrac{\partial(x,y)}{\partial(u,v)}=x_uy_v-x_vy_u$ son subdeterminantes
jacobianos. Escribiendo $\dif y\wedge\dif z\defeq\tfrac{\partial(y,z)}{\partial(u,v)}\,\dif u\,\dif v$
(y cíclicamente, $x\to y\to z\to x$), con $\vc{F}=(P,Q,R)$ queda
$\iint_S\vc{F}\cdot\dS=\iint_SP\,\dif y\wedge\dif z+Q\,\dif z\wedge\dif x+R\,\dif x\wedge\dif y$,
una \emph{2-forma diferencial}; de momento $\wedge$ es solo una notación para
estos determinantes~\cite{NAVAS}.
\end{observacion}

\begin{ejemplo}[Flujo del campo inverso del cuadrado]\label{ej:flujo-r3}
Sea $\vc{F}=\rho\,\vc{r}/r^{3}$ ($\rho$ constante, $r=\norm{\vc{r}}$) y
$S_R:x^{2}+y^{2}+z^{2}=R^{2}$. La normal unitaria a una esfera es el radio
normalizado, $\un=\vc{r}/r$, de modo que $\vc{F}\cdot\un=\rho\,r^{2}/r^{4}=\rho/R^{2}$
es constante sobre $S_R$ y
\[
  \oiint_{S_R}\vc{F}\cdot\un\,\dsc=\frac{\rho}{R^{2}}\cdot4\pi R^{2}=4\pi\rho ,
\]
independiente de $R$. Para un campo eléctrico ($\rho\propto q$) es la ley de
Gauss; su explicación profunda (la divergencia de este campo es una delta de
Dirac en el origen) se verá en la \cref{sec:delta}.
\end{ejemplo}

% Figura: flujo de un campo a través de una superficie. (a) superficie
% ondulada inmersa en un campo; (b) elemento de área dirigido y componente normal.
\begin{figure}[ht]
\centering
\begin{minipage}[t]{0.56\textwidth}\centering
\begin{tikzpicture}[baseline=(current bounding box.north)]
  \begin{axis}[hide axis, width=1.12\textwidth, height=0.95\textwidth, view={-28}{40},
      xmin=-3.2, xmax=3.2, ymin=-2.4, ymax=2.4, zmin=-1.6, zmax=2.2,
      colormap name=verdeamarillo, point meta min=-0.9, point meta max=0.9,
      /pgf/number format/use comma, clip=false]
    % campo por debajo de la superficie
    \addplot3[quiver={u=0.34, v=0.18, w=0.80}, -{Stealth[length=3pt]}, line width=0.7pt,
              acento, domain=-3.0:3.0, domain y=-2.2:2.2, samples=6, samples y=5]
      {-1.2};
    % superficie
    \addplot3[surf, shader=interp, samples=28, samples y=22, domain=-2.5:2.5, domain y=-1.8:1.8,
              opacity=0.93, mesh/ordering=x varies]
      {0.85*sin(deg(1.25*x))*cos(deg(0.5*y))+0.12*y};
    % campo por encima de la superficie
    \addplot3[quiver={u=0.34, v=0.18, w=0.80}, -{Stealth[length=3pt]}, line width=0.7pt,
              acento, domain=-3.0:3.0, domain y=-2.2:2.2, samples=6, samples y=5]
      {1.0};
    \node[font=\small, text=acento] at (axis cs:2.6,0.3,1.9) {$\vc{A}$};
    \node[font=\small, text=black!80, fill=white, fill opacity=0.7, text opacity=1, inner sep=1pt] at (axis cs:-2.0,-1.2,-0.5) {$S$};
  \end{axis}
\end{tikzpicture}\\[2pt]{\small\scshape (a)}
\end{minipage}\hfill
\begin{minipage}[t]{0.40\textwidth}\centering
\begin{tikzpicture}[>=Stealth, x=1cm, y=1cm, baseline=(current bounding box.north)]
  \shade[top color=azulclaro!55, bottom color=azulfondo, shading angle=0]
    (0,0) -- (3.4,0) -- (2.7,1.1) -- (-0.7,1.1) -- cycle;
  \draw[grispizarra!70, line width=0.4pt] (0,0) -- (3.4,0) -- (2.7,1.1) -- (-0.7,1.1) -- cycle;
  \coordinate (O) at (1.35,0.55);
  \draw[->, line width=1.6pt, azulacero] (O) -- ++(0,2.1) node[left, text=azulacero] {$\dS$};
  \draw[->, line width=1.6pt, rojoteo] (O) -- ++(2.0,1.55) node[right, text=rojoteo] {$\vc{A}$};
  % componente normal
  \draw[densely dashed, grispizarra] (3.35,2.1) -- (1.35,2.1);
  \draw[line width=3pt, azulclaro, opacity=0.9] (O) -- ++(0,1.55);
  \node[font=\footnotesize, text=azulprof, anchor=east] at (1.25,1.5) {$\vc{A}\cdot\un$};
  \draw[grispizarra, line width=0.5pt] (1.35,1.15) arc (90:37.8:0.6);
  \node[font=\small] at (1.85,1.25) {$\theta$};
  \fill (O) circle (1.4pt);
\end{tikzpicture}\\[2pt]{\small\scshape (b)}
\end{minipage}
\caption{Flujo de un campo vectorial $\vc{A}$ a través de una superficie $S$.
(a)~Un campo (flechas rojas) atraviesa una superficie curva; el flujo
contabiliza solo la parte del campo que corta a la superficie. (b)~En un
entorno de cada punto la superficie es plana: $\dS=\un\,\dsc$ es un vector de
módulo igual al área del parche y perpendicular a él, y el flujo
elemental es $\vc{A}\cdot\dS=(\vc{A}\cdot\un)\,\dsc=\norm{\vc{A}}\cos\theta\,\dsc$.}
\label{fig:flujo}
\end{figure}


% -------------------------------------------------------------------------
\subsection{El teorema de Stokes (del rotacional)}\label{sec:stokes}
% -------------------------------------------------------------------------
Una superficie $S$ con borde y su borde $\partial S$ son objetos orientados, y
sus orientaciones no son independientes.

\begin{definicion}[Orientación compatible del borde: regla de la mano derecha]\label{def:orientacion-compatible}
Dada una superficie orientada por el vector normal unitario continuo $\un$, la
\emph{orientación compatible} de su borde $\partial S$ es la del sentido de
circulación que, al acercarnos al borde, induce $\un$ según la \emph{regla del
pulgar}: con la mano derecha, el pulgar apunta según $\un$ y los demás dedos se
curvan en el sentido en que se recorre $\partial S$.
\end{definicion}

En la práctica: \emph{eligiendo el sentido en que se recorre la curva frontera
(cómo se orienta su parametrización) queda determinado el sentido de la normal}, y
viceversa.

\begin{teorema}[Teorema de Stokes]\label{teo:stokes}
Sea $\vc{A}$ un campo vectorial de clase $C^{1}$ en un entorno de una superficie
acotada y regular $S\subset\R^{3}$ cuyo borde $L=\partial S$ es una o varias
curvas cerradas. Si las orientaciones de $S$ y de $\partial S$ son compatibles,
\begin{equation}\label{eq:stokes}
  \boxed{\;\iint_{S}(\nab\times\vc{A})\cdot\dS=\oint_{\partial S}\vc{A}\cdot\dl\;}
\end{equation}
Es decir, el flujo del rotacional a través de $S$ es igual a la circulación del
campo a lo largo de su borde.
\end{teorema}
\begin{proof}[Idea de la demostración]
Sea primero $\Delta S=\Delta x\,\Delta y$ un pequeño rectángulo del plano $xy$ con
normal $\uz$ y borde recorrido en sentido antihorario. La circulación de
$\vc{A}$ alrededor de él es, a primer orden,
\begin{align*}
  \oint\vc{A}\cdot\dl&\approx A_x(x,y)\,\Delta x+A_y(x+\Delta x,y)\,\Delta y
     -A_x(x,y+\Delta y)\,\Delta x-A_y(x,y)\,\Delta y\\
  &\approx\Bigl(\pd{A_y}{x}-\pd{A_x}{y}\Bigr)\Delta x\,\Delta y
   =(\nab\times\vc{A})\cdot\uz\,\Delta S .
\end{align*}
Es decir, \emph{la circulación por unidad de área es la componente normal del
rotacional}. Una superficie cualquiera se subdivide en muchos de estos parches
infinitesimales orientados (\cref{fig:stokes}); como cada lado interior aparece
dos veces, con orientaciones opuestas, las circulaciones sobre las paredes
interiores se cancelan por parejas y solo sobrevive la circulación sobre el borde
exterior $\partial S$. Para parches curvos se cambia de variable mediante la
parametrización $\sigma$, que transforma $\uz$ en $\un$ (este es el teorema de
Green en el plano, llevado a la superficie)~\cite{NAVAS}.
\end{proof}

Para una región plana en el plano $xy$ (con $\un=\uz$), \cref{eq:stokes} se
reduce al \emph{teorema de Green}:
$\oint_{\partial S}P\,\dif x+Q\,\dif y=\iint_S\bigl(\partial_xQ-\partial_yP\bigr)\dif x\,\dif y$.

\begin{corolario}[Superficies cerradas]\label{cor:stokes-cerrada}
Si $S$ es una superficie cerrada ($\partial S=\varnothing$), el flujo del
rotacional de cualquier campo $\vc{A}$ a través de ella es nulo:
\begin{equation}\label{eq:flujo-rot-nulo}
  \oiint_{S}(\nab\times\vc{A})\cdot\dS=0 .
\end{equation}
\end{corolario}
\begin{proof}
Con $\partial S=\varnothing$ no hay circulación sobre el borde, y
\cref{eq:stokes} da $0$. El resultado es coherente con el teorema de Gauss
(\cref{teo:gauss}): si $S=\partial V$, entonces
$\oiint_S(\nab\times\vc{A})\cdot\dS=\iiint_V\nab\cdot(\nab\times\vc{A})\,\dV=0$ por
la \cref{prop:fundamentales}.
\end{proof}

Obsérvese qué ocurre en la frontera cuando $\partial S\neq\varnothing$: allí los
vórtices infinitesimales de los parches no tienen «vecino» con el que
cancelarse, por eso \emph{sí} queda la circulación sobre $L=\partial S$.

\begin{proposicion}[Rotacional y vorticidad]\label{prop:vorticidad}
La componente del rotacional según un eje unitario $\un$ en un punto $\vc{p}$ es
la circulación por unidad de área (\emph{vorticidad}) alrededor de ese eje:
\begin{equation}\label{eq:vorticidad}
  (\nab\times\vc{F})\cdot\un\,(\vc{p})=\lim_{G\to\{\vc{p}\}}\frac{1}{a(G)}\oint_{\partial G}\vc{F}\cdot\dl ,
\end{equation}
donde $G$ es un disco (o superficie similar) de área $a(G)$ contenido en un plano
ortogonal a $\un$ y que se contrae al punto $\vc{p}$.
\end{proposicion}
\begin{proof}
Por la \cref{eq:stokes}, $\frac{1}{a(G)}\oint_{\partial G}\vc{F}\cdot\dl=\frac{1}{a(G)}\iint_G(\nab\times\vc{F})\cdot\un\,\dsc$,
el valor medio de la componente normal del rotacional sobre $G$, que tiende a su
valor en $\vc{p}$ cuando $G$ se encoge a un punto (por la continuidad del
integrando).
\end{proof}

La \cref{eq:vorticidad} proporciona una definición del rotacional independiente
del sistema de coordenadas: tomando $\un=\ux,\uy,\uz$ se recuperan las tres
componentes. Con la ruedecita de paletas: la vorticidad alrededor de su eje es
proporcional a su velocidad angular de giro.

% Figura: teorema de Stokes. La circulación en los bordes interiores de los
% parches infinitesimales se cancela y solo queda la del borde L = dS.
\begin{figure}[ht]
\centering
\begin{tikzpicture}[>=Stealth, x=1cm, y=1cm]
  \def\blob{plot[smooth cycle, tension=0.75] coordinates {(-4.3,0.1) (-3.1,1.25) (-1.2,1.5)
     (0.9,1.05) (2.7,1.45) (4.1,0.8) (4.5,-0.25) (3.3,-1.0) (1.0,-1.3) (-1.4,-1.1) (-3.3,-0.9)}}
  \path[name path=superficie] \blob;
  % superficie con degradado azul
  \shade[left color=azulclaro!35, right color=azulclaro!85] \blob;
  \begin{scope}
    \clip \blob;
    % parches con su circulación (elipses por la perspectiva)
    \foreach \j in {-3,-2,...,3}{\foreach \i in {-8,-7,...,8}{%
      \pgfmathsetmacro{\cx}{0.55*\i+0.12*\j}\pgfmathsetmacro{\cy}{0.33*\j-0.03*\i*\i/8}%
      \draw[black!70, line width=0.55pt,
            decoration={markings, mark=at position 0.2 with {\arrow{Stealth[length=2.6pt]}},
                                  mark=at position 0.7 with {\arrow{Stealth[length=2.6pt]}}},
            postaction={decorate}] (\cx,\cy) ellipse (0.27 and 0.145);}}
  \end{scope}
  % borde L con orientación
  \draw[rojoteo, line width=1.5pt,
        decoration={markings, mark=between positions 0.04 and 0.97 step 0.1 with {\arrow{Stealth[length=6pt]}}},
        postaction={decorate}] \blob;
  \draw[azulacero, line width=0.5pt, densely dashed] (0,0.1) -- ++(0,0.05);
  \draw[->, line width=1.6pt, azulacero] (0.15,0.25) -- ++(0,1.9) node[above, text=azulacero] {$\dS$};
  \node[font=\Large, text=azulprof] at (-3.4,0.1) {$S$};
  \node[font=\Large, text=rojoteo, anchor=west] at (4.6,-0.55) {$L=\partial S$};
\end{tikzpicture}
\caption{Idea del teorema de Stokes. La superficie $S$ se divide en parches
infinitesimales orientados, cada uno con su circulación (elipses). En las
paredes interiores cada lado se recorre dos veces en sentidos opuestos y la
circulación se cancela; solo sobrevive la circulación sobre el borde exterior
$L=\partial S$ (en rojo). La normal $\dS$ y el sentido de recorrido de $L$ están
ligados por la regla de la mano derecha.}
\label{fig:stokes}
\end{figure}


% -------------------------------------------------------------------------
\subsection{El teorema de la divergencia (de Gauss)}\label{sec:gauss}
% -------------------------------------------------------------------------
\begin{definicion}[Borde de un sólido y su orientación]\label{def:borde-solido}
El borde $\partial V$ de un sólido regular $V\subset\R^{3}$ es una o varias
superficies cerradas. La orientación positiva de $\partial V$ es la determinada
por el campo normal unitario \emph{exterior} (el que apunta fuera de $V$). Si $V$
tiene huecos o «burbujas», la normal exterior respecto a $V$ es la normal
\emph{interior} respecto a la burbuja.
\end{definicion}

\begin{teorema}[Teorema de la divergencia o de Gauss]\label{teo:gauss}
Sea $\vc{A}$ un campo vectorial de clase $C^{1}$ en un entorno de un sólido
acotado $V\subset\R^{3}$ cuyo borde $S=\partial V$ es una o varias superficies
cerradas orientadas con la normal exterior. Entonces
\begin{equation}\label{eq:gauss}
  \boxed{\;\iiint_{V}\nab\cdot\vc{A}\,\dV=\oiint_{\partial V}\vc{A}\cdot\dS\;}
\end{equation}
El flujo de $\vc{A}$ a través de la superficie cerrada que limita a $V$ es la
integral de su divergencia en el interior.
\end{teorema}
\begin{proof}[Idea de la demostración]
La fórmula es lineal en $\vc{A}$, así que basta probarla para campos paralelos
a los ejes: sea $\vc{A}=R\,\uz$. Para el cubo unidad $V=[0,1]^{3}$, solo las
tapas $z=1$ (normal $+\uz$) y $z=0$ (normal $-\uz$) contribuyen, pues en las
otras cuatro caras $\vc{A}\cdot\un=0$:
\begin{align*}
  \oiint_{\partial V}\vc{A}\cdot\dS&=\int_0^1\!\!\int_0^1\bigl[R(x,y,1)-R(x,y,0)\bigr]\dif x\,\dif y\\
  &=\int_0^1\!\!\int_0^1\!\!\int_0^1\pd{R}{z}\,\dif z\,\dif x\,\dif y=\iiint_V\nab\cdot(R\,\uz)\,\dV ,
\end{align*}
donde se ha usado el teorema fundamental del cálculo en $z$. Para un sólido
general se pega una malla de cubitos (o se cambia de variable): cada cara
interior aparece dos veces con normales opuestas, y los flujos interiores se
cancelan, de modo que solo queda el flujo a través del borde exterior
(\cref{fig:gauss-vol}).
\end{proof}

\begin{proposicion}[La divergencia como densidad de flujo]\label{prop:div-intrinseca}
Para un campo $\vc{F}$ de clase $C^{1}$ y un punto $\vc{p}$,
\begin{equation}\label{eq:div-intrinseca}
  (\nab\cdot\vc{F})(\vc{p})=\lim_{V\to\{\vc{p}\}}\frac{1}{v(V)}\oiint_{\partial V}\vc{F}\cdot\un\,\dsc ,
\end{equation}
siendo $V$ un pequeño sólido (por ejemplo una bola centrada en $\vc{p}$) de volumen
$v(V)$ que se contrae a $\vc{p}$.
\end{proposicion}
\begin{proof}
Por el \cref{teo:gauss}, $\frac{1}{v(V)}\oiint_{\partial V}\vc{F}\cdot\un\,\dsc=\frac{1}{v(V)}\iiint_V\nab\cdot\vc{F}\,\dV$,
que es el valor medio de la divergencia sobre $V$ y converge a $(\nab\cdot\vc{F})(\vc{p})$
al contraer $V$ (continuidad).
\end{proof}

La \cref{eq:div-intrinseca} es la definición general de la divergencia,
independiente del sistema de coordenadas, y explica su nombre y su
interpretación: es el flujo neto que sale por unidad de volumen. El teorema de
Gauss es entonces la expresión matemática de las cancelaciones de flujo en las
paredes interiores: la integral de la divergencia es la «suma» de todas las
fuentes y sumideros del interior, y solo sobrevive el flujo por el borde. Un
campo de divergencia nula (solenoidal, o incompresible) tiene flujo neto nulo a
través de toda superficie cerrada.

% Figura: teorema de la divergencia. Campo radial y elemento de volumen en
% coordenadas esféricas.
% Elemento de volumen esférico: caras como polígonos con aristas muestreadas.
% \pt{r}{theta}{phi} define las coordenadas cartesianas en \X, \Y, \Z
\newcommand{\sph}[3]{({(#1)*sin(#2)*cos(#3)},{(#1)*sin(#2)*sin(#3)},{(#1)*cos(#2)})}
% cara a r fijo: #1 = r, recorre theta (ta..tb) y phi (pa..pb)
\newcommand{\caraR}[1]{%
  \fill[ocreteo, opacity=0.38]
  ({#1*sin(\ta)*cos(\pa)},{#1*sin(\ta)*sin(\pa)},{#1*cos(\ta)})
  \foreach \k in {1,...,6}{ -- ({#1*sin(\ta+\k*(\tb-\ta)/6)*cos(\pa)},{#1*sin(\ta+\k*(\tb-\ta)/6)*sin(\pa)},{#1*cos(\ta+\k*(\tb-\ta)/6)}) }
  \foreach \k in {1,...,6}{ -- ({#1*sin(\tb)*cos(\pa+\k*(\pb-\pa)/6)},{#1*sin(\tb)*sin(\pa+\k*(\pb-\pa)/6)},{#1*cos(\tb)}) }
  \foreach \k in {1,...,6}{ -- ({#1*sin(\tb-\k*(\tb-\ta)/6)*cos(\pb)},{#1*sin(\tb-\k*(\tb-\ta)/6)*sin(\pb)},{#1*cos(\tb-\k*(\tb-\ta)/6)}) }
  \foreach \k in {1,...,6}{ -- ({#1*sin(\ta)*cos(\pb-\k*(\pb-\pa)/6)},{#1*sin(\ta)*sin(\pb-\k*(\pb-\pa)/6)},{#1*cos(\ta)}) }
  -- cycle;}
% cara a theta fijo: #1 = theta
\newcommand{\caraT}[1]{%
  \fill[ocreteo, opacity=0.38]
  ({\ra*sin(#1)*cos(\pa)},{\ra*sin(#1)*sin(\pa)},{\ra*cos(#1)})
  -- ({\rb*sin(#1)*cos(\pa)},{\rb*sin(#1)*sin(\pa)},{\rb*cos(#1)})
  \foreach \k in {1,...,6}{ -- ({\rb*sin(#1)*cos(\pa+\k*(\pb-\pa)/6)},{\rb*sin(#1)*sin(\pa+\k*(\pb-\pa)/6)},{\rb*cos(#1)}) }
  -- ({\ra*sin(#1)*cos(\pb)},{\ra*sin(#1)*sin(\pb)},{\ra*cos(#1)})
  \foreach \k in {1,...,6}{ -- ({\ra*sin(#1)*cos(\pb-\k*(\pb-\pa)/6)},{\ra*sin(#1)*sin(\pb-\k*(\pb-\pa)/6)},{\ra*cos(#1)}) }
  -- cycle;}
% cara a phi fijo: #1 = phi
\newcommand{\caraP}[1]{%
  \fill[ocreteo, opacity=0.38]
  ({\ra*sin(\ta)*cos(#1)},{\ra*sin(\ta)*sin(#1)},{\ra*cos(\ta)})
  -- ({\rb*sin(\ta)*cos(#1)},{\rb*sin(\ta)*sin(#1)},{\rb*cos(\ta)})
  \foreach \k in {1,...,6}{ -- ({\rb*sin(\ta+\k*(\tb-\ta)/6)*cos(#1)},{\rb*sin(\ta+\k*(\tb-\ta)/6)*sin(#1)},{\rb*cos(\ta+\k*(\tb-\ta)/6)}) }
  -- ({\ra*sin(\tb)*cos(#1)},{\ra*sin(\tb)*sin(#1)},{\ra*cos(\tb)})
  \foreach \k in {1,...,6}{ -- ({\ra*sin(\tb-\k*(\tb-\ta)/6)*cos(#1)},{\ra*sin(\tb-\k*(\tb-\ta)/6)*sin(#1)},{\ra*cos(\tb-\k*(\tb-\ta)/6)}) }
  -- cycle;}
% arista a (r,theta) fijos recorriendo phi, etc.
\newcommand{\aristaP}[2]{\draw[rojoteo, dashed, line width=0.5pt]
  ({#1*sin(#2)*cos(\pa)},{#1*sin(#2)*sin(\pa)},{#1*cos(#2)})
  \foreach \k in {1,...,6}{ -- ({#1*sin(#2)*cos(\pa+\k*(\pb-\pa)/6)},{#1*sin(#2)*sin(\pa+\k*(\pb-\pa)/6)},{#1*cos(#2)}) };}
\newcommand{\aristaT}[2]{\draw[rojoteo, dashed, line width=0.5pt]
  ({#1*sin(\ta)*cos(#2)},{#1*sin(\ta)*sin(#2)},{#1*cos(\ta)})
  \foreach \k in {1,...,6}{ -- ({#1*sin(\ta+\k*(\tb-\ta)/6)*cos(#2)},{#1*sin(\ta+\k*(\tb-\ta)/6)*sin(#2)},{#1*cos(\ta+\k*(\tb-\ta)/6)}) };}
\begin{figure}[ht]
\centering
\tdplotsetmaincoords{66}{128}
\begin{minipage}[t]{0.48\textwidth}\centering
\begin{tikzpicture}[tdplot_main_coords, scale=0.78, baseline=(current bounding box.north)]
  \draw[eje] (0,0,0) -- (3.3,0,0) node[below left] {\small $x$};
  \draw[eje] (0,0,0) -- (0,3.3,0) node[right] {\small $y$};
  \draw[eje] (0,0,0) -- (0,0,3.3) node[above] {\small $z$};
  % campo radial (malla gruesa, módulo ~ 1/r^0: flechas unitarias coloreadas por r)
  \foreach \z in {-1.8,-0.6,0.6,1.8}{\foreach \y in {1.8,0.6,-0.6,-1.8}{\foreach \x in {-1.8,-0.6,0.6,1.8}{%
    \pgfmathsetmacro{\fm}{sqrt(\x*\x+\y*\y+\z*\z)}%
    \pgfmathsetmacro{\ft}{\fm/3.12}\mapcolor{\ft}%
    \draw[-{Stealth[length=3.4pt,width=3.4pt]}, line width=0.7pt, mapc]
      (\x,\y,\z) -- ++({0.36*\x/\fm*1.6},{0.36*\y/\fm*1.6},{0.36*\z/\fm*1.6});}}}
  % elemento de volumen: r in [1.4,2.5], theta in [38,62], phi in [8,42]
  \def\ra{1.4}\def\rb{2.5}\def\ta{38}\def\tb{62}\def\pa{8}\def\pb{42}
  \caraR{\ra}\caraR{\rb}
  \caraT{\ta}\caraT{\tb}
  \caraP{\pa}\caraP{\pb}
  \aristaP{\ra}{\ta}\aristaP{\ra}{\tb}\aristaP{\rb}{\ta}\aristaP{\rb}{\tb}
  \aristaT{\ra}{\pa}\aristaT{\ra}{\pb}\aristaT{\rb}{\pa}\aristaT{\rb}{\pb}
  \foreach \tt in {\ta,\tb}{\foreach \pp in {\pa,\pb}{\draw[rojoteo, dashed, line width=0.5pt]
      \sph{\ra}{\tt}{\pp} -- \sph{\rb}{\tt}{\pp};}}
\end{tikzpicture}\\[2pt]{\small\scshape (a)}
\end{minipage}\hfill
\begin{minipage}[t]{0.48\textwidth}\centering
\begin{tikzpicture}[>=Stealth, x=1cm, y=1cm, baseline=(current bounding box.north)]
  \foreach \a in {0,18,...,342}{%
    \draw[-{Stealth[length=5pt,width=5pt]}, azulacero, line width=1pt]
      (\a:0.55) -- (\a:1.65);
    \draw[-{Stealth[length=5pt,width=5pt]}, azulacero, line width=1pt]
      (\a:1.95) -- (\a:2.55);
    \draw[-{Stealth[length=5pt,width=5pt]}, azulacero!70, line width=1pt]
      (\a:2.85) -- (\a:3.2);}
  \fill[azulprof] (0,0) circle (1.6pt);
  % elemento de área (corte del elemento de volumen): sector anular
  \draw[rojoteo, dashed, line width=0.9pt, fill=ocreteo, fill opacity=0.25]
     (60:1.7) arc (60:84:1.7) -- (84:2.9) arc (84:60:2.9) -- cycle;
  \draw[rojoteo, line width=0.4pt] (0,0) -- (60:1.7);
  \node[font=\footnotesize, text=rojoteo] at (72:3.35) {$r\,\dif\theta$};
  \node[font=\footnotesize, text=rojoteo, anchor=west] at (50:2.35) {$\dif r$};
\end{tikzpicture}\\[2pt]{\small\scshape (b)}
\end{minipage}\\[6pt]
\begin{minipage}[t]{0.95\textwidth}\centering
\begin{tikzpicture}[>=Stealth, x=1cm, y=1cm, baseline=(current bounding box.north)]
  % elemento de volumen (paralelepípedo en proyección oblicua)
  \def\ax{2.8}\def\ay{1.7}\def\ox{1.0}\def\oy{0.55}
  \draw[rojoteo, dashed, line width=0.7pt] (\ox,\oy) -- (\ox+\ax,\oy) (\ox,\oy) -- (\ox,\oy+\ay) (\ox,\oy) -- (0,0);
  \fill[ocreteo, opacity=0.28] (0,0) rectangle (\ax,\ay);
  \draw[rojoteo, line width=0.9pt] (0,0) rectangle (\ax,\ay);
  \fill[ocreteo, opacity=0.18] (0,\ay) -- (\ax,\ay) -- (\ax+\ox,\ay+\oy) -- (\ox,\ay+\oy) -- cycle;
  \fill[ocreteo, opacity=0.40] (\ax,0) -- (\ax,\ay) -- (\ax+\ox,\ay+\oy) -- (\ax+\ox,\oy) -- cycle;
  \draw[rojoteo, line width=0.9pt] (0,\ay) -- (\ox,\ay+\oy) -- (\ax+\ox,\ay+\oy) -- (\ax+\ox,\oy) -- (\ax,0)
        (\ax,\ay) -- (\ax+\ox,\ay+\oy);
  % terna de versores locales
  \coordinate (T) at (-3.7,1.0);
  \draw[->, line width=1.2pt, verdeteo] (T) -- ++(0.95,0.52) node[above right, font=\small] {$\ur$};
  \draw[->, line width=1.2pt, verdeteo] (T) -- ++(0,-1.15) node[below, font=\small] {$\uth$};
  \draw[->, line width=1.2pt, verdeteo] (T) -- ++(1.15,0) node[right, font=\small] {$\uph$};
  \fill[verdeteo] (T) circle (1.3pt);
  % cotas
  \node[font=\small, text=rojoteo, anchor=north] at (\ax/2,-0.08) {$r\sin\theta\,\dif\varphi$};
  \node[font=\small, text=rojoteo, anchor=east] at (-0.08,\ay/2) {$r\,\dif\theta$};
  \node[font=\small, text=rojoteo, anchor=south west] at (\ax+0.45,\ay+0.5) {$\dif r$};
  \node[font=\small, text=black!80, anchor=west] at (\ax+\ox+0.5,\ay/2+\oy/2) {$\dV=r^{2}\sin\theta\,\dif r\,\dif\theta\,\dif\varphi$};
\end{tikzpicture}\\[2pt]{\small\scshape (c)}
\end{minipage}
\caption{Teorema de la divergencia con simetría radial. (a)~Campo radial
$\vc{A}=A(r)\,\ur$ y un elemento de volumen $\dV$ limitado por
$r,r+\dif r$; $\theta,\theta+\dif\theta$; $\varphi,\varphi+\dif\varphi$.
(b)~Corte meridiano: el elemento es un sector anular de lados $\dif r$ y
$r\,\dif\theta$. (c)~El elemento, ampliado: es un «ladrillo» curvilíneo de aristas
$\dif r$, $r\,\dif\theta$ y $r\sin\theta\,\dif\varphi$ según la terna ortonormal
$(\ur,\uth,\uph)$. El flujo neto por sus seis caras es $\nab\cdot\vc{A}\,\dV$.}
\label{fig:gauss-vol}
\end{figure}


% -------------------------------------------------------------------------
\subsection{Resumen y comprobación en campos sencillos}\label{sec:resumen-teoremas}
% -------------------------------------------------------------------------
Los tres teoremas, en la notación que se empleará en el resto del curso, son:
\begin{align}
  \int_{\vc{a}}^{\vc{b}}\nab f\cdot\dl&=f(\vc{b})-f(\vc{a}) &&\text{(teorema del gradiente)},\label{eq:resumen-grad}\\
  \int_{V}\nab\cdot\vc{A}\,\dV&=\oint_{S}\vc{A}\cdot\dS &&\text{(teorema de la divergencia, Gauss)},\label{eq:resumen-gauss}\\
  \int_{S}(\nab\times\vc{A})\cdot\dS&=\oint_{L}\vc{A}\cdot\dl &&\text{(teorema del rotacional, Stokes)},\label{eq:resumen-stokes}
\end{align}
con $S=\partial V$ en la segunda y $L=\partial S$ en la tercera. Se reconoce en
todos ellos la estructura «integral de la derivada $=$ integral sobre el borde».

\begin{ejemplo}[Comprobación en los campos de referencia]\label{ej:comprobacion-campos}
Tomamos los campos de la \cref{fig:campos_div} y los dos objetos de la
\cref{fig:campos_int}: un cubo unidad $V=[1,2]\times[1{,}5,2{,}5]\times[0,1]$ y un
circuito rectangular $C=\partial S$ en el plano $z=0$ con
$S=[1,2]\times[0{,}5,1]$, recorrido en sentido horario (normal $-\uz$, de
área $0{,}5$).
\begin{itemize}
  \item $\vc{F}_2=x\,\ux$ ($\nab\cdot\vc{F}_2=1$, $\nab\times\vc{F}_2=\vc{0}$).
        Flujo por las tapas $x=2$ y $x=1$: $(2-1)\cdot1=1=\iiint_V1\,\dV$;
        el resto de caras no contribuye. Circulación: los lados horizontales
        se cancelan y los verticales son perpendiculares al campo, luego es
        nula, de acuerdo con $\nab\times\vc{F}_2=\vc{0}$.
  \item $\vc{F}_3=y\,\ux$ ($\nab\cdot\vc{F}_3=0$, $\nab\times\vc{F}_3=-\uz$).
        El flujo por el cubo es nulo (lo que entra por una cara sale por la
        opuesta), y la circulación horaria es
        $\oint\vc{F}_3\cdot\dl=(1-0{,}5)\cdot1=0{,}5=\iint_S(-\uz)\cdot(-\uz)\,\dsc$
        (el lado superior, a $y=1$, está a favor del campo, el inferior, a
        $y=0{,}5$, en contra y más débil).
\end{itemize}
\end{ejemplo}

% Figura: los campos de referencia con un volumen V (cubo) y un circuito C.
\newcommand{\cuboV}[2]{%  esquina inferior izquierda (#1,#2): cubo de lado 1
  \fill[grisclaro] (#1,#2) rectangle (#1+1,#2+1);
  \fill[grispizarra!40] (#1,#2+1) -- (#1+0.35,#2+1.3) -- (#1+1.35,#2+1.3) -- (#1+1,#2+1) -- cycle;
  \fill[grispizarra!55] (#1+1,#2) -- (#1+1.35,#2+0.3) -- (#1+1.35,#2+1.3) -- (#1+1,#2+1) -- cycle;
  \draw[rojoteo, line width=0.8pt] (#1,#2) rectangle (#1+1,#2+1);
  \draw[rojoteo, line width=0.8pt] (#1,#2+1) -- (#1+0.35,#2+1.3) -- (#1+1.35,#2+1.3) -- (#1+1.35,#2+0.3) -- (#1+1,#2);
  \draw[rojoteo, line width=0.8pt] (#1+1,#2+1) -- (#1+1.35,#2+1.3);
  \node[font=\footnotesize, text=rojoteo] at (#1+0.5,#2+0.5) {$V$};}
\newcommand{\circuitoC}[4]{%  rectángulo [#1,#2]x[#3,#4] recorrido en sentido horario
  \fill[white, opacity=0.55] (#1,#3) rectangle (#2,#4);\draw[rojoteo, line width=1pt] (#1,#3) rectangle (#2,#4);
  \draw[-{Stealth[length=6pt]}, rojoteo, line width=1pt] ({(#1+#2)/2-0.15},#4) -- ({(#1+#2)/2+0.05},#4);
  \draw[-{Stealth[length=6pt]}, rojoteo, line width=1pt] (#2,{(#3+#4)/2+0.12}) -- (#2,{(#3+#4)/2-0.05});
  \draw[-{Stealth[length=6pt]}, rojoteo, line width=1pt] ({(#1+#2)/2+0.15},#3) -- ({(#1+#2)/2-0.05},#3);
  \draw[-{Stealth[length=6pt]}, rojoteo, line width=1pt] (#1,{(#3+#4)/2-0.12}) -- (#1,{(#3+#4)/2+0.05});
  \node[font=\footnotesize, text=rojoteo, anchor=west] at (#2+0.05,{(#3+#4)/2}) {$C$};}
\begin{figure}[ht]
\centering
\begin{minipage}[t]{0.325\textwidth}\centering
\begin{tikzpicture}[x=0.68cm, y=0.68cm, baseline=(current bounding box.north)]
  \marcoD{-2.7}{2.7}{-2.7}{2.7}{-2,-1,0,1,2}{-2,-1,0,1,2}
  \ejesD{-2.7}{2.7}{-2.7}{2.7}
  \foreach \x in {-2.4,-1.8,...,2.4}{\foreach \y in {-2.4,-1.8,...,2.4}{%
    \draw[-{Stealth[length=4.5pt,width=4.5pt]}, line width=0.7pt, azulprof]
      (\x-0.17,\y-0.17) -- ++(0.34,0.34);}}
\end{tikzpicture}\\[3pt]{\small $\vc{F}_1=\ux+\uy$}
\end{minipage}\hfill
\begin{minipage}[t]{0.325\textwidth}\centering
\begin{tikzpicture}[x=1.12cm, y=1.12cm, baseline=(current bounding box.north)]
  \marcoD{0}{3.2}{0}{3.2}{0,1,2,3}{0,1,2,3}
  \ejesD{0}{3.2}{0}{3.2}
  \clipD{0}{3.2}{0}{3.2}{\flechas{\x}{0}{0.2,0.575,...,3.1}{0.2,0.575,...,3.1}{0.30}{3.1}}
  \cuboV{1}{1.5}
  \circuitoC{1}{2}{0.5}{1}
\end{tikzpicture}\\[3pt]{\small $\vc{F}_2=x\,\ux$}
\end{minipage}\hfill
\begin{minipage}[t]{0.325\textwidth}\centering
\begin{tikzpicture}[x=1.12cm, y=1.12cm, baseline=(current bounding box.north)]
  \marcoD{0}{3.2}{0}{3.2}{0,1,2,3}{0,1,2,3}
  \ejesD{0}{3.2}{0}{3.2}
  \clipD{0}{3.2}{0}{3.2}{\flechas{\y}{0}{0.2,0.575,...,3.1}{0.2,0.575,...,3.1}{0.30}{3.1}}
  \cuboV{1}{1.5}
  \circuitoC{1}{2}{0.5}{1}
\end{tikzpicture}\\[3pt]{\small $\vc{F}_3=y\,\ux$}
\end{minipage}
\caption{Campos de referencia con un volumen $V$ (cubo) y un circuito $C$
(rectángulo recorrido en sentido horario). Para $\vc{F}_2$ el flujo neto por la
superficie del cubo es no nulo ($=1$) y la circulación por $C$ es nula; para
$\vc{F}_3$ ocurre lo contrario: el flujo es nulo y la circulación vale $0{,}5$.}
\label{fig:campos_int}
\end{figure}


\begin{ejemplo}[Rotacional uniforme: la circulación solo depende del área]\label{ej:rot-uniforme}
Los campos $\vc{F}=y\,\ux-x\,\uy$ y $\vc{F}=y\,\ux$ tienen rotacional
\emph{uniforme} (el mismo en todos los puntos): $-2\,\uz$ y $-\uz$,
respectivamente. Por el teorema de Stokes, la circulación horaria (normal $-\uz$)
alrededor de un circuito de área $a$ vale $\oint\vc{F}\cdot\dl=\iint(\nab\times\vc{F})\cdot(-\uz)\,\dsc=2a$
(resp. $a$), \emph{sin importar dónde esté colocado el circuito}. En la
\cref{fig:campos_circ}, los dos rectángulos de área $a=1{,}4$ dan $2{,}8$ (resp.
$1{,}4$) aunque el campo valga cosas muy distintas sobre cada uno. Se comprueba
directamente en el segundo caso: el campo es horizontal, los lados verticales
no aportan, y los horizontales dan $(y_{\mathrm{sup}}-y_{\mathrm{inf}})\cdot1{,}4=1\cdot1{,}4=1{,}4$,
con independencia de la altura $y$ a la que esté el rectángulo.
\end{ejemplo}

% Figura: campos de rotacional uniforme y dos circuitos iguales en posiciones distintas.
\begin{figure}[ht]
\centering
\begin{minipage}[t]{0.47\textwidth}\centering
\begin{tikzpicture}[x=0.74cm, y=0.74cm, baseline=(current bounding box.north)]
  \marcoD{-2.7}{2.7}{-2.7}{2.7}{-2,-1,0,1,2}{-2,-1,0,1,2}
  \ejesD{-2.7}{2.7}{-2.7}{2.7}
  \clipD{-2.7}{2.7}{-2.7}{2.7}{\flechas{\y}{-\x}{-2.4,-1.8,...,2.4}{-2.4,-1.8,...,2.4}{0.34}{3.4}}
  \circuitoC{-0.7}{0.7}{-0.45}{0.55}
  \circuitoC{-0.7}{0.7}{-2.0}{-1.0}
\end{tikzpicture}\\[3pt]{\small $\vc{F}=y\,\ux-x\,\uy$,\quad $\nab\times\vc{F}=-2\,\uz$}
\end{minipage}\hfill
\begin{minipage}[t]{0.47\textwidth}\centering
\begin{tikzpicture}[x=0.74cm, y=0.74cm, baseline=(current bounding box.north)]
  \marcoD{-2.7}{2.7}{-2.7}{2.7}{-2,-1,0,1,2}{-2,-1,0,1,2}
  \ejesD{-2.7}{2.7}{-2.7}{2.7}
  \clipD{-2.7}{2.7}{-2.7}{2.7}{\flechas{\y}{0}{-2.4,-1.8,...,2.4}{-2.4,-1.8,...,2.4}{0.34}{2.5}}
  \circuitoC{-0.7}{0.7}{-0.45}{0.55}
  \circuitoC{-0.7}{0.7}{-2.0}{-1.0}
\end{tikzpicture}\\[3pt]{\small $\vc{F}=y\,\ux$,\quad $\nab\times\vc{F}=-\uz$}
\end{minipage}
\caption{Campos de rotacional uniforme y dos circuitos rectangulares iguales
(área $1{,}4$, recorridos en sentido horario) en posiciones distintas. Aunque
el campo es muy diferente sobre uno y otro circuito, la circulación es la misma
en ambos: $2{,}8$ para $\vc{F}=y\,\ux-x\,\uy$ y $1{,}4$ para $\vc{F}=y\,\ux$.
Solo depende del rotacional y del área encerrada.}
\label{fig:campos_circ}
\end{figure}



% =========================================================================
\section{Sistemas de coordenadas y operadores diferenciales}\label{sec:coordenadas}
% =========================================================================
Las leyes del electromagnetismo no dependen de los ejes que elijamos, pero los
\emph{cálculos} sí: ante una distribución con simetría esférica o cilíndrica,
lo inteligente es usar coordenadas adaptadas a esa simetría. En esta sección
repasamos los tres sistemas ortogonales fundamentales (cartesiano, cilíndrico y
esférico) con sus elementos de línea, de área y de volumen, sus cambios de base,
y los cuatro operadores diferenciales en cada uno de ellos~\cite{FORMULARIO}.
Los obtenemos todos a la vez a partir de un marco general.

% -------------------------------------------------------------------------
\subsection{Coordenadas curvilíneas ortogonales}\label{sec:curvilineas}
% -------------------------------------------------------------------------
Un sistema de coordenadas curvilíneas $(q_1,q_2,q_3)$ viene dado por el
vector de posición $\vc{r}=\vc{r}(q_1,q_2,q_3)$. Variando solo $q_i$ se
recorre la \emph{línea coordenada} $i$, cuya tangente es $\partial\vc{r}/\partial q_i$.

\begin{definicion}[Coordenadas ortogonales y factores de escala]\label{def:ortogonales}
El sistema es \emph{ortogonal} si los tres vectores $\partial\vc{r}/\partial q_i$
son perpendiculares dos a dos. Se definen entonces los \emph{factores de escala}
$h_i$ y los versores $\uh{i}$ por
\[
  h_i\defeq\Bigl\lVert\pd{\vc{r}}{q_i}\Bigr\rVert,
  \qquad
  \uh{i}\defeq\frac{1}{h_i}\pd{\vc{r}}{q_i},
\]
y $(\uh{1},\uh{2},\uh{3})$ es una base ortonormal \emph{local} (depende del
punto), que supondremos orientada a derechas: $\uh{1}\times\uh{2}=\uh{3}$.
\end{definicion}

El factor $h_i$ convierte un incremento de la coordenada, $\dif q_i$, en una
\emph{longitud} $h_i\,\dif q_i$ (por ejemplo, $\dif\varphi$ es un ángulo y
$\rho\,\dif\varphi$ un arco). De ahí se obtienen los elementos geométricos.

\begin{proposicion}[Elementos de línea, área y volumen]\label{prop:elementos}
En coordenadas ortogonales,
\begin{align}
  \dl&=h_1\,\dif q_1\,\uh{1}+h_2\,\dif q_2\,\uh{2}+h_3\,\dif q_3\,\uh{3}, \label{eq:dl-gen}\\
  \dif S_i&=h_jh_k\,\dif q_j\,\dif q_k\quad(i,j,k\ \text{cíclicos}), \label{eq:dS-gen}\\
  \dV&=h_1h_2h_3\,\dif q_1\,\dif q_2\,\dif q_3, \label{eq:dV-gen}
\end{align}
donde $\dif S_i$ es el área de la cara del elemento de volumen perpendicular a
$\uh{i}$ (a $q_i$ constante). El jacobiano del cambio de variable
$(q_1,q_2,q_3)\mapsto(x,y,z)$ es
$\bigl|\partial(x,y,z)/\partial(q_1,q_2,q_3)\bigr|=h_1h_2h_3$.
\end{proposicion}
\begin{proof}
Al variar $q_i$ en $\dif q_i$ el punto se desplaza
$(\partial\vc{r}/\partial q_i)\dif q_i=h_i\,\dif q_i\,\uh{i}$, y el desplazamiento
total es la suma, \cref{eq:dl-gen}. El elemento de volumen es el paralelepípedo
rectangular de aristas $h_i\,\dif q_i$, lo que da \eqref{eq:dV-gen}, y cada cara es
el rectángulo de las otras dos aristas, \cref{eq:dS-gen}. El jacobiano es el
determinante de la matriz cuyas columnas son $\partial\vc{r}/\partial q_i=h_i\uh{i}$,
que, por ser ortogonales, vale el producto de sus longitudes.
\end{proof}

\begin{teorema}[Operadores diferenciales en coordenadas ortogonales]\label{teo:operadores-ortogonales}
Sean $f$ un campo escalar y $\vc{A}=A_1\uh{1}+A_2\uh{2}+A_3\uh{3}$ un campo vectorial.
Con $\partial_i\equiv\partial/\partial q_i$ y $H=h_1h_2h_3$,
\begin{align}
  \nab f&=\sum_{i=1}^{3}\frac{1}{h_i}\,\partial_if\;\uh{i}, \label{eq:grad-gen}\\
  \nab\cdot\vc{A}&=\frac{1}{H}\Bigl[\partial_1(h_2h_3A_1)+\partial_2(h_3h_1A_2)+\partial_3(h_1h_2A_3)\Bigr], \label{eq:div-gen}\\
  \nab\times\vc{A}&=\frac{1}{H}
  \begin{vmatrix}
    h_1\uh{1} & h_2\uh{2} & h_3\uh{3}\\
    \partial_1 & \partial_2 & \partial_3\\
    h_1A_1 & h_2A_2 & h_3A_3
  \end{vmatrix}, \label{eq:rot-gen}\\
  \Lap f&=\frac{1}{H}\Bigl[\partial_1\Bigl(\frac{h_2h_3}{h_1}\partial_1f\Bigr)
     +\partial_2\Bigl(\frac{h_3h_1}{h_2}\partial_2f\Bigr)
     +\partial_3\Bigl(\frac{h_1h_2}{h_3}\partial_3f\Bigr)\Bigr]. \label{eq:lap-gen}
\end{align}
\end{teorema}
\begin{proof}
\emph{Gradiente.} Por un lado $\dif f=\sum_i\partial_if\,\dif q_i$; por otro, por
la \cref{eq:df-compacta} y la \cref{eq:dl-gen},
$\dif f=\nab f\cdot\dl=\sum_i(\nab f)_i\,h_i\,\dif q_i$. Comparando coeficientes de los
$\dif q_i$ independientes, $(\nab f)_i=\partial_if/h_i$.

\emph{Divergencia.} Usamos la \cref{prop:div-intrinseca} sobre el elemento de
volumen curvilíneo. El flujo por las dos caras perpendiculares a $\uh{1}$ (de
área $h_2h_3\,\dif q_2\,\dif q_3$, evaluada en $q_1$ y en $q_1+\dif q_1$) es
\[
  \bigl[A_1h_2h_3\bigr]_{q_1+\dif q_1}\dif q_2\dif q_3-\bigl[A_1h_2h_3\bigr]_{q_1}\dif q_2\dif q_3
  \approx\partial_1(h_2h_3A_1)\,\dif q_1\dif q_2\dif q_3 ,
\]
y análogamente para las otras dos parejas de caras. Dividiendo el flujo total
por $\dV=H\,\dif q_1\dif q_2\dif q_3$ se obtiene \eqref{eq:div-gen}.

\emph{Rotacional.} Usamos la \cref{prop:vorticidad} sobre la cara perpendicular a
$\uh{1}$, de área $h_2h_3\,\dif q_2\,\dif q_3$. Su circulación (lados de longitud
$h_2\dif q_2$ y $h_3\dif q_3$) es, a primer orden,
$\bigl[\partial_2(h_3A_3)-\partial_3(h_2A_2)\bigr]\dif q_2\dif q_3$; dividiendo por el
área, $(\nab\times\vc{A})_1=\tfrac{1}{h_2h_3}\bigl[\partial_2(h_3A_3)-\partial_3(h_2A_2)\bigr]$.
Las otras dos componentes se obtienen permutando cíclicamente los índices, y el
conjunto es el desarrollo del determinante \eqref{eq:rot-gen}.

\emph{Laplaciano.} Es $\nab\cdot(\nab f)$: basta sustituir $A_i=\partial_if/h_i$
en la \cref{eq:div-gen}.
\end{proof}

Los factores de escala de los tres sistemas que usaremos son
\begin{center}
\renewcommand{\arraystretch}{1.3}
\begin{tabular}{@{}lccc@{}}
\toprule
Sistema & $(q_1,q_2,q_3)$ & $(h_1,h_2,h_3)$ & $H=h_1h_2h_3$\\
\midrule
cartesiano & $(x,y,z)$ & $(1,\,1,\,1)$ & $1$\\
cilíndrico & $(\rho,\varphi,z)$ & $(1,\,\rho,\,1)$ & $\rho$\\
esférico & $(r,\theta,\varphi)$ & $(1,\,r,\,r\sen\theta)$ & $r^{2}\sen\theta$\\
\bottomrule
\end{tabular}
\end{center}

% -------------------------------------------------------------------------
\subsection{Coordenadas cartesianas}\label{sec:cartesianas}
% -------------------------------------------------------------------------
\begin{figure}[ht]
\centering
\tdplotsetmaincoords{72}{120}
\begin{tikzpicture}[tdplot_main_coords, scale=1.45]
  \def\px{1.8}\def\py{2.3}\def\pz{2.1}
  \draw[eje] (0,0,0) -- (3.3,0,0) node[below left] {\small $x$};
  \draw[eje] (0,0,0) -- (0,3.6,0) node[right] {\small $y$};
  \draw[eje] (0,0,0) -- (0,0,3.2) node[above] {\small $z$};
  \fill[ocreteo, opacity=0.18] (0,0,0) -- (3.0,0,0) -- (3.0,3.2,0) -- (0,3.2,0) -- cycle;
  \draw[punteado] (\px,0,0) -- (\px,\py,0) -- (0,\py,0) (\px,\py,0) -- (\px,\py,\pz);
  \draw[punteado] (0,0,\pz) -- (\px,0,\pz) -- (\px,\py,\pz) (0,\py,0) -- (0,\py,\pz) -- (\px,\py,\pz);
  \draw[->, line width=1.5pt, azulprof] (0,0,0) -- (\px,\py,\pz);
  \node[azulprof, font=\small] at (0.55,0.5,1.6) {$\vc{r}$};
  \uvect{\px}{\py}{\pz}{(0.95,0,0)}{right}
  \uvect{\px}{\py}{\pz}{(0,0.95,0)}{right}
  \uvect{\px}{\py}{\pz}{(0,0,0.95)}{above}
  \node[rojoteo, font=\scriptsize] at (3.15,2.1,1.7) {$\ux$};
  \node[rojoteo, font=\scriptsize] at (2.0,3.15,1.7) {$\uy$};
  \node[rojoteo, font=\scriptsize] at (2.0,2.1,2.95) {$\uz$};
  \fill[azulprof] (\px,\py,\pz) circle (1.1pt);
  \node[font=\scriptsize] at (\px,-0.25,0) {$x$};
  \node[font=\scriptsize] at (-0.25,\py,0) {$y$};
  \node[font=\scriptsize] at (0,-0.15,\pz) {$z$};
\end{tikzpicture}
\caption{Coordenadas cartesianas $(x,y,z)$ y su base ortonormal constante $(\ux,\uy,\uz)$ en un punto $P$ de vector de posición $\vc{r}=x\,\ux+y\,\uy+z\,\uz$. Las líneas discontinuas muestran las proyecciones de $P$ sobre los ejes y los planos coordenados.}
\label{fig:coord-cart}
\end{figure}

Son las coordenadas $(x,y,z)$ asociadas a una base ortonormal \emph{constante}
$(\ux,\uy,\uz)$. Con $h_i=1$:
\begin{align}
  \vc{r}&=x\,\ux+y\,\uy+z\,\uz, & \dl&=\dif x\,\ux+\dif y\,\uy+\dif z\,\uz,\nonumber\\
  \dif S_x&=\dif y\,\dif z,\quad\dif S_y=\dif x\,\dif z,\quad\dif S_z=\dif x\,\dif y, & \dV&=\dif x\,\dif y\,\dif z.\nonumber
\end{align}
Los operadores son los ya conocidos:
\begin{align}
  \nab f&=\pd{f}{x}\,\ux+\pd{f}{y}\,\uy+\pd{f}{z}\,\uz, &
  \nab\cdot\vc{F}&=\pd{F_x}{x}+\pd{F_y}{y}+\pd{F_z}{z},\nonumber\\
  \nab\times\vc{F}&=\Bigl(\pd{F_z}{y}-\pd{F_y}{z}\Bigr)\ux+\Bigl(\pd{F_x}{z}-\pd{F_z}{x}\Bigr)\uy+\Bigl(\pd{F_y}{x}-\pd{F_x}{y}\Bigr)\uz, &
  \Lap f&=\pdd{f}{x}+\pdd{f}{y}+\pdd{f}{z}.\nonumber
\end{align}

% -------------------------------------------------------------------------
\subsection{Coordenadas cilíndricas}\label{sec:cilindricas}
% -------------------------------------------------------------------------
Se usan $(\rho,\varphi,z)$ con $\rho\ge0$ la distancia al eje $z$,
$\varphi\in[0,2\pi)$ el ángulo azimutal y $z\in\R$. (En el formulario del curso
la distancia al eje se denota $r$; aquí empleamos $\rho$ para no confundirla con
la distancia al origen $r$ de las coordenadas esféricas.)

\begin{figure}[ht]
\centering
\tdplotsetmaincoords{72}{120}
\begin{tikzpicture}[tdplot_main_coords, scale=1.45]
  \def\rh{2.3}\def\ph{38}\def\pz{1.7}
  \pgfmathsetmacro{\px}{\rh*cos(\ph)}\pgfmathsetmacro{\py}{\rh*sin(\ph)}
  \draw[eje] (0,0,0) -- (3.3,0,0) node[below left] {\small $x$};
  \draw[eje] (0,0,0) -- (0,3.6,0) node[right] {\small $y$};
  \draw[eje] (0,0,0) -- (0,0,3.2) node[above] {\small $z$};
  % pared cilíndrica
  \fill[ocreteo, opacity=0.20]
    plot[variable=\t, domain=-15:105, samples=14] ({\rh*cos(\t)},{\rh*sin(\t)},-0.3)
    -- plot[variable=\t, domain=105:-15, samples=14] ({\rh*cos(\t)},{\rh*sin(\t)},2.6) -- cycle;
  \draw[ocreteo!70, line width=0.4pt] plot[variable=\t, domain=-15:105, samples=14] ({\rh*cos(\t)},{\rh*sin(\t)},\pz);
  \draw[punteado] (0,0,0) -- (\px,\py,0) -- (\px,\py,\pz);
  \draw[->, line width=1.5pt, azulprof] (0,0,0) -- (\px,\py,\pz);
  \node[azulprof, font=\small] at (0.35,0.1,1.45) {$\vc{r}$};
  \draw[grispizarra, line width=0.5pt] plot[variable=\t, domain=0:\ph, samples=8] ({0.9*cos(\t)},{0.9*sin(\t)},0);
  \node[font=\scriptsize] at ({1.25*cos(\ph/2)},{1.25*sin(\ph/2)},0) {$\varphi$};
  \node[font=\scriptsize] at ({0.72*\rh*cos(\ph)+0.30*sin(\ph)},{0.72*\rh*sin(\ph)-0.30*cos(\ph)},0) {$\rho$};
  \node[font=\scriptsize] at (\px+0.05,\py+0.25,0.5*\pz) {$z$};
  \uvect{\px}{\py}{\pz}{({0.95*cos(\ph)},{0.95*sin(\ph)},0)}{right}
  \uvect{\px}{\py}{\pz}{({-0.95*sin(\ph)},{0.95*cos(\ph)},0)}{above}
  \uvect{\px}{\py}{\pz}{(0,0,0.95)}{above}
  \node[rojoteo, font=\scriptsize] at ({\px+1.25*cos(\ph)},{\py+1.25*sin(\ph)},\pz-0.1) {$\urho$};
  \node[rojoteo, font=\scriptsize] at ({\px-1.2*sin(\ph)},{\py+1.2*cos(\ph)},\pz+0.05) {$\uph$};
  \node[rojoteo, font=\scriptsize] at (\px,\py,\pz+1.25) {$\uz$};
  \fill[azulprof] (\px,\py,\pz) circle (1.8pt);
\end{tikzpicture}
\caption{Coordenadas cilíndricas $(\rho,\varphi,z)$. El versor $\urho$ apunta alejándose del eje $z$ y $\uph$ es tangente a la circunferencia de radio $\rho$ (en el sentido en que crece $\varphi$); $\uz$ es el mismo que en cartesianas. La pared cilíndrica es la superficie $\rho=\text{cte}$. Al cambiar de punto, $\urho$ y $\uph$ giran.}
\label{fig:coord-cil}
\end{figure}


\begin{align}
  x&=\rho\cos\varphi, & y&=\rho\sen\varphi, & z&=z,\nonumber\\
  \rho&=\sqrt{x^{2}+y^{2}}, & \varphi&=\tan^{-1}(y/x), & z&=z .\nonumber
\end{align}
Los versores y su relación con la base cartesiana son
\begin{align}
  \urho&=\cos\varphi\,\ux+\sen\varphi\,\uy, & \ux&=\cos\varphi\,\urho-\sen\varphi\,\uph,\nonumber\\
  \uph&=-\sen\varphi\,\ux+\cos\varphi\,\uy, & \uy&=\sen\varphi\,\urho+\cos\varphi\,\uph,\label{eq:base-cil}\\
  \uz&=\uz, & \uz&=\uz.\nonumber
\end{align}
La posición es $\vc{r}=\rho\,\urho+z\,\uz$ y los elementos geométricos
(\cref{prop:elementos} con $h=(1,\rho,1)$):
\begin{align}
  \dl&=\dif\rho\,\urho+\rho\,\dif\varphi\,\uph+\dif z\,\uz, \nonumber\\
  \dif S_\rho&=\rho\,\dif\varphi\,\dif z,\quad
  \dif S_\varphi=\dif\rho\,\dif z,\quad
  \dif S_z=\rho\,\dif\rho\,\dif\varphi,\qquad
  \dV=\rho\,\dif\rho\,\dif\varphi\,\dif z .\nonumber
\end{align}
Los operadores (\cref{teo:operadores-ortogonales}) son
\begin{align}
  \nab f&=\pd{f}{\rho}\,\urho+\frac{1}{\rho}\pd{f}{\varphi}\,\uph+\pd{f}{z}\,\uz,\nonumber\\
  \nab\cdot\vc{F}&=\frac{1}{\rho}\pd{(\rho F_\rho)}{\rho}+\frac{1}{\rho}\pd{F_\varphi}{\varphi}+\pd{F_z}{z},\nonumber\\
  \nab\times\vc{F}&=\Bigl(\frac{1}{\rho}\pd{F_z}{\varphi}-\pd{F_\varphi}{z}\Bigr)\urho
    +\Bigl(\pd{F_\rho}{z}-\pd{F_z}{\rho}\Bigr)\uph
    +\frac{1}{\rho}\Bigl(\pd{(\rho F_\varphi)}{\rho}-\pd{F_\rho}{\varphi}\Bigr)\uz,\label{eq:ops-cil}\\
  \Lap f&=\frac{1}{\rho}\pd{}{\rho}\Bigl(\rho\pd{f}{\rho}\Bigr)+\frac{1}{\rho^{2}}\pdd{f}{\varphi}+\pdd{f}{z}.\nonumber
\end{align}

% -------------------------------------------------------------------------
\subsection{Coordenadas esféricas}\label{sec:esfericas}
% -------------------------------------------------------------------------
Se usan $(r,\theta,\varphi)$ con $r\ge0$ la distancia al origen,
$\theta\in[0,\pi]$ el ángulo polar (medido desde el eje $z$) y
$\varphi\in[0,2\pi)$ el azimut. Obsérvese que $\sen\theta\ge0$ en todo el
rango, hecho que usaremos al calcular jacobianos.

\begin{figure}[ht]
\centering
\tdplotsetmaincoords{72}{120}
\begin{tikzpicture}[tdplot_main_coords, scale=1.45]
  \def\rr{3.0}\def\th{52}\def\ph{38}
  \pgfmathsetmacro{\px}{\rr*sin(\th)*cos(\ph)}\pgfmathsetmacro{\py}{\rr*sin(\th)*sin(\ph)}\pgfmathsetmacro{\pz}{\rr*cos(\th)}
  \draw[eje] (0,0,0) -- (3.5,0,0) node[below left] {\small $x$};
  \draw[eje] (0,0,0) -- (0,3.7,0) node[right] {\small $y$};
  \draw[eje] (0,0,0) -- (0,0,3.5) node[above] {\small $z$};
  % casquete esférico
  \fill[ocreteo, opacity=0.20]
    plot[variable=\t, domain=0:90, samples=10] ({\rr*sin(\t)*cos(-10)},{\rr*sin(\t)*sin(-10)},{\rr*cos(\t)})
    -- plot[variable=\t, domain=-10:100, samples=14] ({\rr*cos(\t)},{\rr*sin(\t)},0)
    -- plot[variable=\t, domain=90:0, samples=10] ({\rr*sin(\t)*cos(100)},{\rr*sin(\t)*sin(100)},{\rr*cos(\t)}) -- cycle;
  \draw[punteado] (0,0,0) -- (\px,\py,0) -- (\px,\py,\pz);
  \draw[->, line width=1.5pt, azulprof] (0,0,0) -- (\px,\py,\pz);
  \node[azulprof, font=\small] at ({0.45*\px+0.25},{0.45*\py-0.1},{0.45*\pz+0.25}) {$\vc{r}$};
  % arcos de theta y phi
  \draw[grispizarra, line width=0.5pt] plot[variable=\t, domain=0:\th, samples=8]
        ({1.0*sin(\t)*cos(\ph)},{1.0*sin(\t)*sin(\ph)},{1.0*cos(\t)});
  \node[font=\scriptsize] at ({1.35*sin(\th/2)*cos(\ph)},{1.35*sin(\th/2)*sin(\ph)},{1.35*cos(\th/2)}) {$\theta$};
  \draw[grispizarra, line width=0.5pt] plot[variable=\t, domain=0:\ph, samples=8] ({1.0*cos(\t)},{1.0*sin(\t)},0);
  \node[font=\scriptsize] at ({1.35*cos(\ph/2)},{1.35*sin(\ph/2)},0) {$\varphi$};
  \uvect{\px}{\py}{\pz}{({0.95*sin(\th)*cos(\ph)},{0.95*sin(\th)*sin(\ph)},{0.95*cos(\th)})}{above}
  \uvect{\px}{\py}{\pz}{({0.95*cos(\th)*cos(\ph)},{0.95*cos(\th)*sin(\ph)},{-0.95*sin(\th)})}{below}
  \uvect{\px}{\py}{\pz}{({-0.95*sin(\ph)},{0.95*cos(\ph)},0)}{right}
  \node[rojoteo, font=\scriptsize] at ({\px+1.2*sin(\th)*cos(\ph)},{\py+1.2*sin(\th)*sin(\ph)},{\pz+1.2*cos(\th)}) {$\ur$};
  \node[rojoteo, font=\scriptsize] at ({\px+1.25*cos(\th)*cos(\ph)},{\py+1.25*cos(\th)*sin(\ph)},{\pz-1.25*sin(\th)}) {$\uth$};
  \node[rojoteo, font=\scriptsize] at ({\px-1.2*sin(\ph)},{\py+1.2*cos(\ph)},\pz) {$\uph$};
  \fill[azulprof] (\px,\py,\pz) circle (1.1pt);
\end{tikzpicture}
\caption{Coordenadas esféricas $(r,\theta,\varphi)$. El versor $\ur$ es radial, $\uth$ apunta hacia la dirección en que aumenta $\theta$ (hacia el sur) y $\uph$ es el mismo versor azimutal de las cilíndricas. El casquete es la superficie $r=\text{cte}$; el vector de posición es $\vc{r}=r\,\ur$.}
\label{fig:coord-esf}
\end{figure}

\begin{align}
  x&=r\sen\theta\cos\varphi, & y&=r\sen\theta\sen\varphi, & z&=r\cos\theta,\nonumber\\
  r&=\sqrt{x^{2}+y^{2}+z^{2}}, & \theta&=\tan^{-1}\!\Bigl(\sqrt{x^{2}+y^{2}}/z\Bigr), & \varphi&=\tan^{-1}(y/x).\nonumber
\end{align}
Los versores esféricos en función de los cartesianos, y viceversa, son
\begin{align}
  \ur&=\sen\theta\cos\varphi\,\ux+\sen\theta\sen\varphi\,\uy+\cos\theta\,\uz,\nonumber\\
  \uth&=\cos\theta\cos\varphi\,\ux+\cos\theta\sen\varphi\,\uy-\sen\theta\,\uz,\nonumber\\
  \uph&=-\sen\varphi\,\ux+\cos\varphi\,\uy,\label{eq:base-esf}\\[4pt]
  \ux&=\sen\theta\cos\varphi\,\ur+\cos\theta\cos\varphi\,\uth-\sen\varphi\,\uph,\nonumber\\
  \uy&=\sen\theta\sen\varphi\,\ur+\cos\theta\sen\varphi\,\uth+\cos\varphi\,\uph,\nonumber\\
  \uz&=\cos\theta\,\ur-\sen\theta\,\uth.\nonumber
\end{align}
La posición es simplemente $\vc{r}=r\,\ur$ (¡la dirección está en $\ur$!) y
\begin{align}
  \dl&=\dif r\,\ur+r\,\dif\theta\,\uth+r\sen\theta\,\dif\varphi\,\uph, \nonumber\\
  \dif S_r&=r^{2}\sen\theta\,\dif\theta\,\dif\varphi,\quad
  \dif S_\theta=r\sen\theta\,\dif r\,\dif\varphi,\quad
  \dif S_\varphi=r\,\dif r\,\dif\theta,\qquad
  \dV=r^{2}\sen\theta\,\dif r\,\dif\theta\,\dif\varphi .\nonumber
\end{align}
Los operadores son
\begin{align}
  \nab f&=\pd{f}{r}\,\ur+\frac{1}{r}\pd{f}{\theta}\,\uth+\frac{1}{r\sen\theta}\pd{f}{\varphi}\,\uph,\nonumber\\
  \nab\cdot\vc{F}&=\frac{1}{r^{2}}\pd{(r^{2}F_r)}{r}+\frac{1}{r\sen\theta}\pd{(F_\theta\sen\theta)}{\theta}+\frac{1}{r\sen\theta}\pd{F_\varphi}{\varphi},\nonumber\\
  \nab\times\vc{F}&=\frac{1}{r\sen\theta}\Bigl(\pd{(F_\varphi\sen\theta)}{\theta}-\pd{F_\theta}{\varphi}\Bigr)\ur
    +\frac{1}{r}\Bigl(\frac{1}{\sen\theta}\pd{F_r}{\varphi}-\pd{(rF_\varphi)}{r}\Bigr)\uth
    +\frac{1}{r}\Bigl(\pd{(rF_\theta)}{r}-\pd{F_r}{\theta}\Bigr)\uph,\label{eq:ops-esf}\\
  \Lap f&=\frac{1}{r^{2}}\pd{}{r}\Bigl(r^{2}\pd{f}{r}\Bigr)+\frac{1}{r^{2}\sen\theta}\pd{}{\theta}\Bigl(\sen\theta\pd{f}{\theta}\Bigr)+\frac{1}{r^{2}\sen^{2}\theta}\pdd{f}{\varphi}.\nonumber
\end{align}

\begin{ejemplo}[Campos con simetría esférica]\label{ej:simetria-esferica}
Si $\vc{A}=A(r)\,\ur$ solo tiene componente radial y depende solo de $r$, las
\cref{eq:ops-esf} dan
\[
  \nab\cdot\vc{A}=\frac{1}{r^{2}}\frac{\dif}{\dif r}\bigl(r^{2}A\bigr),
  \qquad\nab\times\vc{A}=\vc{0},
\]
y para un campo escalar $f(r)$, $\nab f=f'(r)\,\ur$ y
$\Lap f=\frac{1}{r^{2}}\frac{\dif}{\dif r}\bigl(r^{2}f'\bigr)$. En particular,
$A=r^{n}$ da $\nab\cdot\vc{A}=(n+2)\,r^{n-1}$ y, para $A=1/r^{2}$ (campo de una
carga puntual), la divergencia se anula \emph{fuera del origen}.
\end{ejemplo}

% -------------------------------------------------------------------------
\subsection{Derivadas de los versores y jacobianos}\label{sec:versores-jacobianos}
% -------------------------------------------------------------------------
En cilíndricas y esféricas los versores \emph{dependen del punto}, y derivar un
campo exige derivar también la base. Se obtienen directamente de
las \cref{eq:base-cil,eq:base-esf}.

\begin{lema}[Derivadas de los versores]\label{lema:derivadas-versores}
Las únicas derivadas no nulas son, en cilíndricas,
\[
  \pd{\urho}{\varphi}=\uph,\qquad\pd{\uph}{\varphi}=-\urho,
\]
y en esféricas,
\[
  \pd{\ur}{\theta}=\uth,\quad\pd{\ur}{\varphi}=\sen\theta\,\uph,\quad
  \pd{\uth}{\theta}=-\ur,\quad\pd{\uth}{\varphi}=\cos\theta\,\uph,\quad
  \pd{\uph}{\varphi}=-\sen\theta\,\ur-\cos\theta\,\uth .
\]
\end{lema}
\begin{proof}
Se derivan los desarrollos cartesianos \eqref{eq:base-cil} y \eqref{eq:base-esf} (con
$\ux,\uy,\uz$ constantes) y se reexpresa el resultado en la base local. Por
ejemplo, $\partial_\varphi\urho=-\sen\varphi\,\ux+\cos\varphi\,\uy=\uph$ y
$\partial_\theta\ur=\cos\theta\cos\varphi\,\ux+\cos\theta\sen\varphi\,\uy-\sen\theta\,\uz=\uth$.
\end{proof}

Estas relaciones explican, por ejemplo, por qué el laplaciano vectorial no es
el laplaciano de cada componente cuando se usan curvilíneas
(\cref{def:laplaciano}), y por qué aparecen los factores $\rho$ y $r\sen\theta$
en los operadores.

\begin{proposicion}[Jacobianos]\label{prop:jacobianos}
Para el cambio de variable a coordenadas cilíndricas y esféricas,
\[
  \Bigl|\frac{\partial(x,y,z)}{\partial(\rho,\varphi,z)}\Bigr|=\rho,
  \qquad
  \Bigl|\frac{\partial(x,y,z)}{\partial(r,\theta,\varphi)}\Bigr|=r^{2}\sen\theta,
\]
de modo que $\iiint_Vf\,\dif x\,\dif y\,\dif z$ se transforma, respectivamente, en
\[
  \iiint f\,\rho\,\dif\rho\,\dif\varphi\,\dif z
  \qquad\text{y en}\qquad
  \iiint f\,r^{2}\sen\theta\,\dif r\,\dif\theta\,\dif\varphi .
\]
\end{proposicion}
\begin{proof}
Es el caso particular $|\det|=h_1h_2h_3$ de la \cref{prop:elementos} (el
valor absoluto elimina el signo, y en esféricas $\sen\theta\ge0$ porque
$\theta\in[0,\pi]$). El determinante explícito en esféricas se calcula, paso a
paso, en el \cref{ejer:jacobiano}.
\end{proof}

% =========================================================================
\section{Funciones de la posición relativa, delta de Dirac y desarrollo multipolar}\label{sec:posicion}
% =========================================================================
En electromagnetismo los campos suelen producirlos \emph{fuentes} (cargas,
corrientes) situadas en puntos $\vc{r}\,'$, y se evalúan en otros puntos
$\vc{r}$ donde queremos conocer el campo. Toda la dependencia aparece a través
del vector que une ambos puntos~\cite{FORMULARIO,GRIFFITHS}.

% -------------------------------------------------------------------------
\subsection{Punto fuente, punto campo y vector relativo}\label{sec:relativa}
% -------------------------------------------------------------------------
\begin{definicion}[Posición relativa]\label{def:relativa}
El \emph{punto fuente} es $\vc{r}\,'=(x',y',z')$ y el \emph{punto campo} es
$\vc{r}=(x,y,z)$. El vector que va de la fuente al punto campo, su módulo y su
versor son
\begin{align}
  \vc{R}&=\vc{r}-\vc{r}\,'=(x-x')\,\ux+(y-y')\,\uy+(z-z')\,\uz=R_x\,\ux+R_y\,\uy+R_z\,\uz,\label{eq:R-vec}\\
  R&=\norm{\vc{R}}=\bigl[(x-x')^{2}+(y-y')^{2}+(z-z')^{2}\bigr]^{1/2},\qquad
  \uR=\frac{\vc{R}}{R}\quad\Bigl(\vc{R}=R\,\uR\Bigr).\label{eq:R-mod}
\end{align}
\end{definicion}

% Figura: punto fuente, punto campo y vector que los une.
\begin{figure}[ht]
\centering
\tdplotsetmaincoords{70}{118}
\begin{tikzpicture}[tdplot_main_coords, scale=1.55]
  \def\sx{1.5}\def\sy{-0.6}\def\sz{2.0}   % punto fuente r'
  \def\fx{1.4}\def\fy{2.4}\def\fz{0.4}   % punto campo r
  \draw[eje] (0,0,0) -- (3.0,0,0) node[below left] {\small $x$};
  \draw[eje] (0,0,0) -- (0,3.1,0) node[right] {\small $y$};
  \draw[eje] (0,0,0) -- (0,0,3.1) node[above] {\small $z$};
  % vectores de posición
  \draw[->, line width=1.4pt, azulacero] (0,0,0) -- (\sx,\sy,\sz);
  \draw[->, line width=1.4pt, azulprof] (0,0,0) -- (\fx,\fy,\fz);
  % vector relativo R = r - r'
  \draw[->, line width=1.8pt, rojoteo] (\sx,\sy,\sz) -- (\fx,\fy,\fz);
  % versor u_R en el punto campo
  \pgfmathsetmacro{\Rx}{\fx-\sx}\pgfmathsetmacro{\Ry}{\fy-\sy}\pgfmathsetmacro{\Rz}{\fz-\sz}
  \pgfmathsetmacro{\Rn}{sqrt(\Rx*\Rx+\Ry*\Ry+\Rz*\Rz)}
  \draw[->, line width=1.6pt, verdeteo] (\fx,\fy,\fz) -- ++({0.95*\Rx/\Rn},{0.95*\Ry/\Rn},{0.95*\Rz/\Rn});
  \fill[azulacero] (\sx,\sy,\sz) circle (1.8pt);
  \fill[azulprof] (\fx,\fy,\fz) circle (1.8pt);
  \fill (0,0,0) circle (1.2pt);
  % etiquetas
  \node[azulacero, font=\normalsize] at (0.35,-0.55,1.0) {$\vc{r}\,'$};
  \node[azulprof, font=\normalsize] at (0.55,1.35,-0.15) {$\vc{r}$};
  \node[rojoteo, font=\normalsize, fill=white, fill opacity=0.8, text opacity=1, inner sep=1pt]
     at (1.55,1.15,1.5) {$\vc{R}=\vc{r}-\vc{r}\,'$};
  \node[verdeteo, font=\normalsize] at ({\fx+1.3*\Rx/\Rn},{\fy+1.3*\Ry/\Rn},{\fz+1.3*\Rz/\Rn}) {$\uR$};
  \node[font=\small, text=black!80, anchor=east] at (\sx+0.05,\sy-0.05,\sz+0.1) {fuente};
  \node[font=\small, text=black!80, anchor=north] at (\fx,\fy,\fz-0.12) {campo};
\end{tikzpicture}
\caption{Funciones de la posición relativa. El vector $\vc{r}\,'$ (azul claro)
localiza el \emph{punto fuente}, $\vc{r}$ (azul oscuro) el \emph{punto campo}, y
$\vc{R}=\vc{r}-\vc{r}\,'$ (rojo) apunta de la fuente al punto donde se evalúa
el campo; $\uR=\vc{R}/R$ (verde) es su versor.}
\label{fig:posicion-relativa}
\end{figure}


Interesan campos escalares $f(\vc{R})$ y vectoriales $\vc{A}(\vc{R})$ que dependen
de $\vc{r}$ y $\vc{r}\,'$ \emph{solo} a través de $\vc{R}$. Hay ahora dos
operadores nabla, que actúan sobre variables distintas:
\begin{equation}\label{eq:nabla-nabla-prima}
  \nab=\pd{}{x}\,\ux+\pd{}{y}\,\uy+\pd{}{z}\,\uz
  \qquad\text{y}\qquad
  \nab'=\pd{}{x'}\,\ux+\pd{}{y'}\,\uy+\pd{}{z'}\,\uz ,
\end{equation}
las derivadas respecto de las coordenadas del punto campo y del punto fuente,
respectivamente. Se diferencian en el signo, como se aprecia ya en la primera
derivada parcial:
\begin{equation}\label{eq:dR}
  \pd{R}{x}=\frac{1}{2}\,\frac{2(x-x')}{R}=\frac{R_x}{R},
  \qquad\text{mientras que}\qquad
  \pd{R}{x'}=-\frac{R_x}{R}.
\end{equation}

\begin{lema}[$\nab$ y $\nab'$ difieren en un signo]\label{lema:nabla-prima}
Si $g$ depende de $\vc{r}$ y $\vc{r}\,'$ solo a través de $\vc{R}=\vc{r}-\vc{r}\,'$,
entonces $\nab g=-\nab'g$ (y análogamente para la divergencia y el rotacional de
campos vectoriales de $\vc{R}$).
\end{lema}
\begin{proof}
Como $g$ es función de $x-x'$, la regla de la cadena da
$\partial g/\partial x'=-\partial g/\partial x$, y lo mismo para $y$ y $z$.
\end{proof}

% -------------------------------------------------------------------------
\subsection{Derivadas de potencias de $R$}\label{sec:derivadas-R}
% -------------------------------------------------------------------------
\begin{proposicion}[Gradiente de $R^{n}$]\label{prop:grad-Rn}
Para cualquier $n$ y $R\neq0$,
\begin{equation}\label{eq:grad-Rn}
  \nab(R^{n})=-\nab'(R^{n})=n\,R^{n-1}\,\uR=n\,R^{n-2}\,\vc{R}.
\end{equation}
\end{proposicion}
\begin{proof}
Por la regla de la cadena y la \cref{eq:dR}, cada componente es
\[
  \pd{R^{n}}{x}=\pd{R^{n}}{R}\,\pd{R}{x}=n\,R^{n-1}\,\frac{R_x}{R}=n\,R^{n-2}\,R_x ,
\]
y análogamente para $y,z$. Agrupando, $\nab(R^{n})=nR^{n-2}\bigl(R_x\,\ux+R_y\,\uy+R_z\,\uz\bigr)=nR^{n-2}\vc{R}=nR^{n-1}\uR$.
La igualdad con $-\nab'$ es el \cref{lema:nabla-prima}.
\end{proof}

\begin{proposicion}[Divergencia y rotacional de $R^{n}\uR$]\label{prop:div-rot-Rn}
Para $R\neq0$ y cualquier $n$ real (en particular entero),
\begin{equation}\label{eq:div-rot-Rn}
  \nab\cdot\bigl(R^{n}\,\uR\bigr)=-\nab'\cdot\bigl(R^{n}\,\uR\bigr)=(n+2)\,R^{n-1},
  \qquad
  \nab\times\bigl(R^{n}\,\uR\bigr)=-\nab'\times\bigl(R^{n}\,\uR\bigr)=\vc{0}.
\end{equation}
En particular, para $n=-2$: $\nab\cdot(\uR/R^{2})=0$ si $R\neq0$.
\end{proposicion}
\begin{proof}
Escribimos $R^{n}\uR=R^{n-1}\vc{R}$. Como $\nab\cdot\vc{R}=\partial_x(x-x')+\partial_y(y-y')+\partial_z(z-z')=3$ y
$\nab\times\vc{R}=\vc{0}$ (cada componente de $\vc{R}$ solo depende de su propia
coordenada), las reglas del producto \eqref{eq:prod-c} y \eqref{eq:prod-e} y la
\cref{prop:grad-Rn} dan
\begin{align*}
  \nab\cdot\bigl(R^{n-1}\vc{R}\bigr)&=R^{n-1}\,\nab\cdot\vc{R}+\vc{R}\cdot\nab R^{n-1}
    =3R^{n-1}+(n-1)R^{n-2}\,\vc{R}\cdot\uR=(n+2)\,R^{n-1},\\
  \nab\times\bigl(R^{n-1}\vc{R}\bigr)&=\nab R^{n-1}\times\vc{R}+R^{n-1}\,\nab\times\vc{R}
    =(n-1)R^{n-2}\,\uR\times\vc{R}+\vc{0}=\vc{0},
\end{align*}
pues $\uR\parallel\vc{R}$. Las igualdades con $\nab'$ son el \cref{lema:nabla-prima}.
En el \cref{ejer:div-Rn} se repite el cálculo de la divergencia componente a
componente.
\end{proof}

\begin{observacion}[Comprobaciones y el caso $n=-2$]\label{obs:casos-n}
Para $n=1$, $R\,\uR=\vc{R}$ y $\nab\cdot\vc{R}=3$, como en el \cref{ej:div}; para
$n=0$, $\nab\cdot\uR=2/R$; y para $n=-2$ el factor $(n+2)$ se anula. Pero
atención: el campo $\uR/R^{2}$ \emph{no está definido} en $R=0$ y su divergencia,
nula en todo punto $R\neq0$, tiene una contribución \emph{singular} concentrada
en el origen que la \cref{eq:div-rot-Rn} no recoge. Esa contribución es una delta de
Dirac (\cref{teo:delta}).
\end{observacion}

% -------------------------------------------------------------------------
\subsection{La delta de Dirac}\label{sec:delta}
% -------------------------------------------------------------------------
\begin{definicion}[Delta de Dirac]\label{def:delta}
La \emph{delta de Dirac} unidimensional $\delta(x)$ es el objeto que cumple
\begin{equation}\label{eq:delta1d}
  \delta(x)=\begin{cases}0,&x\neq0\\ \infty,&x=0\end{cases},
  \qquad
  \int_{-\infty}^{\infty}\delta(x)\,\dif x=1,
  \qquad
  \int_{-\infty}^{\infty}f(x)\,\delta(x-a)\,\dif x=f(a)
\end{equation}
para toda función continua $f$ (propiedad de selección o «muestreo»). En tres
dimensiones, $\delta^{3}(\vc{r})=\delta(x)\,\delta(y)\,\delta(z)$ cumple
\begin{equation}\label{eq:delta3d}
  \delta^{3}(\vc{r})=\begin{cases}0,&\vc{r}\neq\vc{0}\\ \infty,&\vc{r}=\vc{0}\end{cases},
  \qquad
  \int_{V_\infty}\delta^{3}(\vc{r})\,\dV=1,
  \qquad
  \int_{V_\infty}f(\vc{r})\,\delta^{3}(\vc{r}-\vc{a})\,\dV=f(\vc{a}),
\end{equation}
donde $V_\infty$ es todo el espacio (o cualquier región que contenga a
$\vc{a}$; si no lo contiene, la integral es nula).
\end{definicion}

\begin{observacion}[Qué es en realidad]\label{obs:delta}
La delta no es una función en el sentido usual (ninguna función es nula salvo
en un punto y tiene integral~$1$), sino una \emph{distribución}: un objeto que
solo tiene sentido dentro de una integral, donde actúa como un «detector» de
valores puntuales. Intuitivamente es el límite de picos cada vez más estrechos y
altos de área unidad, por ejemplo gaussianas
$\delta_\sigma(x)=\tfrac{1}{\sigma\sqrt{2\pi}}\,\eul^{-x^{2}/2\sigma^{2}}$ con
$\sigma\to0$ (\cref{fig:delta}). Representa matemáticamente una carga (o masa)
puntual: densidad nula en todas partes menos en un punto, y carga total finita.
\end{observacion}

% Figura: aproximaciones de la delta de Dirac por gaussianas de área unidad.
\begin{figure}[ht]
\centering
\begin{tikzpicture}
  \begin{axis}[tema, width=0.62\textwidth, height=5.6cm, xlabel={$x$}, ylabel={$\delta_\sigma(x)$},
      xmin=-2.2, xmax=2.2, ymin=0, ymax=3.6, samples=200, domain=-2.2:2.2,
      legend pos=north east, legend cell align=left, grid=major]
    \addplot[azulclaro!80!azulacero, line width=1.1pt] {1/(1*sqrt(2*pi))*exp(-x^2/(2*1^2))};
    \addplot[azulacero, line width=1.1pt] {1/(0.5*sqrt(2*pi))*exp(-x^2/(2*0.5^2))};
    \addplot[azulprof, line width=1.1pt] {1/(0.25*sqrt(2*pi))*exp(-x^2/(2*0.25^2))};
    \addplot[rojoteo, line width=1.1pt, samples=400] {1/(0.12*sqrt(2*pi))*exp(-x^2/(2*0.12^2))};
    \legend{$\sigma=1$, $\sigma=0{,}5$, $\sigma=0{,}25$, $\sigma=0{,}12$}
  \end{axis}
\end{tikzpicture}
\caption{La delta de Dirac como límite de picos cada vez más estrechos y altos
con área unidad: $\delta_\sigma(x)=\frac{1}{\sigma\sqrt{2\pi}}\,\eul^{-x^{2}/2\sigma^{2}}$.
Cuando $\sigma\to0$, $\int\delta_\sigma\,\dif x=1$ se conserva y el pico se
concentra en $x=0$.}
\label{fig:delta}
\end{figure}


\begin{teorema}[Divergencia del campo coulombiano y laplaciano de $1/R$]\label{teo:delta}
En el sentido de distribuciones,
\begin{equation}\label{eq:delta-div}
  \nab\cdot\Bigl(\frac{\uR}{R^{2}}\Bigr)=-\nab'\cdot\Bigl(\frac{\uR}{R^{2}}\Bigr)=4\pi\,\delta^{3}(\vc{R}),
  \qquad
  \Lap\Bigl(\frac{1}{R}\Bigr)=\Lapp\Bigl(\frac{1}{R}\Bigr)=-4\pi\,\delta^{3}(\vc{R}).
\end{equation}
\end{teorema}
\begin{proof}
\emph{Primera igualdad.} Por la \cref{prop:div-rot-Rn}, $\nab\cdot(\uR/R^{2})=0$ para
$R\neq0$: la distribución, si no es nula, está concentrada en $R=0$. Calculemos
su integral sobre una región $V$ que contenga a la fuente. Quitamos una bola
pequeña $B_\varepsilon$ de radio $\varepsilon$ alrededor de $\vc{r}\,'$; en
$V\setminus B_\varepsilon$ el campo es de clase $C^{1}$ y de divergencia nula, así
que por el teorema de Gauss (\cref{teo:gauss})
\[
  0=\iiint_{V\setminus B_\varepsilon}\nab\cdot\frac{\uR}{R^{2}}\,\dV
   =\oiint_{\partial V}\frac{\uR}{R^{2}}\cdot\dS-\oiint_{\partial B_\varepsilon}\frac{\uR}{R^{2}}\cdot\dS ,
\]
donde la normal exterior de la bola es $\uR$ (con el signo $-$ porque, para
$V\setminus B_\varepsilon$, la normal es la interior a la bola). Sobre
$\partial B_\varepsilon$, $\uR\cdot\uR/\varepsilon^{2}=1/\varepsilon^{2}$ es constante y
el flujo vale $\tfrac{1}{\varepsilon^{2}}\cdot4\pi\varepsilon^{2}=4\pi$. Por tanto
el flujo del campo por \emph{cualquier} superficie cerrada que rodee la fuente
es $4\pi$, esto es, la integral de su divergencia sobre cualquier región que
contenga a $\vc{r}\,'$ vale $4\pi$, independientemente del tamaño y la forma: es
exactamente la propiedad $\int\delta^{3}\,\dV=1$ multiplicada por $4\pi$.
Si $\varphi$ es una función continua cualquiera, el mismo argumento con
$\varepsilon\to0$ da $\int\varphi\,\nab\cdot(\uR/R^{2})\,\dV=4\pi\varphi(\vc{r}\,')$, que es
\eqref{eq:delta3d}.

\emph{Segunda igualdad.} Por la \cref{prop:grad-Rn} con $n=-1$,
$\nab(1/R)=-\uR/R^{2}$, de modo que $\Lap(1/R)=\nab\cdot\nab(1/R)=-\nab\cdot(\uR/R^{2})=-4\pi\,\delta^{3}(\vc{R})$.
Finalmente $\Lapp(1/R)=\nab'\cdot\nab'(1/R)=(-\nab)\cdot(-\nab)(1/R)=\Lap(1/R)$.
\end{proof}

En electrostática esta identidad es la ley de Gauss en forma local: el campo de
una carga puntual es proporcional a $\uR/R^{2}$ y su divergencia es una delta
centrada en la carga.

% -------------------------------------------------------------------------
\subsection{Polinomios de Legendre y desarrollo de $1/|\vc{r}-\vc{r}\,'|$}\label{sec:legendre}
% -------------------------------------------------------------------------
El desarrollo del inverso de la distancia en serie de potencias de
$r'/r$ es la base de la expansión multipolar. Sus coeficientes son los
polinomios de Legendre.

\begin{definicion}[Polinomios de Legendre]\label{def:legendre}
Los polinomios de Legendre $P_n(x)$, $x\in[-1,1]$, vienen dados por la fórmula de
Rodrigues $P_n(x)=\frac{1}{2^{n}n!}\frac{\dif^{n}}{\dif x^{n}}(x^{2}-1)^{n}$. Los
primeros son
\[
  P_0(x)=1,\qquad P_1(x)=x,\qquad P_2(x)=\tfrac{1}{2}(3x^{2}-1),\qquad P_3(x)=\tfrac{1}{2}(5x^{3}-3x),
\]
y cumplen $P_n(1)=1$ y la recurrencia
$(n+1)P_{n+1}=(2n+1)\,x\,P_n-n\,P_{n-1}$.
\end{definicion}

\begin{teorema}[Desarrollo en polinomios de Legendre]\label{teo:legendre}
Si $r'<r$ y $\theta'$ es el ángulo que forman $\vc{r}$ y $\vc{r}\,'$,
\begin{equation}\label{eq:legendre}
  \frac{1}{\abs{\vc{r}-\vc{r}\,'}}=\frac{1}{r}\sum_{n=0}^{\infty}\Bigl(\frac{r'}{r}\Bigr)^{n}P_n(\cos\theta').
\end{equation}
\end{teorema}
\begin{proof}
Por el teorema del coseno,
$\abs{\vc{r}-\vc{r}\,'}^{2}=r^{2}+r'^{2}-2rr'\cos\theta'=r^{2}\bigl(1-2xt+t^{2}\bigr)$
con $t=r'/r<1$ y $x=\cos\theta'$. Por tanto
$\frac{1}{\abs{\vc{r}-\vc{r}\,'}}=\frac{1}{r}(1+u)^{-1/2}$ con $u=t^{2}-2xt$ pequeño,
y el desarrollo binomial $(1+u)^{-1/2}=1-\tfrac{u}{2}+\tfrac{3u^{2}}{8}-\tfrac{5u^{3}}{16}+\cdots$
ordenado por potencias de $t$ da
\[
  1+x\,t+\tfrac{1}{2}(3x^{2}-1)\,t^{2}+\tfrac{1}{2}(5x^{3}-3x)\,t^{3}+\cdots
  =P_0+P_1\,t+P_2\,t^{2}+P_3\,t^{3}+\cdots ,
\]
es decir, la función generatriz $(1-2xt+t^{2})^{-1/2}=\sum_nP_n(x)\,t^{n}$ de los
polinomios de Legendre~\cite{ARFKEN}.
\end{proof}

% Figura: polinomios de Legendre P0..P3.
\begin{figure}[ht]
\centering
\begin{tikzpicture}
  \begin{axis}[tema, width=0.62\textwidth, height=6.2cm, xlabel={$x=\cos\theta'$}, ylabel={$P_n(x)$},
      xmin=-1, xmax=1, ymin=-1.05, ymax=1.15, samples=100, domain=-1:1,
      xtick={-1,-0.5,0,0.5,1}, legend pos=south east, legend cell align=left]
    \addplot[grispizarra, line width=1.1pt] {1};
    \addplot[azulclaro!70!azulacero, line width=1.1pt] {x};
    \addplot[azulacero, line width=1.1pt] {0.5*(3*x^2-1)};
    \addplot[azulprof, line width=1.1pt] {0.5*(5*x^3-3*x)};
    \legend{$P_0$, $P_1$, $P_2$, $P_3$}
  \end{axis}
\end{tikzpicture}
\caption{Primeros polinomios de Legendre en $[-1,1]$. Todos valen $1$ en $x=1$
(dirección de $\vc{r}\,'$ paralela a $\vc{r}$).}
\label{fig:legendre}
\end{figure}


Cada término tiene un significado físico: $n=0$ es el campo de un monopolo
($\propto1/r$); $n=1$, un dipolo ($\propto r'\cos\theta'/r^{2}$), etc. Si $r'>r$
basta intercambiar los papeles de $r$ y $r'$.

% =========================================================================
\section{Seminarios prácticos}\label{sec:seminarios}
% =========================================================================
Esta sección reúne los ejercicios del tema resueltos con todo detalle. Cada
ejercicio sigue la misma estructura: un \emph{enunciado}, la \emph{idea} que
guía la resolución (qué estamos buscando y por qué el método funciona), la
\emph{resolución} paso a paso y, cuando merece la pena, un breve comentario con
la \emph{interpretación} geométrica o física del resultado. Los seminarios
\ref{sem:pendiente}--\ref{sem:levi} trabajan con operadores diferenciales;
los demás verifican, con campos concretos, los tres teoremas integrales de la
\cref{sec:integral}.

% =========================================================================
\subsection{Seminario 1. Campos escalares: la pendiente máxima de una campana}\label{sem:pendiente}
% =========================================================================
\begin{ejercicio}[Pendiente máxima]\label{ejer:pendiente}
Para el campo escalar $f(\vc{r})=\eul^{-(x^{2}+y^{2})}$, calcular a qué distancia
del origen la pendiente es máxima.
\end{ejercicio}

\begin{resolucion}
\textbf{Idea.} Imaginamos $f:\R^{2}\to\R$ como una superficie $z=f(x,y)$: una
campana centrada en el origen. En un punto $P=(x,y)$ podemos movernos en
cualquier dirección del plano y la derivada direccional nos dice cuál es la
pendiente de la superficie en esa dirección. Por la \cref{prop:maximo}, la
mayor de todas esas pendientes en $P$ es exactamente $\norm{\nab f(P)}$, y se
alcanza al ir en la dirección del gradiente. La pregunta se transforma
entonces en: \emph{¿dónde es máximo $\norm{\nab f}$?}

\textbf{Paso 1: el gradiente.} Por la regla de la cadena
$[g(f(x))]'=g'(f(x))\,f'(x)$,
\[
  \pd{f}{x}=-2x\,\eul^{-(x^{2}+y^{2})},\qquad
  \pd{f}{y}=-2y\,\eul^{-(x^{2}+y^{2})},
  \quad\Longrightarrow\quad
  \nab f=-2\,\eul^{-(x^{2}+y^{2})}\,(x\,\ux+y\,\uy),
\]
cuyo módulo es $\norm{\nab f}=2\,\eul^{-(x^{2}+y^{2})}\sqrt{x^{2}+y^{2}}$.

\textbf{Paso 2: aprovechar la simetría.} El módulo depende solo de la distancia
al origen. Pasando a polares, $x=r\cos\varphi$, $y=r\sen\varphi$, con
$r=\sqrt{x^{2}+y^{2}}$ ($x^{2}+y^{2}=r^{2}$), queda
$\norm{\nab f(r,\varphi)}=\norm{\nab f(r)}=2r\,\eul^{-r^{2}}$. El problema
bidimensional se ha convertido en uno unidimensional: maximizar
\[
  g(r)=2r\,\eul^{-r^{2}},\qquad r\ge0 .
\]

\textbf{Paso 3: puntos críticos.} Derivamos:
$g'(r)=2\eul^{-r^{2}}+2r(-2r)\eul^{-r^{2}}=2\,\eul^{-r^{2}}(1-2r^{2})$. En un punto
crítico $g'(r)=0$; como $\eul^{-r^{2}}>0$ para todo $r$, necesariamente
$1-2r^{2}=0$, es decir,
\[
  \boxed{\,r=\frac{1}{\sqrt{2}}\approx0{,}707\,}
\]
(en unidades arbitrarias del problema).

\textbf{Paso 4: comprobar que es un máximo.} La segunda derivada es
$g''(r)=4r\,\eul^{-r^{2}}(2r^{2}-3)$ (se obtiene derivando
$2\eul^{-r^{2}}(1-2r^{2})$ por el producto). En $r=1/\sqrt{2}$,
$g''(1/\sqrt{2})=4\tfrac{1}{\sqrt{2}}\,\eul^{-1/2}(1-3)<0$, así que efectivamente
es un máximo, de valor $g(1/\sqrt{2})=\sqrt{2}\,\eul^{-1/2}\approx0{,}858$.
\end{resolucion}

\begin{observacion}[Interpretación]
El resultado es razonable. En el origen $\nab f(\vc{0})=\vc{0}$: la cima de la
campana es plana, la pendiente es nula. Muy lejos, $\eul^{-r^{2}}\to0$ y la
superficie vuelve a aplanarse. Por tanto, tiene que haber una zona intermedia
donde la pendiente sea máxima: crece desde $0$, alcanza su máximo en
$r=1/\sqrt{2}$ y vuelve a decrecer hasta anularse (\cref{fig:sem-pendiente}).
Es, además, el \emph{punto de inflexión} del perfil $\eul^{-r^{2}}$: allí
$f''(r)=(4r^{2}-2)\eul^{-r^{2}}=0$, como corresponde a un punto donde la
pendiente $f'$ es extremal.
\end{observacion}

% Figura: pendiente máxima de la campana gaussiana
\begin{figure}[ht]
\centering
\begin{minipage}[t]{0.52\textwidth}\centering
\begin{tikzpicture}[baseline=(current bounding box.north)]
  \begin{axis}[tema, width=\textwidth, height=5.8cm, xlabel={$r$}, xmin=0, xmax=3, ymin=0, ymax=1.12,
      samples=120, domain=0:3, xtick={0,0.5,0.707,1,1.5,2,3}, xticklabels={0,0{,}5,$\tfrac{1}{\sqrt2}$,1,1{,}5,2,3},
      legend style={at={(0.98,0.97)}, anchor=north east}]
    \addplot[azulacero, line width=1.1pt] {exp(-x^2)};
    \addplot[rojoteo, line width=1.1pt] {2*x*exp(-x^2)};
    \addplot[grispizarra, densely dashed, line width=0.6pt] coordinates {(0.7071,0) (0.7071,0.858)};
    \addplot[only marks, mark=*, mark size=2pt, rojoteo] coordinates {(0.7071,0.8578)};
    \node[font=\scriptsize, rojoteo, anchor=west] at (axis cs:0.78,0.98) {máximo: $\sqrt2\,\eul^{-1/2}\approx0{,}858$};
    \legend{$f(r)=\eul^{-r^{2}}$, $\norm{\nab f}=g(r)=2r\,\eul^{-r^{2}}$}
  \end{axis}
\end{tikzpicture}\\[2pt]{\small\scshape (a)}
\end{minipage}\hfill
\begin{minipage}[t]{0.44\textwidth}\centering
\begin{tikzpicture}[x=0.78cm, y=0.78cm, baseline=(current bounding box.north)]
  \mapaGauss{%
    \draw[white, dashed, line width=1.1pt] (0,0) circle (0.7071);
    \foreach \a in {0,45,...,315}{%
      \draw[-{Stealth[length=4pt]}, line width=0.9pt, white] ({0.7071*cos(\a)},{0.7071*sin(\a)}) -- ++({-0.62*cos(\a)},{-0.62*sin(\a)});}
    \node[white, font=\scriptsize, fill=cm0, fill opacity=0.6, text opacity=1, inner sep=1pt] at (1.45,1.45) {$r=\tfrac{1}{\sqrt2}$};}
\end{tikzpicture}\\[2pt]{\small\scshape (b)}
\end{minipage}
\caption{Pendiente máxima de $f=\eul^{-(x^{2}+y^{2})}$. (a)~Perfil radial $f(r)$ y
módulo de su gradiente $g(r)=\norm{\nab f}$, que crece desde cero en el origen,
alcanza su máximo en $r=1/\sqrt2$ y decae. (b)~Sobre el mapa de calor, la
circunferencia $r=1/\sqrt2$ (discontinua) es el lugar geométrico de máxima
pendiente; las flechas son $\nab f$ allí.}
\label{fig:sem-pendiente}
\end{figure}


% =========================================================================
\subsection{Seminario 2. Gradiente y divergencia: cartesianas frente a esféricas}\label{sem:inverso}
% =========================================================================
\begin{ejercicio}[El campo del inverso del cuadrado]\label{ejer:inverso}
Sea $f(\vc{r})=1/r=(x^{2}+y^{2}+z^{2})^{-1/2}$ (a)~Calcular $\nab f$ en
coordenadas cartesianas y en esféricas. (b)~Para el campo
$\vc{A}(\vc{r})=\ur/r^{2}$ calcular su divergencia en ambos sistemas.
\end{ejercicio}

\begin{resolucion}
\textbf{(a) Gradiente.} En cartesianas, $\partial_x(x^{2}+y^{2}+z^{2})^{-1/2}=-\tfrac12(x^{2}+y^{2}+z^{2})^{-3/2}\cdot2x=-x/r^{3}$, y lo mismo para $y,z$:
\[
  \nab f=\pd{f}{x}\,\ux+\pd{f}{y}\,\uy+\pd{f}{z}\,\uz
  =-\frac{1}{r^{3}}\bigl(x\,\ux+y\,\uy+z\,\uz\bigr)=-\frac{\vc{r}}{r^{3}}=-\frac{\ur}{r^{2}},
\]
donde hemos usado $\vc{r}=r\,\ur$. En esféricas, como $f=f(r)$ solo depende de
$r$, la fórmula del gradiente (\cref{eq:ops-esf}) conserva un único término:
$\nab f=\pd{f}{r}\,\ur=-\dfrac{1}{r^{2}}\,\ur$. Mismo resultado, pero obtenido
sin cálculos: ¡la simetría esférica hace que las esféricas sean el sistema
natural!

\textbf{(b) Divergencia.} El campo $\vc{A}=\ur/r^{2}=\vc{r}/r^{3}$ tiene
componentes cartesianas
$\vc{A}=\dfrac{x\,\ux+y\,\uy+z\,\uz}{(x^{2}+y^{2}+z^{2})^{3/2}}$ (es, salvo signo,
el campo $\nab f$ anterior). En cartesianas,
\[
  \nab\cdot\vc{A}=\pd{A_x}{x}+\pd{A_y}{y}+\pd{A_z}{z},
\]
y cada sumando se calcula con la regla del cociente:
\[
  \pd{A_x}{x}=\pd{}{x}\Bigl[\frac{x}{(x^{2}+y^{2}+z^{2})^{3/2}}\Bigr]
  =\frac{(x^{2}+y^{2}+z^{2})^{3/2}-x\cdot\frac{3}{2}(x^{2}+y^{2}+z^{2})^{1/2}\cdot2x}{(x^{2}+y^{2}+z^{2})^{3}}
  =\frac{1}{r^{3}}-\frac{3x^{2}}{r^{5}} .
\]
Sumando las tres,
\[
  \nab\cdot\vc{A}=\frac{3}{r^{3}}-\frac{3\,(x^{2}+y^{2}+z^{2})}{r^{5}}=\frac{3}{r^{3}}-\frac{3}{r^{3}}=0 .
\]
En esféricas, con $A_r=1/r^{2}$ y $A_\theta=A_\varphi=0$,
\[
  \nab\cdot\vc{A}=\frac{1}{r^{2}}\pd{}{r}\Bigl(r^{2}\cdot\frac{1}{r^{2}}\Bigr)=\frac{1}{r^{2}}\pd{(1)}{r}=0 .
\]
\end{resolucion}

\begin{observacion}[La singularidad del origen]
La divergencia es nula \emph{en todo punto excepto en el origen}, donde el
campo no está definido (diverge) y la divergencia tiene una contribución
singular: una delta de Dirac. El hecho de que, a pesar de ello, el flujo por
cualquier esfera centrada en el origen valga $4\pi$ (\cref{ej:flujo-r3}) y
\emph{no dependa de su radio} es la huella de esa fuente puntual, como muestra
la \cref{fig:sem-inverso} y se formaliza en el \cref{teo:delta}.
\end{observacion}

% Figura: campo del inverso del cuadrado y flujo constante
\begin{figure}[ht]
\centering
\begin{minipage}[t]{0.47\textwidth}\centering
\begin{tikzpicture}[x=0.84cm, y=0.84cm, baseline=(current bounding box.north)]
  \marcoD{-2.7}{2.7}{-2.7}{2.7}{-2,-1,0,1,2}{-2,-1,0,1,2}
  \ejesD{-2.7}{2.7}{-2.7}{2.7}
  \clipD{-2.7}{2.7}{-2.7}{2.7}{\flechas{\x/((\x*\x+\y*\y)^1.5)}{\y/((\x*\x+\y*\y)^1.5)}%
     {-2.4,-1.8,-1.2,-0.6,0.6,1.2,1.8,2.4}{-2.4,-1.8,-1.2,-0.6,0.6,1.2,1.8,2.4}{0.42}{1.4}}
  \fill[azulprof] (0,0) circle (2pt);
\end{tikzpicture}\\[2pt]{\small\scshape (a)}\ $\vc{A}=\ur/r^{2}$
\end{minipage}\hfill
\begin{minipage}[t]{0.50\textwidth}\centering
\begin{tikzpicture}[baseline=(current bounding box.north)]
  \begin{axis}[tema, width=\textwidth, height=5.4cm, xlabel={$r$}, xmin=0, xmax=3, ymin=0, ymax=22,
      samples=150, domain=0.26:3, legend style={at={(0.97,0.97)}, anchor=north east}]
    \addplot[azulacero, line width=1.1pt] {1/x^2};
    \addplot[rojoteo, line width=1.1pt, domain=0.26:3] {4*pi};
    \legend{$\norm{\vc{A}}=1/r^{2}$, $\Phi(r)=\oiint_{S_r}\vc{A}\cdot\dS=4\pi$}
  \end{axis}
\end{tikzpicture}\\[2pt]{\small\scshape (b)}
\end{minipage}
\caption{El campo $\vc{A}=\ur/r^{2}$. (a)~Flechas radiales cuyo módulo cae como
$1/r^{2}$ al alejarse de la fuente (el origen). (b)~El módulo decrece, pero el
flujo por una esfera de radio $r$, $\Phi=\norm{\vc{A}}\cdot4\pi r^{2}$, es constante
e igual a $4\pi$: el campo «se diluye» exactamente al ritmo en que crece la
superficie, de modo que no hay fuentes salvo la del origen.}
\label{fig:sem-inverso}
\end{figure}


% =========================================================================
\subsection{Seminario 3. Cálculo con índices: el tensor de Levi-Civita}\label{sem:levi}\label{sec:levi-civita}
% =========================================================================
Varias identidades vectoriales de la \cref{prop:productos} y de la
\cref{prop:fundamentales} son tediosas de comprobar componente a componente.
La notación con índices las convierte en cálculos cortos y mecánicos. Para ello
necesitamos una herramienta.

\begin{definicion}[Tensor de Levi-Civita]\label{def:levi-civita}
El \emph{tensor de Levi-Civita} (o símbolo de permutación) en tres dimensiones es
\[
  \varepsilon_{ijk}=\begin{cases}
    +1 & \text{si }(i,j,k)\text{ es una permutación par de }(1,2,3),\\
    -1 & \text{si }(i,j,k)\text{ es una permutación impar de }(1,2,3),\\
    \ \ 0 & \text{si dos índices cualesquiera son iguales.}
  \end{cases}
\]
Es decir, $\varepsilon_{123}=\varepsilon_{231}=\varepsilon_{312}=+1$ (permutaciones
cíclicas) y $\varepsilon_{132}=\varepsilon_{213}=\varepsilon_{321}=-1$.
\end{definicion}

% Figura: tensor de Levi-Civita como orientación cíclica
\begin{figure}[ht]
\centering
\begin{tikzpicture}[>=Stealth, x=1cm, y=1cm]
  \foreach \i/\ang in {1/90,2/-30,3/210}{%
    \node[circle, draw=azulprof, fill=azulclaro!50, minimum size=0.8cm, font=\large, text=azulprof] (n\i) at (\ang:1.5) {\i};}
  \draw[->, line width=1.4pt, azulprof] (n1) to[bend left=18] (n2);
  \draw[->, line width=1.4pt, azulprof] (n2) to[bend left=18] (n3);
  \draw[->, line width=1.4pt, azulprof] (n3) to[bend left=18] (n1);
  \node[font=\small, text=azulprof, align=center] at (0,0) {sentido\\ cíclico};
  \node[anchor=west, align=left, font=\small] at (2.6,0.55)
     {$\textcolor{azulprof}{\varepsilon_{123}=\varepsilon_{231}=\varepsilon_{312}=+1}$\quad(a favor de las flechas)};
  \node[anchor=west, align=left, font=\small] at (2.6,-0.15)
     {$\textcolor{rojoteo}{\varepsilon_{132}=\varepsilon_{213}=\varepsilon_{321}=-1}$\quad(en contra)};
  \node[anchor=west, align=left, font=\small] at (2.6,-0.85)
     {$\textcolor{grispizarra}{\varepsilon_{112}=\varepsilon_{233}=\varepsilon_{121}=\cdots=0}$\quad(índices repetidos)};
\end{tikzpicture}
\caption{El tensor de Levi-Civita codifica la orientación: es $+1$ si los
índices recorren $1,2,3$ en el sentido cíclico de las flechas, $-1$ si lo hacen
en sentido contrario y $0$ si algún índice se repite. Así, $\vc{A}\times\vc{B}$
con $(\vc{A}\times\vc{B})_i=\varepsilon_{ijk}A_jB_k$ reproduce
$\ux\times\uy=\uz$ y los demás productos de la base.}
\label{fig:levi}
\end{figure}


Con el convenio de sumación de Einstein (\cref{not:convenios}), el producto
vectorial, el rotacional, la divergencia y el gradiente se escriben
\begin{equation}\label{eq:indices}
  (\vc{A}\times\vc{B})_i=\varepsilon_{ijk}A_jB_k,\qquad
  (\nab\times\vc{A})_i=\varepsilon_{ijk}\partial_jA_k,\qquad
  \nab\cdot\vc{A}=\partial_iA_i,\qquad
  (\nab f)_i=\partial_if ,
\end{equation}
donde $\partial_i\equiv\partial/\partial x_i$ (derivada \emph{espacial}). En la
primera expresión, $c_i=\sum_{j,k=1}^{3}\varepsilon_{ijk}A_jB_k$ es la componente
$i$-ésima, y los índices repetidos $j,k$ se suman.

\begin{lema}[Propiedades del tensor de Levi-Civita]\label{lema:levi}
\begin{enumerate}[label=(\alph*)]
  \item \emph{Antisimetría:} intercambiar dos índices cualesquiera cambia el
        signo, por ejemplo $\varepsilon_{ijk}=-\varepsilon_{jik}=-\varepsilon_{ikj}$.
  \item \emph{Invariancia cíclica:} $\varepsilon_{ijk}=\varepsilon_{jki}=\varepsilon_{kij}$.
  \item \emph{Constancia:} los valores de $\varepsilon_{ijk}$ son números fijos, así
        que $\partial_l\varepsilon_{ijk}=0$ y puede sacarse de cualquier derivada.
\end{enumerate}
\end{lema}
\begin{proof}
(a) Intercambiar dos índices convierte una permutación par en impar y viceversa
(una transposición cambia la paridad), y deja en cero las de índices repetidos.
(b) Una permutación cíclica de tres elementos es par, luego no cambia el signo.
(c) Es constante por definición.
\end{proof}

\begin{lema}[Identidad $\varepsilon$--$\delta$]\label{lema:eps-delta}
Con $\delta_{ij}$ la delta de Kronecker,
\[
  \varepsilon_{ijk}\,\varepsilon_{ilm}=\delta_{jl}\,\delta_{km}-\delta_{jm}\,\delta_{kl}.
\]
\end{lema}
\begin{proof}
Para $j,k$ fijos, el producto $\varepsilon_{ijk}\varepsilon_{ilm}$ (sumado en $i$) solo
es no nulo si $j\neq k$ y si $\{l,m\}=\{j,k\}$, pues entonces $i$ es el tercer
valor distinto. Si $(l,m)=(j,k)$ ambos factores son iguales y el producto es
$(\pm1)^{2}=+1$; si $(l,m)=(k,j)$ difieren en un intercambio y el producto es
$-1$. Eso es exactamente lo que valen $\delta_{jl}\delta_{km}-\delta_{jm}\delta_{kl}$.
\end{proof}

\begin{lema}[Contracción de un tensor simétrico con otro antisimétrico]\label{lema:sim-anti}
Sean $S_{ij}$ simétrico ($S_{ij}=S_{ji}$) y $T_{ij}$ antisimétrico
($T_{ij}=-T_{ji}$). Entonces $C\equiv S_{ij}\,T_{ij}=\sum_{i,j}S_{ij}T_{ij}=0$.
\end{lema}
\begin{proof}
Los índices sumados son «mudos»: podemos renombrarlos libremente sin alterar la
suma ($i\leftrightarrow j$). Entonces
$C=S_{ji}T_{ji}$. Por simetría y antisimetría, $S_{ji}=S_{ij}$ y $T_{ji}=-T_{ij}$, luego
$C=-S_{ij}T_{ij}=-C$, y por tanto $C=0$.
\end{proof}

Además, para campos de clase $C^{2}$ las derivadas parciales conmutan
(teorema de Schwarz), de modo que el operador $\partial_i\partial_j$ es
\emph{simétrico} en $(i,j)$.

\begin{ejercicio}[Divergencia de un producto vectorial]\label{ejer:div-prod-vect}
Demostrar que $\nab\cdot(\vc{A}\times\vc{B})=\vc{B}\cdot(\nab\times\vc{A})-\vc{A}\cdot(\nab\times\vc{B})$.
\end{ejercicio}
\begin{resolucion}
\textbf{Idea.} Escribimos todo con índices, aplicamos la regla del producto a la
derivada (el símbolo $\varepsilon$ es constante) y reconocemos, en cada término,
la componente de un rotacional.

Por la \cref{eq:indices}, la divergencia es $\partial_i(\cdot)_i$ aplicada a
$(\vc{A}\times\vc{B})_i=\varepsilon_{ijk}A_jB_k$:
\[
  \nab\cdot(\vc{A}\times\vc{B})=\partial_i\bigl(\varepsilon_{ijk}A_jB_k\bigr)
  =\varepsilon_{ijk}\bigl[(\partial_iA_j)\,B_k+A_j\,(\partial_iB_k)\bigr].
\]
\emph{Primer sumando.} Por la invariancia cíclica, $\varepsilon_{ijk}=\varepsilon_{kij}$, así que
$\varepsilon_{ijk}\,\partial_iA_j=\varepsilon_{kij}\,\partial_iA_j=(\nab\times\vc{A})_k$
(la definición del rotacional, con $k$ como índice libre). Por tanto
$\varepsilon_{ijk}(\partial_iA_j)B_k=(\nab\times\vc{A})_kB_k=\vc{B}\cdot(\nab\times\vc{A})$.

\emph{Segundo sumando.} Usamos antisimetría para reordenar:
$\varepsilon_{ijk}=-\varepsilon_{jik}$, luego
$\varepsilon_{ijk}\,\partial_iB_k=-\varepsilon_{jik}\,\partial_iB_k=-(\nab\times\vc{B})_j$, y
$\varepsilon_{ijk}A_j(\partial_iB_k)=-A_j(\nab\times\vc{B})_j=-\vc{A}\cdot(\nab\times\vc{B})$.

Sumando ambos resultados se obtiene la identidad pedida.
\end{resolucion}

\begin{ejercicio}[Rotacional de un escalar por un vector]\label{ejer:rot-fA}
Demostrar que $\nab\times(f\vc{A})=\nab f\times\vc{A}+f\,\nab\times\vc{A}$.
\end{ejercicio}
\begin{resolucion}
Por la definición del rotacional en notación indicial y la regla del producto,
\begin{align*}
  \bigl[\nab\times(f\vc{A})\bigr]_i&=\varepsilon_{ijk}\,\partial_j(fA_k)
  =\varepsilon_{ijk}\bigl[(\partial_jf)\,A_k+f\,(\partial_jA_k)\bigr]\\
  &=\varepsilon_{ijk}(\partial_jf)\,A_k+f\,\varepsilon_{ijk}\,\partial_jA_k
  =(\nab f\times\vc{A})_i+f\,(\nab\times\vc{A})_i,
\end{align*}
donde en el último paso reconocemos el producto vectorial
$(\vc{u}\times\vc{A})_i=\varepsilon_{ijk}u_jA_k$ con $u_j=\partial_jf=(\nab f)_j$ y el
rotacional $(\nab\times\vc{A})_i=\varepsilon_{ijk}\partial_jA_k$. Al valer para cada
componente $i$, vale para el vector.
\end{resolucion}

\begin{ejercicio}[Divergencia de un rotacional]\label{ejer:div-rot}
Demostrar que $\nab\cdot(\nab\times\vc{A})=0$ para todo campo $\vc{A}\in C^{2}$.
\end{ejercicio}
\begin{resolucion}
\textbf{Idea.} Es la contracción de algo simétrico con algo antisimétrico.

En notación indicial, $\nab\cdot(\nab\times\vc{A})=\partial_i\bigl(\varepsilon_{ijk}\partial_jA_k\bigr)$
(la expresión entre paréntesis es la componente $i$-ésima del rotacional).
Como $\varepsilon_{ijk}$ es constante (\cref{lema:levi}), sale de la derivada:
\[
  \nab\cdot(\nab\times\vc{A})=\varepsilon_{ijk}\,\partial_i\partial_jA_k .
\]
Para cada $k$ fijo, $T_{ij}=\varepsilon_{ijk}$ es antisimétrico en $(i,j)$
(\cref{lema:levi}), mientras que $S_{ij}=\partial_i\partial_jA_k$ es simétrico
en $(i,j)$ porque $\vc{A}\in C^{2}$ (teorema de Schwarz): $\partial_i\partial_j=\partial_j\partial_i$.
Por el \cref{lema:sim-anti}, la contracción $S_{ij}T_{ij}$ es nula para cada $k$, y
al sumar en $k$ sigue siendo nula:
\[
  \nab\cdot(\nab\times\vc{A})=\varepsilon_{ijk}\,\partial_i\partial_jA_k=0 .
\]
\end{resolucion}

La misma técnica prueba la otra relación fundamental, y la identidad del doble
rotacional.

\begin{ejercicio}[Rotacional de un gradiente]\label{ejer:rot-grad}
Demostrar que $\nab\times(\nab f)=\vc{0}$ para todo $f\in C^{2}$.
\end{ejercicio}
\begin{resolucion}
$[\nab\times(\nab f)]_i=\varepsilon_{ijk}\,\partial_j(\partial_kf)=\varepsilon_{ijk}\,\partial_j\partial_kf$:
antisimétrico en $(j,k)$ contra simétrico en $(j,k)$ (Schwarz). Por el
\cref{lema:sim-anti}, es cero para cada $i$.
\end{resolucion}

\begin{ejercicio}[Rotacional del rotacional]\label{ejer:rotrot}
Demostrar que $\nab\times(\nab\times\vc{A})=\nab(\nab\cdot\vc{A})-\Lap\vc{A}$.
\end{ejercicio}
\begin{resolucion}
Aplicando dos veces la definición del rotacional,
\[
  [\nab\times(\nab\times\vc{A})]_i=\varepsilon_{ijk}\,\partial_j\bigl(\varepsilon_{klm}\,\partial_lA_m\bigr)
  =\varepsilon_{ijk}\,\varepsilon_{klm}\,\partial_j\partial_lA_m .
\]
Por la invariancia cíclica, $\varepsilon_{ijk}=\varepsilon_{kij}$, y la identidad $\varepsilon$--$\delta$
(\cref{lema:eps-delta}, con el índice contraído $k$ en primera posición),
\[
  \varepsilon_{kij}\,\varepsilon_{klm}=\delta_{il}\,\delta_{jm}-\delta_{im}\,\delta_{jl} .
\]
Por tanto
\[
  [\nab\times(\nab\times\vc{A})]_i=\partial_j\partial_iA_j-\partial_j\partial_jA_i
  =\partial_i\bigl(\partial_jA_j\bigr)-\Lap A_i
  =\bigl[\nab(\nab\cdot\vc{A})\bigr]_i-\bigl[\Lap\vc{A}\bigr]_i ,
\]
donde se han intercambiado $\partial_i$ y $\partial_j$ por Schwarz.
\end{resolucion}

% =========================================================================
\subsection{Seminario 4. Integrales de línea}\label{sem:linea}
% =========================================================================
\begin{ejercicio}[Trabajo a lo largo de una hélice]\label{ejer:helice}
Calcular el trabajo realizado por la fuerza
$\vc{F}(x,y,z)=-\tfrac12x\,\ux-\tfrac12y\,\uy+\tfrac14\,\uz$ al mover una partícula a
lo largo de la hélice $\vc{r}(t)=\cos t\,\ux+\sen t\,\uy+t\,\uz$ desde el punto
$(1,0,0)$ al punto $(-1,0,3\pi)$.
\end{ejercicio}
\begin{resolucion}
\textbf{Paso 1: parametrizar.} En coordenadas, $\gamma(t)=(\cos t,\sen t,t)$ con
$t\in[0,3\pi]$: en $t=0$ estamos en $(1,0,0)$ y en $t=3\pi$, en
$(\cos3\pi,\sen3\pi,3\pi)=(-1,0,3\pi)$. Es decir,
$x(t)=\cos t$, $y(t)=\sen t$, $z(t)=t$.

\textbf{Paso 2: derivar y restringir el campo.} El vector tangente es
$\gamma'(t)=(-\sen t,\cos t,1)$, y sobre la hélice
\[
  \vc{F}(\gamma(t))=\Bigl(-\tfrac12\cos t,\,-\tfrac12\sen t,\,\tfrac14\Bigr).
\]
\textbf{Paso 3: producto escalar.}
\[
  \vc{F}(\gamma(t))\cdot\gamma'(t)=\tfrac12\sen t\cos t-\tfrac12\sen t\cos t+\tfrac14=\tfrac14 .
\]
\textbf{Paso 4: integrar.}
\[
  W=\int_{C}\vc{F}\cdot\dl=\int_{0}^{3\pi}\frac14\,\dif t=\boxed{\frac{3\pi}{4}} .
\]
\end{resolucion}

\begin{observacion}[Comprobación cualitativa]
Los términos con senos y cosenos se han cancelado, y es lo esperable: la parte
horizontal del campo, $-\tfrac12(x\,\ux+y\,\uy)$, es radial (apunta hacia el eje),
mientras que la velocidad horizontal de la partícula es tangente a la
circunferencia, perpendicular a la radial; por tanto esa parte del campo no
realiza trabajo. Solo cuenta la componente vertical, $F_z=\tfrac14$ constante,
que actúa a lo largo de un ascenso total $\Delta z=3\pi$:
$W=\tfrac14\cdot3\pi$ (\cref{fig:sem-helice}).
\end{observacion}

\begin{ejercicio}[Integral de línea sobre un segmento]\label{ejer:segmento}
Calcular $\displaystyle\int_C(x^{2}-y+3z)\,\dif x$, donde $C$ es el segmento
rectilíneo desde el origen hasta $(1,2,1)$.
\end{ejercicio}
\begin{resolucion}
Usamos la parametrización estándar del segmento (\cref{ej:segmento}), con
$\vc{p}=(0,0,0)$ y $\vc{q}=(1,2,1)$:
$\sigma(t)=(1-t)(0,0,0)+t\,(1,2,1)=(t,2t,t)$, $t\in[0,1]$, $\sigma'(t)=(1,2,1)$.
En notación diferencial,
\[
  x=t\Rightarrow\dif x=\dif t,\qquad y=2t\Rightarrow\dif y=2\,\dif t,\qquad
  z=t\Rightarrow\dif z=\dif t .
\]
La 1-forma $(x^{2}-y+3z)\,\dif x$ solo tiene componente en $\dif x$ (equivale a
integrar el campo $(x^{2}-y+3z,\,0,\,0)$), de modo que
\[
  \int_C(x^{2}-y+3z)\,\dif x=\int_{0}^{1}\bigl(t^{2}-2t+3t\bigr)\dif t
  =\int_{0}^{1}(t^{2}+t)\,\dif t=\Bigl[\frac{t^{3}}{3}+\frac{t^{2}}{2}\Bigr]_{0}^{1}
  =\frac13+\frac12=\boxed{\frac56}.
\]
\end{resolucion}

\begin{ejercicio}[Una 1-forma sobre una circunferencia]\label{ejer:circulo}
Calcular $\displaystyle\int_Cx\,\dif y-y\,\dif x$, donde $C$ es la circunferencia
$x^{2}+y^{2}=R^{2}$ recorrida en sentido antihorario.
\end{ejercicio}
\begin{resolucion}
Parametrizamos con el ángulo polar, que es el parámetro natural:
$x=R\cos\varphi$, $y=R\sen\varphi$, $\varphi\in[0,2\pi]$. Entonces
$\dif x=-R\sen\varphi\,\dif\varphi$ y $\dif y=R\cos\varphi\,\dif\varphi$, y componiendo
directamente sobre el integrando diferencial,
\[
  x\,\dif y-y\,\dif x=(R\cos\varphi)(R\cos\varphi)\,\dif\varphi-(R\sen\varphi)(-R\sen\varphi)\,\dif\varphi
  =R^{2}(\cos^{2}\varphi+\sen^{2}\varphi)\,\dif\varphi=R^{2}\,\dif\varphi ,
\]
de donde
\[
  \int_Cx\,\dif y-y\,\dif x=\int_{0}^{2\pi}R^{2}\,\dif\varphi=R^{2}\bigl[\varphi\bigr]_{0}^{2\pi}=\boxed{2\pi R^{2}} .
\]
\end{resolucion}

\begin{observacion}[Interpretación geométrica]
La 1-forma $x\,\dif y-y\,\dif x$ es la integral de línea del campo de giro
$\vc{F}=-y\,\ux+x\,\uy$ (\cref{fig:campos_rot}), que sobre la circunferencia es
tangente, de módulo $R$ y en el sentido de recorrido: $\vc{F}\cdot\dl=R\cdot R\,\dif\varphi$.
Por eso la circulación es $R\cdot2\pi R=2\pi R^{2}$: el doble del área del círculo
($=\pi R^{2}$), en acuerdo con el teorema de Green, pues
$\nab\times\vc{F}=2\,\uz$ (\cref{ej:rot}) y $\oint\vc{F}\cdot\dl=\iint2\,\dif S=2\pi R^{2}$
(\cref{fig:sem-circulo}). Que no se anule indica que $\vc{F}$ \emph{no} es un
campo gradiente (\cref{cor:gradiente}).
\end{observacion}

% Figuras: trabajo sobre una hélice y circulación sobre una circunferencia
\begin{figure}[ht]
\centering
\tdplotsetmaincoords{72}{125}
\begin{tikzpicture}[tdplot_main_coords, scale=1.35]
  \draw[eje] (0,0,0) -- (1.9,0,0) node[below left] {\small $x$};
  \draw[eje] (0,0,0) -- (0,1.9,0) node[right] {\small $y$};
  \draw[eje] (0,0,0) -- (0,0,3.3) node[above] {\small $z$};
  % hélice (eje z comprimido: se dibuja 0.28 t en lugar de t)
  \draw[trayecto, domain=0:{3*pi}, samples=140, variable=\t, smooth]
    plot ({cos(deg(\t))},{sin(deg(\t))},{0.28*\t});
  \foreach \k in {0,1,...,6}{%
    \pgfmathsetmacro{\tt}{\k*pi/2*1.0}%
    \pgfmathsetmacro{\px}{cos(deg(\tt))}\pgfmathsetmacro{\py}{sin(deg(\tt))}\pgfmathsetmacro{\pz}{0.28*\tt}%
    % tangente (azul) y campo (ámbar/carmín)
    \draw[->, line width=1pt, azulprof] (\px,\py,\pz) -- ++({-0.55*sin(deg(\tt))},{0.55*cos(deg(\tt))},0.16);
    \draw[->, line width=1pt, cm3] (\px,\py,\pz) -- ++({-0.42*\px},{-0.42*\py},0.30);
    \fill (\px,\py,\pz) circle (1pt);}
  \fill[rojoteo] (1,0,0) circle (2pt) node[below right, text=black] {\small$(1,0,0)$};
  \fill[rojoteo] (-1,0,{0.28*3*pi}) circle (2pt) node[above left, text=black, fill=white, fill opacity=0.8, text opacity=1, inner sep=1pt] {\small$(-1,0,3\pi)$};
  \node[font=\small, text=azulprof, anchor=west] at (1.55,0.35,2.7) {$\gamma'(t)$};
  \node[font=\small, text=cm3, anchor=west] at (1.55,-0.75,2.2) {$\vc{F}$};
  \node[font=\small, text=rojoteo, anchor=west] at (-0.2,1.5,0.7) {$C$};
\end{tikzpicture}
\caption{Hélice del \cref{ejer:helice}. La parte horizontal de
$\vc{F}=-\tfrac12(x\,\ux+y\,\uy)+\tfrac14\,\uz$ es radial y, por tanto,
perpendicular a la velocidad horizontal $\gamma'$ (tangente a la circunferencia):
solo la componente vertical constante $F_z=\tfrac14$ trabaja. (El eje $z$ se
dibuja comprimido.)}
\label{fig:sem-helice}
\end{figure}

\begin{figure}[ht]
\centering
\begin{tikzpicture}[x=1.0cm, y=1.0cm]
  \shade[inner color=azulclaro!25, outer color=azulclaro!60] (0,0) circle (2.2);
  \clipD{-2.9}{2.9}{-2.9}{2.9}{\flechas{-\y}{\x}{-2.0,-1.5,...,2.0}{-2.0,-1.5,...,2.0}{0.36}{2.9}}
  \draw[trayecto] (0,0) circle (2.2);
  \foreach \a in {20,110,200,290}{%
    \draw[-{Stealth[length=7pt]}, rojoteo, line width=1.6pt] ({2.2*cos(\a)},{2.2*sin(\a)}) -- ++({-0.01*sin(\a)*40},{0.01*cos(\a)*40});}
  \draw[->, line width=0.5pt, grispizarra] (0,0) -- (-1.55,-1.55) node[midway, below left, font=\small] {$R$};
  \node[font=\small, text=rojoteo, anchor=west] at (2.3,1.9) {$C$};
  \node[font=\small, text=azulprof] at (0,0.55) {$\nab\times\vc{F}=2\,\uz$};
  \marcoD{-2.9}{2.9}{-2.9}{2.9}{-2,-1,0,1,2}{-2,-1,0,1,2}
  \ejesD{-2.9}{2.9}{-2.9}{2.9}
\end{tikzpicture}
\caption{Circunferencia del \cref{ejer:circulo}. Sobre $C$, el campo de giro
$\vc{F}=-y\,\ux+x\,\uy$ es tangente, de módulo $R$ y en el sentido de recorrido,
de modo que $\oint\vc{F}\cdot\dl=R\cdot2\pi R=2\pi R^{2}$, el doble del área del
disco, coherente con $\nab\times\vc{F}=2\,\uz$.}
\label{fig:sem-circulo}
\end{figure}


% =========================================================================
\subsection{Seminario 5. Verificación del teorema del gradiente}\label{sem:gradiente}
% =========================================================================
\begin{ejercicio}[Dos caminos, un mismo resultado]\label{ejer:verif-gradiente}
Comprobar que se cumple el teorema del gradiente,
$\int_{\vc{a}}^{\vc{b}}\nab f\cdot\dl=f(\vc{b})-f(\vc{a})$, para el campo escalar
$f:\R^{3}\to\R$ definido por $f(x,y,z)=xy^{2}$ y los puntos
$\vc{a}=(0,0,0)$ y $\vc{b}=(2,2,2)$. Hacer la integral por dos caminos
distintos y comprobar que el resultado es el mismo para ambos.
\end{ejercicio}
\begin{resolucion}
\textbf{El gradiente y el lado derecho del teorema.} El gradiente de $f=xy^{2}$ es
\[
  \nab f=\pd{f}{x}\,\ux+\pd{f}{y}\,\uy+\pd{f}{z}\,\uz=y^{2}\,\ux+2xy\,\uy+0\,\uz .
\]
La diferencia de potencial entre los extremos es
$f(\vc{b})-f(\vc{a})=f(2,2,2)-f(0,0,0)=(2)(2^{2})-0=8$. Este es el número que
tienen que reproducir las dos integrales.

\textbf{Camino 1: segmento rectilíneo $\gamma_1$.} La parametrización estándar
del segmento $[\vc{a},\vc{b}]$ es
$\gamma_1(t)=(1-t)\vc{a}+t\vc{b}=(2t,2t,2t)$, $t\in[0,1]$, de modo que
$\dl_1=\gamma_1'(t)\,\dif t=(2,2,2)\,\dif t$. Evaluamos el gradiente sobre la curva:
\[
  \nab f(\gamma_1(t))=\bigl((2t)^{2},\,2(2t)(2t),\,0\bigr)=(4t^{2},\,8t^{2},\,0).
\]
El producto escalar es $\nab f(\gamma_1)\cdot\gamma_1'=4t^{2}\cdot2+8t^{2}\cdot2+0\cdot2=24t^{2}$, y
\[
  \int_{\gamma_1}\nab f\cdot\dl=\int_{0}^{1}24t^{2}\,\dif t=\Bigl[\frac{24}{3}t^{3}\Bigr]_{0}^{1}=8 .
\]
Coincide con $f(\vc{b})-f(\vc{a})$, como asegura el teorema.

\textbf{Camino 2: ruta ortogonal por los ejes $\gamma_2$.} Ahora elegimos una
trayectoria escalonada, paralela a los ejes cartesianos y dividida en tres
subsegmentos $C_x$, $C_y$, $C_z$ (\cref{fig:sem-trayectorias}), y sumamos sus
contribuciones por aditividad (\cref{prop:prop-linea}):
$\gamma_2=C_x\uplus C_y\uplus C_z$.
\begin{itemize}
  \item \emph{Tramo $C_x$ (recorrido en $x$)}, de $\vc{a}=(0,0,0)$ a $(2,0,0)$:
        $\gamma_{2,x}(t)=(2t,0,0)$, $t\in[0,1]$, $\dl_x=(2,0,0)\,\dif t$. Aquí
        $y=0$, de modo que $\nab f=(y^{2},2xy,0)=(0,0,0)$ y
        $\int_{C_x}\nab f\cdot\dl_x=0$.
  \item \emph{Tramo $C_y$ (recorrido en $y$)}, de $(2,0,0)$ a $(2,2,0)$:
        $\gamma_{2,y}(t)=(2,t,0)$, $t\in[0,2]$, $\dl_y=(0,1,0)\,\dif t$. El
        gradiente sobre el tramo es $\nab f=(t^{2},\,2\cdot2\cdot t,\,0)=(t^{2},4t,0)$,
        y el producto escalar con $\dl_y$ vale $4t\,\dif t$:
        \[
          \int_{C_y}\nab f\cdot\dl_y=\int_{0}^{2}4t\,\dif t=\bigl[2t^{2}\bigr]_{0}^{2}=8 .
        \]
  \item \emph{Tramo $C_z$ (recorrido en $z$)}, de $(2,2,0)$ a $(2,2,2)$:
        $\gamma_{2,z}(t)=(2,2,t)$, $t\in[0,2]$, $\dl_z=(0,0,1)\,\dif t$. El gradiente
        es $\nab f=(4,8,0)$, perpendicular a $\dl_z$:
        $\nab f\cdot\dl_z=0$ y $\int_{C_z}\nab f\cdot\dl_z=0$.
\end{itemize}
Sumando, $\int_{\gamma_2}\nab f\cdot\dl=0+8+0=8=f(\vc{b})-f(\vc{a})$.
\end{resolucion}

\begin{observacion}[Qué hemos comprobado]
Las dos integrales dan $8$ aunque los caminos son muy distintos: es la
\emph{independencia del camino} (\cref{cor:gradiente}). Los tramos tienen
sentido físico: $f=xy^{2}$ no depende de $z$, así que el tramo vertical no
cambia $f$; y mientras $y=0$ la función vale $0$ y no cambia, por lo que todo el
incremento $f:0\to8$ ocurre en el tramo $C_y$, donde $f=2y^{2}$ pasa de $0$ a $8$. Además,
al recorrer $\gamma_1$ y volver por $\gamma_2$ en sentido inverso se obtiene un
circuito cerrado con circulación $8-8=0$ (\cref{eq:circulacion-nula}).
\end{observacion}

% Figura: verificación del teorema del gradiente por dos caminos
\begin{figure}[ht]
\centering
\tdplotsetmaincoords{62}{112}
\begin{tikzpicture}[tdplot_main_coords, scale=1.9]
  \draw[eje] (0,0,0) -- (2.8,0,0) node[below left] {\small $x$};
  \draw[eje] (0,0,0) -- (0,2.8,0) node[right] {\small $y$};
  \draw[eje] (0,0,0) -- (0,0,2.8) node[above] {\small $z$};
  % cota de referencia: aristas del cubo
  \draw[punteado] (2,0,0) -- (2,2,0) -- (0,2,0) (2,2,0) -- (2,2,2) (2,0,0) -- (2,0,2) -- (2,2,2) (0,2,0) -- (0,2,2) -- (2,2,2) (0,0,2) -- (2,0,2) (0,0,2) -- (0,2,2);
  % camino 1: segmento recto
  \draw[line width=1.6pt, rojoteo] (0,0,0) -- (2,2,2);
  \draw[->, line width=1.6pt, rojoteo] (0.9,0.9,0.9) -- (1.05,1.05,1.05);
  % camino 2: poligonal por los ejes
  \draw[line width=1.6pt, azulprof] (0,0,0) -- (2,0,0) -- (2,2,0) -- (2,2,2);
  \draw[->, line width=1.6pt, azulprof] (0.9,0,0) -- (1.1,0,0);
  \draw[->, line width=1.6pt, azulprof] (2,0.9,0) -- (2,1.1,0);
  \draw[->, line width=1.6pt, azulprof] (2,2,0.9) -- (2,2,1.1);
  % gradiente grad f = (y^2, 2xy, 0) sobre C_y y sobre el camino 1
  \foreach \t in {0.4,0.8,1.2,1.6}{%
    \pgfmathsetmacro{\gx}{\t*\t}\pgfmathsetmacro{\gy}{4*\t}\pgfmathsetmacro{\gn}{sqrt(\gx*\gx+\gy*\gy)}
    \mapcolor{\gn/6.6}%
    \draw[-{Stealth[length=4pt]}, line width=0.9pt, mapc] (2,\t,0) -- ++({0.09*\gx},{0.09*\gy},0);}
  \foreach \t in {0.4,0.8,1.2,1.6}{%
    \pgfmathsetmacro{\gx}{\t*\t}\pgfmathsetmacro{\gy}{2*\t*\t}\pgfmathsetmacro{\gn}{sqrt(\gx*\gx+\gy*\gy)}
    \mapcolor{\gn/6.6}%
    \draw[-{Stealth[length=4pt]}, line width=0.9pt, mapc] (\t,\t,\t) -- ++({0.09*\gx},{0.09*\gy},0);}
  \fill[azulprof] (0,0,0) circle (2.2pt) node[below left, text=black] {\small $\vc{a}=(0,0,0)$};
  \fill[azulprof] (2,2,2) circle (2.2pt) node[above right, text=black] {\small $\vc{b}=(2,2,2)$};
  \node[font=\small, text=rojoteo, anchor=east] at (0.85,0.55,1.05) {$\gamma_1$};
  \node[font=\small, text=azulprof] at (1,-0.25,0) {$C_x$};
  \node[font=\small, text=azulprof, anchor=west] at (2.15,1.1,0) {$C_y$};
  \node[font=\small, text=azulprof, anchor=west] at (2.12,2,1.0) {$C_z$};
\end{tikzpicture}
\caption{Verificación del teorema del gradiente para $f=xy^{2}$ entre
$\vc{a}=(0,0,0)$ y $\vc{b}=(2,2,2)$. En rojo, el segmento $\gamma_1$; en azul, la
ruta escalonada $\gamma_2=C_x\uplus C_y\uplus C_z$. Se dibuja también
$\nab f=(y^{2},2xy,0)$ (color según su módulo): en $C_x$ es nulo y en $C_z$ es
perpendicular al camino, de modo que el incremento $f=0\to8$ lo aporta solo
$C_y$. Ambos caminos dan $\int\nab f\cdot\dl=8$.}
\label{fig:sem-trayectorias}
\end{figure}


% =========================================================================
\subsection{Seminario 6. Divergencia y teorema de Gauss}\label{sem:gauss}
% =========================================================================
\begin{ejercicio}[Divergencia de $R^{n}\uR$]\label{ejer:div-Rn}
Demostrar que $\nab\cdot(R^{n}\,\uR)=(n+2)\,R^{n-1}$ para $R\neq0$
(con $n\neq-2$ para que el resultado no sea idénticamente nulo), donde
$\vc{R}=\vc{r}-\vc{r}\,'$.
\end{ejercicio}
\begin{resolucion}
\textbf{Idea.} Ya lo demostramos con reglas de derivación del producto
(\cref{prop:div-rot-Rn}); ahora lo hacemos a mano, componente a componente,
para ver el mecanismo.

Escribimos $R^{n}\uR=R^{n-1}\vc{R}=R^{n-1}\bigl(R_x\,\ux+R_y\,\uy+R_z\,\uz\bigr)$ con
$R_x=x-x'$, etc. La divergencia es la suma de tres derivadas parciales; la primera
requiere la regla del producto:
\[
  \pd{}{x}\bigl[R^{n-1}(x-x')\bigr]=(x-x')\,\pd{R^{n-1}}{x}+R^{n-1}\,\pd{(x-x')}{x} .
\]
Como $\partial R/\partial x=(x-x')/R$ (\cref{eq:dR}), la regla de la cadena da
\[
  \pd{R^{n-1}}{x}=\pd{R^{n-1}}{R}\,\pd{R}{x}=(n-1)R^{n-2}\,\frac{x-x'}{R}
  =(n-1)\,R^{n-3}(x-x') .
\]
Por tanto
$\pd{}{x}\bigl[R^{n-1}(x-x')\bigr]=(n-1)R^{n-3}(x-x')^{2}+R^{n-1}$, y análogamente para
$y$ y $z$:
\[
  \pd{}{y}\bigl[R^{n-1}(y-y')\bigr]=(n-1)R^{n-3}(y-y')^{2}+R^{n-1},\qquad
  \pd{}{z}\bigl[R^{n-1}(z-z')\bigr]=(n-1)R^{n-3}(z-z')^{2}+R^{n-1}.
\]
Sumando los tres, y usando $(x-x')^{2}+(y-y')^{2}+(z-z')^{2}=R^{2}$,
\[
  \nab\cdot\bigl(R^{n}\uR\bigr)=(n-1)\,R^{n-3}R^{2}+3R^{n-1}=(n-1)R^{n-1}+3R^{n-1}=\boxed{(n+2)\,R^{n-1}} .
\]
Este resultado vale para $R\neq0$, es decir, $\vc{r}\neq\vc{r}\,'$. Para $n=-2$,
$\nab\cdot(R^{-2}\uR)=(-2+2)R^{-3}=0$: la divergencia de $\uR/R^{2}$ es nula en todo
punto distinto de la fuente (con la contribución singular de \cref{teo:delta}
en $R=0$).
\end{resolucion}

\begin{ejercicio}[Jacobiano de las coordenadas esféricas]\label{ejer:jacobiano}
Calcular el determinante jacobiano $\partial(x,y,z)/\partial(r,\theta,\varphi)$ del
cambio a coordenadas esféricas.
\end{ejercicio}
\begin{resolucion}
Con $x=r\sen\theta\cos\varphi$, $y=r\sen\theta\sen\varphi$, $z=r\cos\theta$, la
matriz de las derivadas parciales respecto de $(r,\theta,\varphi)$ es
\[
  \frac{\partial(x,y,z)}{\partial(r,\theta,\varphi)}=
  \begin{pmatrix}
    \sen\theta\cos\varphi & r\cos\theta\cos\varphi & -r\sen\theta\sen\varphi\\
    \sen\theta\sen\varphi & r\cos\theta\sen\varphi & \ r\sen\theta\cos\varphi\\
    \cos\theta & -r\sen\theta & 0
  \end{pmatrix}.
\]
Desarrollamos el determinante por la tercera fila (la que tiene un cero):
\begin{align*}
  \det&=\cos\theta\begin{vmatrix}r\cos\theta\cos\varphi & -r\sen\theta\sen\varphi\\ r\cos\theta\sen\varphi & r\sen\theta\cos\varphi\end{vmatrix}
  -(-r\sen\theta)\begin{vmatrix}\sen\theta\cos\varphi & -r\sen\theta\sen\varphi\\ \sen\theta\sen\varphi & r\sen\theta\cos\varphi\end{vmatrix}+0.
\end{align*}
El primer menor vale
$r^{2}\cos\theta\sen\theta\cos^{2}\varphi+r^{2}\cos\theta\sen\theta\sen^{2}\varphi=r^{2}\sen\theta\cos\theta$,
y el segundo
$r\sen^{2}\theta\cos^{2}\varphi+r\sen^{2}\theta\sen^{2}\varphi=r\sen^{2}\theta$
(en ambos, $\cos^{2}\varphi+\sen^{2}\varphi=1$). Entonces
\[
  \det=\cos\theta\,(r^{2}\sen\theta\cos\theta)+r\sen\theta\,(r\sen^{2}\theta)
  =r^{2}\sen\theta\,(\cos^{2}\theta+\sen^{2}\theta)=r^{2}\sen\theta .
\]
El jacobiano (valor absoluto) es $J=|r^{2}\sen\theta|=r^{2}\sen\theta$, porque
$\theta\in[0,\pi]$ implica $\sen\theta\ge0$. Concuerda con el producto de
factores de escala $h_1h_2h_3=1\cdot r\cdot r\sen\theta$ (\cref{prop:elementos}).
\end{resolucion}

\begin{ejercicio}[Teorema de Gauss en una esfera]\label{ejer:gauss-esfera}
Comprobar que se cumple el teorema de la divergencia,
$\int_V\nab\cdot\vc{A}\,\dV=\oint_S\vc{A}\cdot\dS$, siendo $\vc{A}(\vc{r})=r^{2}\,\ur$
(en coordenadas esféricas) y $V$ una esfera de radio $R$ centrada en el origen.
\end{ejercicio}
\begin{resolucion}
\textbf{Idea.} El campo es radial, crece como $r^{2}$ y «sale» de la esfera por
todas partes. Calculamos por separado el lado de volumen y el lado de superficie
y comprobamos que coinciden.

\textbf{Lado del volumen.} La divergencia del campo $\vc{A}=(A_r,A_\theta,A_\varphi)=(r^{2},0,0)$ en esféricas
(\cref{eq:ops-esf}) solo conserva el primer término, pues $A_\theta=A_\varphi=0$:
\[
  \nab\cdot\vc{A}=\frac{1}{r^{2}}\pd{}{r}\bigl(r^{2}A_r\bigr)=\frac{1}{r^{2}}\pd{}{r}\bigl(r^{4}\bigr)=\frac{4r^{3}}{r^{2}}=4r .
\]
La esfera sólida es
$V=\{(x,y,z):x^{2}+y^{2}+z^{2}\le R^{2}\}=\{(r,\theta,\varphi):0\le r\le R,\ 0\le\theta\le\pi,\ 0\le\varphi\le2\pi\}$,
y con el elemento de volumen $\dV=r^{2}\sen\theta\,\dif r\,\dif\theta\,\dif\varphi$
(\cref{ejer:jacobiano}),
\[
  \iiint_V\nab\cdot\vc{A}\,\dV=\int_{0}^{2\pi}\!\!\int_{0}^{\pi}\!\!\int_{0}^{R}(4r)\,(r^{2}\sen\theta)\,\dif r\,\dif\theta\,\dif\varphi .
\]
El integrando es un producto de funciones de una sola variable y los límites son
constantes: por el teorema de Fubini, la integral se factoriza,
\begin{align*}
  &=\Bigl(\int_0^{2\pi}\dif\varphi\Bigr)\Bigl(\int_0^{\pi}\sen\theta\,\dif\theta\Bigr)\Bigl(\int_0^{R}4r^{3}\,\dif r\Bigr)
  =\bigl[\varphi\bigr]_0^{2\pi}\cdot\bigl[-\cos\theta\bigr]_0^{\pi}\cdot\bigl[r^{4}\bigr]_0^{R}\\
  &=2\pi\cdot\bigl[-\cos\pi+\cos0\bigr]\cdot R^{4}=2\pi\cdot2\cdot R^{4}=4\pi R^{4}.
\end{align*}

\textbf{Lado de la superficie.} Sobre la esfera cerrada $S=\partial V$ el radio es
constante, $r=R$, y el campo es $\vc{A}=R^{2}\ur$, de módulo constante y radial. La
normal unitaria exterior de la esfera es $\un=\ur$ (el radio normalizado), de modo
que $\vc{A}\cdot\un=R^{2}\,\ur\cdot\ur=R^{2}$ (producto de dos versores de una base
ortonormal: $\langle\uh{i},\uh{j}\rangle=\delta_{ij}$). El elemento de área en esféricas a $r=R$ es
$\dsc=R^{2}\sen\theta\,\dif\theta\,\dif\varphi$ (\cref{sec:esfericas}), por lo que
\[
  \oiint_S\vc{A}\cdot\un\,\dsc=\int_0^{2\pi}\!\!\int_0^{\pi}R^{2}\cdot R^{2}\sen\theta\,\dif\theta\,\dif\varphi
  =R^{4}\Bigl(\int_0^{2\pi}\dif\varphi\Bigr)\Bigl(\int_0^{\pi}\sen\theta\,\dif\theta\Bigr)
  =R^{4}\cdot2\pi\cdot2=4\pi R^{4}.
\]

\textbf{Conclusión.} Ambos lados valen $4\pi R^{4}$, así que
$\iiint_V\nab\cdot\vc{A}\,\dV=\oiint_S\vc{A}\cdot\dS=4\pi R^{4}$ y el teorema de Gauss se
verifica para este campo radial y geometría esférica.
\end{resolucion}

\begin{observacion}
Es instructivo ver por qué sale así. El campo crece hacia fuera ($A_r=r^{2}$) y
cada capa esférica «deja pasar» más flujo hacia fuera de la que recibe de
dentro; esa diferencia es lo que mide la divergencia ($4r>0$ en todo punto,
fuente distribuida). La suma de todas esas contribuciones infinitesimales,
ponderadas por el volumen de cada capa, reproduce el flujo total que atraviesa
la superficie exterior (\cref{fig:sem-esfera}).
\end{observacion}

% Figura: verificación del teorema de la divergencia en una esfera
\begin{figure}[ht]
\centering
\begin{minipage}[t]{0.52\textwidth}\centering
\tdplotsetmaincoords{70}{125}
\begin{tikzpicture}[tdplot_main_coords, scale=1.05, baseline=(current bounding box.north)]
  \def\RR{2.1}
  \draw[eje] (0,0,0) -- (3.0,0,0) node[below left] {\small $x$};
  \draw[eje] (0,0,0) -- (0,3.1,0) node[right] {\small $y$};
  \draw[eje] (0,0,0) -- (0,0,3.1) node[above] {\small $z$};
  % esfera: contorno aparente y círculos máximos
  \begin{scope}[tdplot_screen_coords]
    \shade[ball color=azulclaro!55, opacity=0.35] (0,0) circle (\RR);
    \draw[azulprof!60, line width=0.6pt] (0,0) circle (\RR);
  \end{scope}
  \draw[azulprof!45, line width=0.4pt, domain=0:360, samples=72] plot ({\RR*cos(\x)},{\RR*sin(\x)},0);
  \foreach \ph in {0,45,90,135}{\draw[azulprof!35, line width=0.4pt, domain=0:360, samples=72]
     plot ({\RR*sin(\x)*cos(\ph)},{\RR*sin(\x)*sin(\ph)},{\RR*cos(\x)});}
  % campo radial A = R^2 u_r (flechas salientes, módulo constante)
  \foreach \th/\ph in {35/20,35/200,70/110,70/290,110/20,110/200,145/110,145/290}{%
    \pgfmathsetmacro{\px}{\RR*sin(\th)*cos(\ph)}\pgfmathsetmacro{\py}{\RR*sin(\th)*sin(\ph)}\pgfmathsetmacro{\pz}{\RR*cos(\th)}%
    \draw[-{Stealth[length=4pt]}, line width=1pt, acento] (\px,\py,\pz) -- ++({0.55*sin(\th)*cos(\ph)},{0.55*sin(\th)*sin(\ph)},{0.55*cos(\th)});}
  % elemento de superficie y normal
  \def\ta{48}\def\tb{62}\def\pa{28}\def\pb{46}
  \fill[ocreteo, opacity=0.7]
    ({\RR*sin(\ta)*cos(\pa)},{\RR*sin(\ta)*sin(\pa)},{\RR*cos(\ta)}) --
    ({\RR*sin(\tb)*cos(\pa)},{\RR*sin(\tb)*sin(\pa)},{\RR*cos(\tb)}) --
    ({\RR*sin(\tb)*cos(\pb)},{\RR*sin(\tb)*sin(\pb)},{\RR*cos(\tb)}) --
    ({\RR*sin(\ta)*cos(\pb)},{\RR*sin(\ta)*sin(\pb)},{\RR*cos(\ta)}) -- cycle;
  \draw[-{Stealth[length=5pt]}, line width=1.5pt, azulacero]
    ({\RR*sin(55)*cos(37)},{\RR*sin(55)*sin(37)},{\RR*cos(55)}) --
    ({(\RR+1.0)*sin(55)*cos(37)},{(\RR+1.0)*sin(55)*sin(37)},{(\RR+1.0)*cos(55)});
  \node[font=\small, text=azulacero, anchor=south west] at ({(\RR+1.0)*sin(55)*cos(37)},{(\RR+1.0)*sin(55)*sin(37)},{(\RR+1.0)*cos(55)}) {$\dS$};
  \fill (0,0,0) circle (1.2pt);
\end{tikzpicture}\\[2pt]{\small\scshape (a)}
\end{minipage}\hfill
\begin{minipage}[t]{0.44\textwidth}\centering
\begin{tikzpicture}[baseline=(current bounding box.north)]
  \begin{axis}[tema, width=\textwidth, height=6.3cm, xlabel={$r$}, xmin=0, xmax=1.6, ymin=0, ymax=42,
      samples=100, domain=0:1.6, xtick={0,0.5,1,1.5}, xticklabels={0,0{,}5,$R$,1{,}5},
      xtick={0,0.5,1.0,1.5}, legend style={at={(0.03,0.97)}, anchor=north west}]
    \addplot[azulacero, line width=1.1pt] {4*pi*x^4};
    \addplot[acento, line width=1.1pt] {16*pi*x^3};
    \addplot[name path=curva, acento, draw=none, domain=0:1] {16*pi*x^3};
    \path[name path=eje] (axis cs:0,0) -- (axis cs:1,0);
    \addplot[acento!25] fill between[of=curva and eje];
    \addplot[grispizarra, densely dashed] coordinates {(1,0) (1,12.57)};
    \node[font=\scriptsize, text=azulacero, anchor=west] at (axis cs:1.05,9.5) {$4\pi R^{4}$};
    \legend{$\Phi(r)=4\pi r^{4}$, $16\pi r^{3}=\nab\cdot\vc{A}\,\cdot4\pi r^{2}$}
  \end{axis}
\end{tikzpicture}\\[2pt]{\small\scshape (b)}
\end{minipage}
\caption{Teorema de Gauss para $\vc{A}=r^{2}\,\ur$ en una esfera de radio
$R$. (a)~El campo es radial y de módulo constante sobre la esfera; la normal
exterior es $\ur$, de modo que $\vc{A}\cdot\un=R^{2}$ y el flujo es
$R^{2}\cdot4\pi R^{2}$. (b)~Visión «capa a capa»: cada cáscara de radio $r$ y
espesor $\dif r$ aporta $\nab\cdot\vc{A}\,\dV=4r\cdot4\pi r^{2}\,\dif r$; el área
bajo $16\pi r^{3}$ (sombreada) hasta $R=1$ es el flujo total
$\Phi(R)=4\pi R^{4}$.}
\label{fig:sem-esfera}
\end{figure}


% =========================================================================
\subsection{Seminario 7. Verificación del teorema de Stokes}\label{sem:stokes}
% =========================================================================
\begin{ejercicio}[Stokes en un triángulo]\label{ejer:stokes-triangulo}
Comprobar que se cumple el teorema del rotacional,
$\iint_S(\nab\times\vc{A})\cdot\dS=\oint_{\partial S}\vc{A}\cdot\dl$, siendo
$\vc{A}(\vc{r})=xy\,\ux+2yz\,\uy+3zx\,\uz$ y $S$ el área encerrada por el triángulo
rectángulo de vértices $\vc{a}=(0,0,0)$, $\vc{b}=(2,0,0)$ y $\vc{c}=(0,2,0)$ (en el
plano $z=0$).
\end{ejercicio}
\begin{resolucion}
\textbf{Idea.} Calculamos el flujo del rotacional por el triángulo y la
circulación del campo por sus tres lados, y comprobamos que coinciden. Antes de
empezar fijamos la orientación: recorremos el borde $L=\partial S$ en el orden
$\vc{a}\to\vc{b}\to\vc{c}\to\vc{a}$ (sentido antihorario visto desde $+z$); por la
regla de la mano derecha, la normal compatible es $+\uz$ (\cref{fig:sem-stokes}).

\textbf{1. Integral de superficie.} \emph{El rotacional:}
\[
  \nab\times\vc{A}=
  \begin{vmatrix}
    \ux & \uy & \uz\\
    \partial_x & \partial_y & \partial_z\\
    xy & 2yz & 3zx
  \end{vmatrix}
  =(0-2y)\,\ux-(3z-0)\,\uy+(0-x)\,\uz=-2y\,\ux-3z\,\uy-x\,\uz .
\]
\emph{La superficie:} $S$ está en el plano $z=0$; la parametrizamos con
$\sigma(x,y)=(x,y,0)$, con $(x,y)\in S^{*}$ la región plana limitada por los catetos y la
hipotenusa $y=-x+2$:
\[
  S^{*}=\{(x,y):0\le x\le2,\ 0\le y\le-x+2\}.
\]
Las derivadas parciales son $\partial_x\sigma=(1,0,0)=\ux$ y
$\partial_y\sigma=(0,1,0)=\uy$, y el vector normal asociado es
$\vc{n}_\sigma=\partial_x\sigma\times\partial_y\sigma=\ux\times\uy=\uz$, compatible con la orientación
elegida. \emph{El integrando:} evaluado sobre la superficie ($z=0$),
$(\nab\times\vc{A})(\sigma(x,y))=(-2y,\,0,\,-x)$, y su producto escalar con la normal
es $(\nab\times\vc{A})\cdot\uz=-x$. Por tanto
\[
  \iint_S(\nab\times\vc{A})\cdot\dS=\int_0^2\!\!\int_0^{-x+2}(-x)\,\dif y\,\dif x
  =\int_0^2-x\bigl[y\bigr]_0^{-x+2}\dif x=\int_0^2(x^{2}-2x)\,\dif x
  =\Bigl[\frac{x^{3}}{3}-x^{2}\Bigr]_0^2=\frac83-4=-\frac43 .
\]

\textbf{2. Integral de línea (circulación).} Dividimos el contorno en tres tramos,
$L=\Gamma_1\uplus\Gamma_2\uplus\Gamma_3$:
\begin{itemize}
  \item \emph{Tramo $\Gamma_1$} ($\vc{a}\to\vc{b}$, a lo largo del eje $x$): $\gamma_1(t)=(t,0,0)$, $t\in[0,2]$,
        $\dl=(1,0,0)\,\dif t$. Aquí $y=z=0$, luego $\vc{A}=(xy,2yz,3zx)=\vc{0}$ y $\int_{\Gamma_1}\vc{A}\cdot\dl=0$.
  \item \emph{Tramo $\Gamma_2$} ($\vc{b}\to\vc{c}$, la hipotenusa): $z=0$, $y=-x+2$, y por
        tanto $\dif y=-\dif x$. Parametrizamos con $t\in[0,2]$ mediante
        $\gamma_2(t)=(2-t,\,t,\,0)$ (es el segmento $(1-s)\vc{b}+s\,\vc{c}$ con $s=t/2$), de modo que
        $\gamma_2'(t)=(-1,1,0)$ y $\dl=(-1,1,0)\,\dif t$. Sobre el tramo, $\vc{A}=\bigl((2-t)t,\,0,\,0\bigr)$ (porque $z=0$), y
        $\vc{A}\cdot\gamma_2'=-(2-t)t=t^{2}-2t$, así que
        \[
          \int_{\Gamma_2}\vc{A}\cdot\dl=\int_0^2(t^{2}-2t)\,\dif t=\Bigl[\frac{t^{3}}{3}-t^{2}\Bigr]_0^2=\frac83-4=-\frac43 .
        \]
  \item \emph{Tramo $\Gamma_3$} ($\vc{c}\to\vc{a}$, a lo largo del eje $y$): $\gamma_3(t)=(0,2-t,0)$, $t\in[0,2]$,
        $\dl=(0,-1,0)\,\dif t$. Aquí $x=z=0$, luego $\vc{A}=\vc{0}$ y la integral es nula.
\end{itemize}
Sumando,
\[
  \oint_{\partial S}\vc{A}\cdot\dl=0-\frac43+0=-\frac43 .
\]

\textbf{Conclusión.} Ambos lados valen $-4/3$:
\[
  \iint_S(\nab\times\vc{A})\cdot\dS=\oint_{\partial S}\vc{A}\cdot\dl=-\frac43 ,
\]
y el teorema de Stokes se verifica formalmente.
\end{resolucion}

\begin{observacion}[Lectura geométrica]
Solo el lado $\Gamma_2$ contribuye a la circulación, pues sobre los ejes
coordenados ($y=0$ o $x=0$, con $z=0$) el campo se anula. En el interior del
triángulo, la componente que cuenta del rotacional, $(\nab\times\vc{A})\cdot\uz=-x$, es
negativa y más intensa hacia el vértice $\vc{b}$ (donde $x$ es mayor) (\cref{fig:sem-stokes}):
el campo tiende a hacer girar una rueda en sentido \emph{horario} visto desde
$+z$, justo al contrario del sentido en que recorremos el borde, por eso el
resultado es negativo.
\end{observacion}

% Figura: verificación del teorema de Stokes en un triángulo
\begin{figure}[ht]
\centering
\begin{minipage}[t]{0.48\textwidth}\centering
\begin{tikzpicture}[>=Stealth, x=1.9cm, y=1.9cm, baseline=(current bounding box.north)]
  \draw[eje] (-0.25,0) -- (2.55,0) node[right, text=black] {$x$};
  \draw[eje] (0,-0.25) -- (0,2.55) node[above, text=black] {$y$};
  \fill[azulclaro!55] (0,0) -- (2,0) -- (0,2) -- cycle;
  \draw[trayecto] (0,0) -- (2,0) -- (0,2) -- cycle;
  \draw[-{Stealth[length=8pt]}, rojoteo, line width=1.6pt] (0.85,0) -- (1.15,0);
  \draw[-{Stealth[length=8pt]}, rojoteo, line width=1.6pt] (1.07,0.93) -- (0.93,1.07);
  \draw[-{Stealth[length=8pt]}, rojoteo, line width=1.6pt] (0,1.15) -- (0,0.85);
  \fill[azulprof] (0,0) circle (1.8pt) node[below left, text=black] {$\vc{a}$};
  \fill[azulprof] (2,0) circle (1.8pt) node[below right, text=black] {$\vc{b}$};
  \fill[azulprof] (0,2) circle (1.8pt) node[above left, text=black] {$\vc{c}$};
  \node[font=\small, text=rojoteo, anchor=north] at (1,-0.05) {$\Gamma_1$};
  \node[font=\small, text=rojoteo, anchor=south west] at (1.05,1.05) {$\Gamma_2:\ y=-x+2$};
  \node[font=\small, text=rojoteo, anchor=east] at (-0.06,1) {$\Gamma_3$};
  \node[font=\small, text=azulprof] at (0.48,0.42) {$S$};
  % normal saliendo del plano
  \draw[azulacero, line width=0.8pt] (0.62,0.62) circle (0.075);
  \fill[azulacero] (0.62,0.62) circle (0.02);
  \node[font=\footnotesize, text=azulacero, anchor=south west] at (0.66,0.64) {$\un=\uz$};
  \node[font=\scriptsize, text=black!75, anchor=north west] at (0.15,-0.28) {$\vc{A}=\vc{0}$ sobre $\Gamma_1$ y $\Gamma_3$};
\end{tikzpicture}\\[2pt]{\small\scshape (a)}\ contorno y orientación
\end{minipage}\hfill
\begin{minipage}[t]{0.48\textwidth}\centering
\begin{tikzpicture}[>=Stealth, x=1.9cm, y=1.9cm, baseline=(current bounding box.north)]
  \foreach \i in {0,...,79}{%
    \pgfmathsetmacro{\xa}{2*\i/80}\pgfmathsetmacro{\xb}{2*(\i+1)/80}%
    \pgfmathsetmacro{\tt}{(\xa+\xb)/4}\mapcolor{\tt}%
    \fill[mapc, draw=mapc, line width=0.2pt] (\xa,0) -- (\xb,0) -- (\xb,{2-\xb}) -- (\xa,{2-\xa}) -- cycle;}
  \draw[eje] (-0.25,0) -- (2.55,0) node[right, text=black] {$x$};
  \draw[eje] (0,-0.25) -- (0,2.55) node[above, text=black] {$y$};
  \draw[black!60, line width=0.8pt] (0,0) -- (2,0) -- (0,2) -- cycle;
  % escala de color
  \foreach \i in {0,...,29}{%
    \pgfmathsetmacro{\yy}{0.2+1.4*\i/30}\pgfmathsetmacro{\tt}{\i/30}\mapcolor{\tt}%
    \fill[mapc] (2.35,\yy) rectangle (2.5,{\yy+1.4/30});}
  \draw[black!60, line width=0.4pt] (2.35,0.2) rectangle (2.5,1.6);
  \node[font=\scriptsize, anchor=west] at (2.5,0.2) {$0$};
  \node[font=\scriptsize, anchor=west] at (2.5,0.9) {$-1$};
  \node[font=\scriptsize, anchor=west] at (2.5,1.6) {$-2$};
  \node[font=\scriptsize, anchor=south] at (2.42,1.68) {$(\nab\times\vc{A})\cdot\uz$};
  \node[font=\small, text=white, fill=cm0, fill opacity=0.65, text opacity=1, inner sep=2pt] at (0.62,0.55) {$\iint_S=-\tfrac43$};
\end{tikzpicture}\\[2pt]{\small\scshape (b)}\ integrando $(\nab\times\vc{A})\cdot\uz=-x$
\end{minipage}
\caption{Verificación del teorema de Stokes para $\vc{A}=xy\,\ux+2yz\,\uy+3zx\,\uz$ sobre el
triángulo de vértices $\vc{a},\vc{b},\vc{c}$. (a)~El borde se recorre
$\vc{a}\to\vc{b}\to\vc{c}\to\vc{a}$ (antihorario visto desde $+z$) y, por la
regla de la mano derecha, la normal es $\uz$. El campo se anula sobre los ejes y
solo $\Gamma_2$ contribuye a la circulación. (b)~La componente normal del
rotacional $-x$ es más intensa (en valor absoluto) hacia el vértice $\vc{b}$; su
integral sobre el triángulo es $-4/3$, igual a la circulación $-4/3$.}
\label{fig:sem-stokes}
\end{figure}


% =========================================================================
\appendix
\section{Fórmulas de referencia}\label{sec:apendice}
% =========================================================================
Se recogen aquí las expresiones del formulario del curso~\cite{FORMULARIO} que
no han aparecido todavía de forma explícita, con una breve justificación.

% -------------------------------------------------------------------------
\subsection{Doble producto de vectores}\label{sec:doble-producto}
% -------------------------------------------------------------------------
\begin{proposicion}\label{prop:doble-producto}
Para tres vectores $\vc{A},\vc{B},\vc{C}$ de $\R^{3}$,
\begin{align}
  \vc{A}\cdot(\vc{B}\times\vc{C})&=\vc{B}\cdot(\vc{C}\times\vc{A})=\vc{C}\cdot(\vc{A}\times\vc{B}), \label{eq:triple-escalar}\\
  \vc{A}\times(\vc{B}\times\vc{C})&=(\vc{A}\cdot\vc{C})\,\vc{B}-(\vc{A}\cdot\vc{B})\,\vc{C}. \label{eq:bac-cab}
\end{align}
\end{proposicion}
\begin{proof}
Con el tensor de Levi-Civita, $\vc{A}\cdot(\vc{B}\times\vc{C})=\varepsilon_{ijk}A_iB_jC_k$, y la
invariancia cíclica $\varepsilon_{ijk}=\varepsilon_{jki}=\varepsilon_{kij}$ da
\eqref{eq:triple-escalar} (el triple producto es el volumen orientado del
paralelepípedo, y no cambia al permutar cíclicamente). Para la segunda,
\[
  [\vc{A}\times(\vc{B}\times\vc{C})]_i=\varepsilon_{ijk}A_j\,\varepsilon_{klm}B_lC_m
  =(\delta_{il}\delta_{jm}-\delta_{im}\delta_{jl})A_jB_lC_m
  =B_i\,(\vc{A}\cdot\vc{C})-C_i\,(\vc{A}\cdot\vc{B}),
\]
donde se ha usado $\varepsilon_{ijk}=\varepsilon_{kij}$ y el \cref{lema:eps-delta}.
\end{proof}

% -------------------------------------------------------------------------
\subsection{Cuadro de derivadas vectoriales}\label{sec:cuadro-derivadas}
% -------------------------------------------------------------------------
Reunimos, para consulta, las identidades de las \cref{prop:productos,prop:fundamentales,prop:rotrot}:
\begin{center}
\renewcommand{\arraystretch}{1.55}
\small
\begin{tabular}{@{}ll@{}}
\toprule
\multicolumn{2}{c}{Derivadas de productos}\\
\midrule
$\nab(fg)=f\,\nab g+g\,\nab f$ & $\nab\cdot(f\vc{A})=f\,\nab\cdot\vc{A}+\vc{A}\cdot\nab f$\\
$\nab\times(f\vc{A})=\nab f\times\vc{A}+f\,\nab\times\vc{A}$ & $\nab\cdot(\vc{A}\times\vc{B})=\vc{B}\cdot(\nab\times\vc{A})-\vc{A}\cdot(\nab\times\vc{B})$\\
\multicolumn{2}{@{}l}{$\nab(\vc{A}\cdot\vc{B})=\vc{A}\times(\nab\times\vc{B})+\vc{B}\times(\nab\times\vc{A})+(\vc{A}\cdot\nab)\vc{B}+(\vc{B}\cdot\nab)\vc{A}$}\\
\multicolumn{2}{@{}l}{$\nab\times(\vc{A}\times\vc{B})=(\nab\cdot\vc{B})\vc{A}-(\nab\cdot\vc{A})\vc{B}+(\vc{B}\cdot\nab)\vc{A}-(\vc{A}\cdot\nab)\vc{B}$}\\
\midrule
\multicolumn{2}{c}{Derivadas segundas}\\
\midrule
$\nab\cdot(\nab f)=\Lap f=\pdd{f}{x}+\pdd{f}{y}+\pdd{f}{z}$ & $\nab\cdot(\nab\times\vc{A})=0$\\
$\nab\times(\nab f)=\vc{0}$ & $\nab\times(\nab\times\vc{A})=\nab(\nab\cdot\vc{A})-\Lap\vc{A}$\\
\bottomrule
\end{tabular}
\end{center}

% -------------------------------------------------------------------------
\subsection{Integrales que aparecen con frecuencia}\label{sec:integrales-frecuentes}
% -------------------------------------------------------------------------
En el cálculo de campos de distribuciones de carga aparecen una y otra vez las
siguientes primitivas, con $a>0$ constante:
\begin{align}
  \int\frac{\dif x}{(x^{2}+a^{2})^{1/2}}&=\ln\Bigl(x+\sqrt{x^{2}+a^{2}}\Bigr), &
  \int\frac{x\,\dif x}{(x^{2}+a^{2})^{1/2}}&=\sqrt{x^{2}+a^{2}},\label{eq:int-1}\\
  \int\frac{\dif x}{x^{2}+a^{2}}&=\frac{1}{a}\tan^{-1}\frac{x}{a}, &
  \int\frac{x\,\dif x}{x^{2}+a^{2}}&=\frac{1}{2}\ln\bigl(x^{2}+a^{2}\bigr),\label{eq:int-2}\\
  \int\frac{\dif x}{(x^{2}+a^{2})^{3/2}}&=\frac{x}{a^{2}\sqrt{x^{2}+a^{2}}}, &
  \int\frac{x\,\dif x}{(x^{2}+a^{2})^{3/2}}&=-\frac{1}{\sqrt{x^{2}+a^{2}}}.\label{eq:int-3}
\end{align}
Todas se comprueban derivando el segundo miembro. Por ejemplo, para la primera,
$\frac{\dif}{\dif x}\ln\bigl(x+\sqrt{x^{2}+a^{2}}\bigr)=\frac{1+x/\sqrt{x^{2}+a^{2}}}{x+\sqrt{x^{2}+a^{2}}}=\frac{1}{\sqrt{x^{2}+a^{2}}}$,
y para la quinta,
$\frac{\dif}{\dif x}\frac{x}{a^{2}\sqrt{x^{2}+a^{2}}}=\frac{(x^{2}+a^{2})-x^{2}}{a^{2}(x^{2}+a^{2})^{3/2}}=\frac{1}{(x^{2}+a^{2})^{3/2}}$.

% -------------------------------------------------------------------------
\subsection{Resumen de la delta de Dirac y de los polinomios de Legendre}\label{sec:resumen-delta}
% -------------------------------------------------------------------------
\begin{center}
\renewcommand{\arraystretch}{1.5}
\small
\begin{tabular}{@{}lll@{}}
\toprule
\multicolumn{3}{c}{Delta de Dirac}\\
\midrule
$\displaystyle\delta(x)=\begin{cases}0,&x\neq0\\ \infty,&x=0\end{cases}$ &
$\displaystyle\int_{-\infty}^{\infty}\delta(x)\,\dif x=1$ &
$\displaystyle\int_{-\infty}^{\infty}f(x)\,\delta(x-a)\,\dif x=f(a)$\\
$\displaystyle\delta^{3}(\vc{r})=\begin{cases}0,&\vc{r}\neq\vc{0}\\ \infty,&\vc{r}=\vc{0}\end{cases}$ &
$\displaystyle\int_{V_\infty}\delta^{3}(\vc{r})\,\dV=1$ &
$\displaystyle\int_{V_\infty}f(\vc{r})\,\delta^{3}(\vc{r}-\vc{a})\,\dV=f(\vc{a})$\\
\midrule
\multicolumn{3}{c}{Polinomios de Legendre y desarrollo de $1/\abs{\vc{r}-\vc{r}\,'}$}\\
\midrule
\multicolumn{3}{c}{$\displaystyle\frac{1}{\abs{\vc{r}-\vc{r}\,'}}=\frac{1}{r}\sum_{n=0}^{\infty}\Bigl(\frac{r'}{r}\Bigr)^{n}P_n(\cos\theta')\quad(r'<r)$}\\
$P_0(x)=1$, $\ P_1(x)=x$ & $P_2(x)=\tfrac12(3x^{2}-1)$ & $P_3(x)=\tfrac12(5x^{3}-3x)$\\
\bottomrule
\end{tabular}
\end{center}


% =========================================================================
%                              BIBLIOGRAFÍA
% =========================================================================
\printbibliography[title={Referencias}]

\end{document}
